\documentclass[11pt]{article}
\usepackage[margin=2.6cm]{geometry}
\usepackage[T1]{fontenc}
\usepackage[utf8]{inputenc}
\usepackage{lmodern}
\usepackage{amsmath,amssymb}
\usepackage{array,longtable}
\usepackage{listings}
\usepackage{xcolor}
\usepackage[hidelinks]{hyperref}
\usepackage{parskip}

\lstset{
  basicstyle=\ttfamily\small,
  breaklines=true,
  columns=fullflexible,
  keepspaces=true,
  frame=single,
  framerule=0.2pt,
  xleftmargin=0.6em,
  aboveskip=0.7em,
  belowskip=0.7em,
  literate={—}{{---}}1 {°}{{\ensuremath{^\circ}}}1 {≈}{{\ensuremath{\approx}}}1
}

\newcommand{\kw}[1]{\texttt{\uppercase{#1}}}
\newcommand{\code}[1]{\texttt{#1}}

\title{The \code{posim} Command Language\\[2pt]
       \large Complete Grammar \& Notebook Guide}
\author{physical\_object\_simulator}
\date{July 21, 2026}

\begin{document}
\maketitle
\tableofcontents

\section{What is this document about?}

\textbf{posim} is the front end of the \code{physical\_object\_simulator}.
Instead of writing Rust code to create objects and run physics, you type
short \emph{commands} --- one per line --- into a \emph{notebook}. Each
command line is:
\begin{enumerate}
\item chopped into \emph{tokens} by a \emph{lexer} (like the classic Unix
      tool \code{lex}/\code{flex});
\item checked against a \emph{grammar} and compiled into a small program
      by a \emph{parser} (like \code{yacc}/\code{bison});
\item executed by a \emph{stack machine} --- a tiny computer whose only
      memory is a stack of values --- which reads and writes the simulator
      exclusively through its public get/set functions.
\end{enumerate}
You never see steps 1--3; you type a line, press Enter (or shift-enter in
JupyterLab), and see the result. This document specifies exactly what you
may type.

There are three ways to talk to the same language:

\begin{center}
\begin{tabular}{lll}
\textbf{mode} & \textbf{start with} & \textbf{who it is for}\\\hline
notebook REPL & \code{cargo run} (or \code{posim}) & humans at a terminal\\
script        & \code{posim --script file}         & batch runs, reproducibility\\
dynamic notebook & \code{posim --notebook file}    & load a notebook, open its scene window, continue live\\
machine       & \code{posim --machine}             & programs, and the JupyterLab kernel\\
\end{tabular}
\end{center}

\section{Lexical structure (what the lexer sees)}

A command line is a sequence of \emph{tokens} separated by optional
whitespace. Everything from a \code{\#} to the end of the line is a
\emph{comment} and is ignored.

\subsection{Token kinds}

\begin{center}
\begin{tabular}{lp{4.2cm}p{6.6cm}}
\textbf{token} & \textbf{examples} & \textbf{rules}\\\hline
keyword & \code{NEW}, \code{set}, \code{Run} & case-insensitive; list in \S\ref{sec:keywords}\\
identifier & \code{obj0}, \code{position} & a letter or \code{\_}, then letters, digits, \code{\_}\\
string & \code{"ball"}, \code{"dumbell0"} & double-quoted, with the escapes \code{\textbackslash"} \code{\textbackslash\textbackslash} \code{\textbackslash n}; strings name objects (\kw{new} \ldots{} \kw{as}, \S\ref{sec:new}) and are passed to user functions (\S\ref{sec:userdefs})\\
number & \code{2}, \code{0.5}, \code{.5}, \code{1e-3} & 64-bit float; scientific notation allowed; a leading \code{-} is \emph{not} part of the number --- it is the negation operator\\
punctuation & \code{[ ] \{ \} ( ) , . =} & brackets build vectors, braces enclose initializers, \code{.} builds dotted paths\\
operators & \code{+ - * /} & ordinary arithmetic\\
\end{tabular}
\end{center}

If the lexer meets a character it does not know (say \code{@}), it stops
with an error naming the \emph{column} (character position, starting at~1):
\begin{lstlisting}
Err[1]: lexical error at column 10: unexpected character `@`
\end{lstlisting}

\subsection{Keywords}\label{sec:keywords}

\code{NEW SET GET DEL DELETE LIST STEP RUN STEPS METHOD ADAMS BDF SPRK
ENERGY COM MOMENTUM ANGMOM LAPLACE HELP POINT SPHERE CUBOID TORUS DISK
CYLINDER BOX RESET}

and, for the graphical scene window (documented in \code{scene\_info.pdf};
Examples 11--12 show live sessions):

\code{SCENE CREATE CLOSE DESTROY TRANSLATE ROTATE ZOOM IN OUT HIDE SHOW
REFRESH REDRAW START STOP PAUSE REVERSE SET\_TIME\_STEP SETTIMESTEP
STATUS EVENTS ALL}

and, for rigid-body collisions (\S\ref{sec:collide}):

\code{COLLIDE CONTACTS ON OFF}

and, for the quantum families (\S\ref{sec:qm}, \S\ref{sec:qm2},
\S\ref{sec:qm3}):

\code{QM QM2 QM3}

and, for named objects, session variables and user functions
(\S\ref{sec:new}, \S\ref{sec:userdefs}):

\code{AS LET FUNCS DUMBBELL}

(\kw{destroy} is an alias for \kw{close}; \kw{settimestep} for
\kw{set\_time\_step}; the spellings \kw{cube} and \kw{disc} are lexer
aliases for \kw{cuboid} and \kw{disk}; \kw{functions} for \kw{funcs};
and the single-b spelling \kw{dumbell} for \kw{dumbbell}.
\code{collisions} is deliberately \textbf{not} a keyword so the field
path \code{system.collisions} keeps its own spelling. \kw{def} is not
a keyword either: a line starting \code{DEF } is a \textbf{line form}
recognized before the ordinary grammar --- see \S3 and
\S\ref{sec:userdefs}.)

Keywords are reserved \emph{at the start of paths}, but \textbf{field
names may reuse keyword spellings} --- \code{obj0.momentum} and
\code{system.method} work even though \kw{momentum} and \kw{method} are
keywords, because after a dot (or inside \code{NEW \{ ... \}}) the parser
accepts either an identifier or a keyword as a field name. The same
holds for the scene keywords: making \kw{show} or \kw{in} reserved does
not stop you from ever using those spellings as field names after a
\code{.}.

\subsection{Magics}

A line whose first character is \code{\%} is a \emph{notebook magic}
(\S\ref{sec:magics}) and is handled by the notebook itself; it never
reaches the lexer.

\section{The grammar (what the parser accepts)}

The grammar in EBNF (Extended Backus--Naur Form: \code{[x]} means
optional, \code{\{x\}} means zero-or-more repetitions, \code{|} separates
alternatives):

\begin{lstlisting}
line     := command | magic ;

command  := "NEW" shape [ "AS" (IDENT | STRING) ]
                         [ "{" init { "," init } "}" ]
                                              (* AS registers a user
                                                 name for the object *)
          | "SET" path "=" expr
          | "GET" path
          | "DEL" NUMBER
          | "LIST"
          | "STEP" expr                       (* advance by dt      *)
          | "RUN" expr [ "STEPS" NUMBER ]     (* advance by t, n outs *)
          | "METHOD" ( "ADAMS" | "BDF" | "SPRK" IDENT [ NUMBER ] )
          | "ENERGY" | "COM" | "MOMENTUM" | "ANGMOM"
          | "LAPLACE" NUMBER
          | "RESET" | "HELP"
          | "SCENE" scenecmd                  (* graphical scene     *)
          | "COLLIDE" [ "ON" | "OFF" ]        (* bare: report status *)
          | "CONTACTS"                        (* list last contacts  *)
          | "LET" IDENT "=" expr              (* session variable    *)
          | "FUNCS"                           (* list user functions *)
          | "SHOW" IDENT                      (* print a function    *)
          | "BOX" [ "OFF" | expr ]            (* rigid bounding box:
                                                 expr = inner side
                                                 length; bare = status;
                                                 OFF removes it      *)
          | expr ;                            (* bare expression     *)
scenecmd := "CREATE" [ NUMBER ]               (* open window [port]  *)
          | "CLOSE"                           (* aka DESTROY         *)
          | "TRANSLATE" term term [ term ]    (* camera dx dy [dz]   *)
          | "ROTATE" term term                (* camera dyaw dpitch  *)
          | "ZOOM" ( "IN" | "OUT" | term )    (* factor > 1 zooms in *)
          | "HIDE" [ NUMBER | "ALL" ]         (* default: ALL        *)
          | "SHOW" [ NUMBER | "ALL" ]
          | "REFRESH"                         (* re-sync from state  *)
          | "REDRAW"                          (* re-send full scene  *)
          | "START" | "STOP" | "PAUSE" | "REVERSE"
          | "RESET"                           (* re-initialize: all
                                                 values and the time
                                                 return to initial;
                                                 START re-starts    *)
          | "SET_TIME_STEP" term              (* args are term-level:
                                                 -5 is negative five;
                                                 parenthesize sums  *)
          | "STATUS" | "EVENTS" ;
shape    := "POINT" | "SPHERE" | "CUBOID" | "TORUS" | "DISK" | "CYLINDER"
          | "DUMBBELL" ;                      (* two solid spheres +
                                                 a rigid rod, one
                                                 rigid body          *)

(* User-defined functions are a LINE FORM handled before this
   grammar: DEF name(param [= default], ...) { body } — the body is
   newline/;-separated commands using the parameters as variables;
   each body line must itself satisfy this grammar. Invocation uses
   the ordinary call syntax name(arg, ...); trailing parameters
   take their defaults. *)
init     := IDENT "=" expr ;
path     := IDENT { "." IDENT } ;             (* objN.field[.x|y|z|w],
                                                 system.field,
                                                 contactK.field,
                                                 name.field for
                                                 AS-registered names *)
expr     := sum { ("<" | "<=" | ">" | ">=" | "==" | "!=") sum } ;
                                    (* comparisons yield 1 or 0, LOWEST
                                       precedence; 1/0 makes
                                       `(x > a) * (x < b)` an indicator
                                       function -- how a piecewise
                                       potential is written *)
sum      := term { ("+" | "-") term } ;
term     := unary { ("*" | "/") unary } ;
unary    := "-" unary | atom ;
atom     := NUMBER | STRING | "[" expr { "," expr } "]" | "(" expr ")"
          | IDENT "(" [ expr { "," expr } ] ")"   (* builtin or user
                                                     function call   *)
          | path
          | IDENT ;                           (* parameter / LET var *)
\end{lstlisting}

Operator precedence is the usual one: \code{*} and \code{/} bind tighter
than \code{+} and \code{-}; unary minus binds tightest; parentheses
override everything.

\textbf{One subtlety worth learning early:} \kw{get} takes \emph{only a
path}. To do arithmetic with a field, use a \emph{bare expression}:
\begin{lstlisting}
get obj0.position.x - 5        # ERROR -- GET wants just a path
obj0.position.x - 5            # OK -- bare expression
\end{lstlisting}

\textbf{A second subtlety --- \kw{scene} arguments are \emph{terms},
not full expressions.} The scene commands take several numbers in a row
with no commas between them. If those numbers were parsed as full
expressions, \code{scene rotate 15 -5} would be read as \emph{one}
argument $15-5=10$ (whitespace never matters to the lexer) and the
command would then be missing its second argument. So scene arguments
stop one level lower in the grammar, at \code{term}: \code{-5} is
negative five, \code{2*2} and \code{1/3} still work, but a sum or
difference must be parenthesized:

\begin{lstlisting}
scene rotate 15 -5             # yaw +15 deg, pitch -5 deg (two arguments)
scene zoom (1 + 0.5)           # a sum needs parentheses
scene translate 2*2 0 1/2      # * and / are fine unparenthesized
\end{lstlisting}

\textbf{A third subtlety --- bare identifiers compile and resolve at
execution.} An \code{IDENT} alone is an ordinary atom: the constants
\code{pi} and \code{tau}, a function parameter, or a \kw{let}
variable (\S\ref{sec:userdefs}). The name is looked up when the line
\emph{runs}, not when it parses, so a function body may use
parameters freely; a name that is bound to nothing fails at execution
with an error that lists the three ways to bind one (\kw{let}, a
parameter, or a registered object's \code{name.field} path). A call
\code{IDENT(...)} is a builtin (\code{dot}, \code{cross},
\code{norm}, \code{normalize}, \code{sqrt}, \code{abs}, \code{sin},
\code{cos}, \code{exp}, \code{log}) when the name matches one, and a
user-defined function otherwise.

\section{The type system (what values exist)}

Every expression evaluates to one of these \emph{values}:

\begin{center}
\begin{tabular}{llp{7.2cm}}
\textbf{type} & \textbf{written as} & \textbf{notes}\\\hline
number & \code{2}, \code{-0.5}, \code{1e-3} & 64-bit float\\
vec3 & \code{[1, 2, 3]} & exactly 3 numeric entries\\
quaternion & \code{[w, x, y, z]} & exactly 4 numeric entries, \textbf{w first}; assigned to \code{orientation} it is renormalized to unit length\\
mat3 & \code{[[a,b,c],[d,e,f],[g,h,i]]} & 3 rows of 3 numbers\\
string & \code{"ball"} & a double-quoted literal (\S2.1) --- names objects (\kw{new} \ldots{} \kw{as}) and feeds user functions; also produced as results like \code{obj0} or run summaries\\
\end{tabular}
\end{center}

Bracket literals are \emph{shape-directed}: 3 numbers make a vec3,
4 numbers make a quaternion, 3 vec3s make a mat3.

Arithmetic is \emph{type-checked}:
\begin{center}
\begin{tabular}{lp{9.5cm}}
\textbf{operation} & \textbf{allowed}\\\hline
\code{+}, \code{-} & number$\pm$number, vec3$\pm$vec3, quat$\pm$quat\\
\code{*} & number$\cdot$number, number$\cdot$vec3, number$\cdot$quat, number$\cdot$mat3, mat3$\cdot$vec3, mat3$\cdot$mat3, quat$\cdot$quat (Hamilton product)\\
\code{/} & number/number, vec3/number\\
vec3\,\code{*}\,vec3 & \textbf{forbidden} --- the error tells you to use \code{dot()} or \code{cross()}\\
\end{tabular}
\end{center}

Builtins: \code{dot(a,b)}, \code{cross(a,b)}, \code{norm(v)} (length;
also quaternions and numbers), \code{normalize(v)}, and scalar
\code{sqrt abs sin cos exp log}. Constants: \code{pi}, \code{tau}
($=2\pi$).

\subsection{Special functions}\label{sec:specialfns}

The same call syntax reaches the \code{special\_functions} library.
Nothing about the grammar changes to accommodate them --- the production
\code{IDENT "(" [expr \{"," expr\}] ")"} already admits any builtin ---
so adding a function is a \emph{registration}, not a grammar change.

\textbf{Orders must be whole numbers.} Where an argument is an integer
order (the $n$, $\ell$, $m$ below), a fractional value is an error
rather than being quietly truncated:

\begin{lstlisting}
In[1]:= hermite_h(2.5, 1)
Err[1]: hermite_h(): argument 1 must be a whole number (an integer order), got 2.5
\end{lstlisting}

That refusal is deliberate. Truncating to \code{hermite\_h(2, 1)} would
return a confident, wrong number, and you would have no way to notice.

\begin{center}
\begin{tabular}{lp{9.4cm}}
\textbf{family} & \textbf{functions}\\
\hline
spherical Bessel & \code{sph\_j(n,x)}, \code{sph\_y(n,x)},
  \code{sph\_j\_prime(n,x)}, \code{sph\_y\_prime(n,x)}\\
Legendre & \code{legendre\_p(n,x)}, \code{legendre\_p\_prime(n,x)},
  \code{assoc\_legendre\_p(l,m,x)}, \code{norm\_assoc\_legendre\_p(l,m,x)}\\
spherical harmonics & \code{sph\_harm(l,m,theta,phi)} $\to$
  \code{[re, im]}, \code{sph\_harm\_real(l,m,theta,phi)}\\
orthogonal polynomials & \code{hermite\_h(n,x)}, \code{hermite\_he(n,x)},
  \code{laguerre\_l(n,x)}, \code{laguerre\_l\_assoc(n,alpha,x)},
  \code{chebyshev\_t(n,x)}, \code{chebyshev\_u(n,x)},
  \code{gegenbauer\_c(n,alpha,x)}, \code{jacobi\_p(n,alpha,beta,x)}\\
cylindrical Bessel & \code{bessel\_j(n,x)},
  \code{bessel\_j\_array(n\_max,x)} $\to$ list\\
cylindrical Bessel, \textbf{complex argument} &
  \code{bessel\_j\_z(n,z)}, \code{bessel\_i\_z(n,z)},
  \code{bessel\_y\_z(n,z)}, \code{bessel\_k\_z(n,z)}\\
cylindrical Bessel, \textbf{any real order} &
  \code{bessel\_j\_nu(nu,z)}, \code{bessel\_i\_nu(nu,z)},
  \code{bessel\_y\_nu(nu,z)}, \code{bessel\_k\_nu(nu,z)}\\
Hankel (travelling waves) &
  \code{hankel\_h1\_z(n,z)}, \code{hankel\_h2\_z(n,z)},
  \code{hankel\_h1\_nu(nu,z)}, \code{hankel\_h2\_nu(nu,z)}\\
Hankel derivatives &
  \code{hankel\_h1\_prime\_z(n,z)}, \code{hankel\_h2\_prime\_z(n,z)},
  \code{hankel\_h1\_prime\_nu(nu,z)},
  \code{hankel\_h2\_prime\_nu(nu,z)}\\
spherical Hankel &
  \code{sph\_hankel\_h1(n,x)}, \code{sph\_hankel\_h2(n,x)},
  \code{sph\_hankel\_h1\_prime(n,x)},
  \code{sph\_hankel\_h2\_prime(n,x)}\\
scaled forms &
  \code{bessel\_j\_scaled(nu,z)}, \code{bessel\_y\_scaled(nu,z)},
  \code{bessel\_i\_scaled(nu,z)}, \code{bessel\_k\_scaled(nu,z)},
  \code{hankel\_h1\_scaled(nu,z)}, \code{hankel\_h2\_scaled(nu,z)}\\
gamma, complex argument &
  \code{gamma\_z(z)}, \code{ln\_gamma\_z(z)}, \code{rgamma\_z(z)}\\
Airy, complex argument &
  \code{airy\_z(z)} $\to$ \code{[Ai, Ai', Bi, Bi']}\\
quadrature & \code{gauss\_legendre(n)} $\to$ \code{[nodes, weights]}\\
eigenproblems & \code{eigenvalues(matrix)} $\to$ list,
  \code{jacobi\_eigen(matrix)} $\to$ \code{[values, vectors]}\\
angular momentum & \code{wigner\_3j(j1,j2,j3,m1,m2,m3)},
  \code{wigner\_6j(j1,j2,j3,j4,j5,j6)},
  \code{wigner\_9j(a,b,c,d,e,f,g,h,i)},
  \code{clebsch\_gordan(j1,m1,j2,m2,j3,m3)}\\
linear algebra & \code{solve\_tridiag(sub, diag, sup, rhs)} $\to$ list,
  \code{solve\_tridiag\_c(\ldots)} $\to$ list,
  \code{solve\_cyclic\_tridiag\_c(sub, diag, sup, rhs, bl, tr)} $\to$ list\\
utility & \code{rel\_err(a, b)}\\
\end{tabular}
\end{center}

\code{wigner\_9j} takes its nine arguments \textbf{row by row} --- it
recouples four angular momenta, and is the overlap between coupling
(1,2) and (3,4) first versus (1,3) and (2,4) first. It is evaluated as a
single sum over 6-j symbols and vanishes when any of the six triads
fails to close.

\textbf{Angular momenta may be half-integers}, so \code{wigner\_3j},
\code{wigner\_6j} and \code{clebsch\_gordan} take plain numbers rather
than demanding whole ones --- \code{clebsch\_gordan(0.5, 0.5, 0.5, -0.5,
1, 0)} is exactly what you want to write for two spin-$\tfrac12$
particles. A coupling that violates a selection rule returns
\textbf{0}, which is the mathematically correct answer, not an error; a
value that is not an angular momentum at all ($j=0.3$, or a negative
$j$) \emph{is} an error.

\subsubsection*{Complex-argument Bessel}

\code{bessel\_j\_z(n, z)} and \code{bessel\_i\_z(n, z)} take and return
\textbf{complex} values --- they exist because the language has
\code{Value::Complex} (\S\ref{sec:complex}). Real arguments promote.

The algorithm is the same Miller downward recurrence as the real
routine, deliberately: the three-term recurrence and the normalisation
$J_0 + 2(J_2+J_4+\cdots) = 1$ are both identities in $z$, real or not,
the second following from the generating function at $t=1$.

\textbf{Accuracy of $J$ falls with the imaginary part}, by cancellation
rather than anything about the recurrence: the terms grow like
$e^{|\mathrm{Im}\,z|}$ while their sum is exactly 1, so about
$|\mathrm{Im}\,z|/\ln 10$ decimal digits are lost. Measured against the
generating-function identity:

\begin{center}
\begin{tabular}{lcccccc}
$|\mathrm{Im}\,z|$ & 0 & 5 & 8 & 12 & 18 & 25\\
\hline
relative error of $J$ & $10^{-16}$ & $10^{-14}$ & $10^{-13}$ &
$10^{-11}$ & $10^{-9}$ & $10^{-6}$\\
\end{tabular}
\end{center}

\textbf{That law governs $J$ alone, and an earlier version of this
manual let it stand for the whole family. It does not.} $I$ is $J$ at
right angles and obeys the mirror law, but $Y$ comes from an ascending
series and $K$ is assembled from $J$ and $Y$ at imaginary argument. A
series fails where a recurrence does not --- on the \textbf{real} axis,
which is precisely where $J$ is at its best. There are four laws:
\[
  \text{relative error} \;\sim\; 10^{-16}e^{L},\qquad
  L = \begin{cases}
    |\mathrm{Im}\,z| & \code{bessel\_j\_z}\\
    |\mathrm{Re}\,z| & \code{bessel\_i\_z}\\
    |z| - |\mathrm{Im}\,z| & \code{bessel\_y\_z}\\
    \max(2|\mathrm{Re}\,z|,\,|z|) + \mathrm{Re}\,z & \code{bessel\_k\_z}
  \end{cases}
\]
Measured against Cephes on the axis where each is at its worst:

\begin{center}
\begin{tabular}{lccccc}
$x$ (real) & 1 & 10 & 20 & 30 & 35\\
\hline
\code{bessel\_j\_z(0,x)}  & 1e-16 & 3e-16 & 5e-16 & 2e-15 & 7e-16\\
\code{bessel\_j\_z(0,ix)} & 2e-16 & 1e-13 & 5e-9  & 5e-5  & 3e-3\\
\code{bessel\_i\_z(0,x)}  & 2e-16 & 1e-13 & 5e-9  & 5e-5  & 3e-3\\
\code{bessel\_y\_z(0,x)}  & 1e-15 & 3e-12 & 3e-8  & 2e-4  & 6e-2\\
\code{bessel\_k\_z(0,x)}  & 7e-16 & 3e-5  & 8e8   & 5e21  & 7e27\\
\end{tabular}
\end{center}

\textbf{That table is now history rather than guidance.} It describes
what each \emph{route} costs, and every failure in it has since been
fixed by giving the function a second route:
\begin{itemize}
\item \code{bessel\_j\_z} near the imaginary axis uses
  $J_\nu(z) = i^\nu I_\nu(-iz)$, putting the argument back near the real
  axis where $I$ has nothing to cancel;
\item \code{bessel\_y\_z} near the imaginary axis uses
  $Y_n(z) = i^{n+1}I_n(w) - (2/\pi)i^{-n}K_n(w)$ with $w = -iz$, instead
  of the upward recurrence in $n$ --- the wrong direction there, because
  $Y_n(iy)$ is mostly $I_n(y)$ and $I$ is the recessive solution of that
  recurrence;
\item \code{bessel\_y\_z} along the real axis and \code{bessel\_k\_z}
  everywhere use their own $1/z$ expansions, single series with nothing
  to cancel;
\item \code{bessel\_j\_z} and \code{bessel\_y\_z} in the wedge either
  side of the \textbf{negative real axis} continue the Hankel expansions
  from the positive one (DLMF 10.11.3, 10.11.4). Those expansions have
  sectors ending at $\arg z = \pi$, so the cut is the one direction
  neither reaches directly, and $w = -z$ puts it back where both are at
  their best. \code{bessel\_k\_z} needed only its sector widening: DLMF
  10.40.2 is valid to $|\arg z| < 3\pi/2$, so the whole principal sheet
  was always interior to it.
\end{itemize}
The routes are chosen by comparing error estimates, so nothing changes
in what you write. What changed is the answers:

\begin{center}
\begin{tabular}{lcc}
 & \textbf{before} & \textbf{after}\\
\hline
\code{bessel\_y\_z(0, 40)} & wrong in the first digit & 1e-15\\
\code{bessel\_k\_z(0, 20)} & out by a factor of 8e8 & 2e-16\\
\code{bessel\_y\_z(2, 29.4e$^{1.6i}$)} & Wronskian 4.5e-6 & below 1e-13\\
\end{tabular}
\end{center}

Measured across $|z|$ from 5 to 40 and $\arg z$ over the upper half
plane, the $J$--$Y$ Wronskian residual is now \textbf{$10^{-10}$ or
better and mostly below $10^{-15}$}.

\textbf{On the cut itself that Wronskian is the wrong instrument.} It is
dominated there by the exponentially \emph{recessive} Hankel member, so
it measures the Stokes phenomenon rather than the answer. The right test
is the pair of exact continuation identities
\[
  J_n(xe^{i\pi}) = (-1)^n J_n(x), \qquad
  Y_n(xe^{i\pi}) = (-1)^n\big[Y_n(x) + 2iJ_n(x)\big],
\]
which relate a point on the cut to one on the positive real axis, where
everything is at its best. Against those, \code{bessel\_j\_z} is
\textbf{exact} and \code{bessel\_y\_z} is $10^{-14}$ out to $|z| = 300$;
the wedge either side is $10^{-12}$ or better, where before it reached
$1.0$ at $|z| = 60$. The branch jump is still exactly $4i(-1)^nJ_n$, and
a test says so --- widening coverage across a cut is only correct if the
cut stays where it was. The $J$ and $I$ rows of the table
agree digit for digit because they are the same measurement ---
$I_n(z)$ \emph{is} $J_n(iz)$.

Two of those defects were found by accident and one by a test written
for something else, and in each case the elementary Wronskian settled
which side was wrong. The $K$ one is worth stating plainly: its identity
$K_n(z) = (\pi/2)i^{n+1}[J_n(iz) + iY_n(iz)]$ \textbf{cancels by
construction} on the real axis --- both terms carry $I_n(x)$, of size
$e^{x}$, and what survives is $e^{-x}$. No ingredient was inaccurate.
The identity was the wrong way to compute a recessive function.

\subsubsection*{$Y_n$ and $K_n$: a different method, because they need one}

$Y_n$ has a \textbf{logarithmic branch point} at the origin, so no
recurrence produces it from $J_n$ alone. It comes from the ascending
series (DLMF 10.8.1) for $Y_0$ and $Y_1$ --- the one carrying the
$(2/\pi)\ln(z/2)J_n(z)$ term and digamma coefficients --- followed by
\textbf{upward} recurrence in $n$.

That direction is the opposite of $J$'s, deliberately: $Y_n$ \emph{grows}
with order while $J_n$ decays, so the stable sweep for one is the
unstable sweep for the other.

$K_n$ follows by identity, $K_n(z) = (\pi/2) i^{n+1}[J_n(iz)+iY_n(iz)]$
(DLMF 10.27.8). Both are verified by Wronskians, whose right-hand sides
are elementary:
\[
  J_{n+1}Y_n - J_nY_{n+1} = \frac{2}{\pi z},\qquad
  I_nK_{n+1} + I_{n+1}K_n = \frac{1}{z}.
\]

Both inherit the branch cut of $\ln$ along the negative real axis, where
they jump by $4iJ_n$ --- crossing changes $\arg$ by $2\pi$, and the
$(2/\pi)\ln(z/2)J$ term turns that into $(2/\pi)(2\pi i)J$. Both are
singular at $z=0$ and report an error there.

\subsubsection*{Non-integer order}

The \code{\_z} family insists on a whole order. The \code{\_nu} family
does not: \code{bessel\_j\_nu(1.3, 2 + 1i)} is $J_{1.3}(2+i)$, and
half-integer orders --- the ones spherical Bessel functions are made of
--- are the common case. Order stays \textbf{real}; complex \emph{order}
is not implemented.

Non-integer order needs no recurrence at all, because the two formulas
that are hardest at integer order become the easy ones here:
\[
  Y_\nu(z) = \frac{J_\nu(z)\cos\nu\pi - J_{-\nu}(z)}{\sin\nu\pi},
  \qquad
  K_\nu(z) = \frac{\pi}{2}\,
    \frac{I_{-\nu}(z) - I_\nu(z)}{\sin\nu\pi}.
\]
Both have $\sin\nu\pi$ downstairs, which is exactly why they are useless
at whole $\nu$ and exactly why $Y_n$ needed its own logarithmic series.
So the whole family reduces to \textbf{one ascending series} (DLMF
10.2.2 for $J$, 10.25.2 for $I$) evaluated at $\pm\nu$, and the
reflections do the rest. Orders within $10^{-9}$ of a whole number are
handed to the \code{\_z} routines --- not as a convenience but because
by then the reflection has already lost nine digits to cancellation.
\code{bessel\_y\_nu(2, z)} gives the right answer; it is
\code{bessel\_y\_z(2, z)} underneath.

The series divides by $\Gamma(\nu+k+1)$ using the \textbf{reciprocal}
gamma, which is \emph{zero} at the poles. That is what makes
$J_{-n}(z) = (-1)^nJ_n(z)$ come out for free at whole $n$, and it keeps
very large orders in range where $\Gamma$ itself would overflow past
about $171$.

\textbf{Accuracy is governed by a different law from the \code{\_z}
family}, and it is important not to carry the \code{\_z} advice over.
The largest term of the series is of size $e^{|z|}$, so what is lost is
the ratio of that to the answer:
\[
  \text{relative error} \;\sim\; 10^{-16}e^{L},\qquad
  L = \begin{cases}
    |z| - |\operatorname{Im}z| & \text{for } J_\nu,\; Y_\nu\\
    |z| + \operatorname{Re}z   & \text{for } I_\nu,\; K_\nu
  \end{cases}
\]

\begin{center}
\begin{tabular}{l|ccccc}
$L$ & 0 & 10 & 20 & 30 & $>35$\\
\hline
relative error & 1e-16 & 1e-12 & 1e-8 & 1e-4 & nothing left\\
\end{tabular}
\end{center}

$J$ and $Y$ are therefore at their \textbf{worst on the real axis} and
at machine precision straight up the imaginary one --- good to $|z|=70$
there. $I$ and $K$ are the mirror image: worst on the \textbf{positive}
real axis, exact along the negative. On the real axis that means $|z|$
up to about $30$ for $J$ and $Y$ but only $15$ for $I$ and $K$, since
$K$ is a \emph{difference} of two $I$s whose leading parts cancel ---
that cancellation is what $K$ \emph{is}, not a defect of the method.

Large \textbf{order} costs nothing: at $\nu = 150$ the order recurrence
still closes to $10^{-13}$, because nothing cancels once $\nu$ exceeds
$|z|$.

Measured, not asserted:
\begin{lstlisting}
cargo run -p special_functions --release --example bessel_nu_accuracy
\end{lstlisting}
prints the full error surface against the closed forms
$J_{1/2}(z)=\sqrt{2/\pi z}\,\sin z$ and
$K_{1/2}(z)=\sqrt{\pi/2z}\,e^{-z}$, which hold for complex $z$ and
share no code with the series. The bound $10^{-14}e^{L}$ is pinned by
the test suite at every point of that surface. Beyond it, use the
integer-order routines where the order allows --- they are Miller
recurrence and reach much further along the real axis. Uniform
asymptotics for large $|z|$ at non-integer order are \textbf{not}
implemented.

\begin{lstlisting}
In[1]:= bessel_j_nu(0.5, 2)
Out[1]= 0.5130161365618277 + 0i
In[2]:= sqrt(2/(pi*2))*sin(2)
Out[2]= 0.5130161365618278
In[3]:= bessel_j_nu(-0.5, 2)
Out[3]= -0.2347857104062484 + 0i
In[4]:= sqrt(2/(pi*2))*cos(2)
Out[4]= -0.2347857104062485
In[5]:= bessel_j_nu(1.3, 2 + 1i)
Out[5]= 0.7246607223671551 + 0.08999135538552433i
In[6]:= bessel_y_nu(2, 1.6 + 0.9i)
Out[6]= -0.6563546517741834 + 0.4094584881356723i
In[7]:= bessel_y_z(2, 1.6 + 0.9i)
Out[7]= -0.6563546517741834 + 0.4094584881356723i
In[8]:= bessel_j_z(1.5, 2)
Err[8]: bessel_j_z(): argument 1 must be a whole number (an integer order), got 1.5
\end{lstlisting}

\code{Out[1]}/\code{Out[2]} and \code{Out[3]}/\code{Out[4]} are the two
half-integer closed forms agreeing to the last digit, computed two
entirely different ways. \code{Out[6]} and \code{Out[7]} are
byte-identical, which is the whole-order handover doing what it claims.
\code{In[8]} is the contrast: the \code{\_z} form still refuses a
fractional order rather than truncating it.

\subsubsection*{Hankel functions --- the travelling-wave pair}

$J$ and $Y$ are the \emph{standing}-wave basis of the Bessel equation.
$H^{(1)} = J + iY$ and $H^{(2)} = J - iY$ are the \emph{travelling}-wave
basis of the same equation, and for wave problems they are the ones you
want:
\[
  H^{(1)}_\nu(x) \sim \sqrt{\frac{2}{\pi x}}\,
    e^{i(x - \nu\pi/2 - \pi/4)},
\]
a pure $e^{+ikr}/\sqrt{r}$ outgoing cylindrical wave under the
$e^{-i\omega t}$ convention. A scattering boundary condition is stated
in terms of $H^{(1)}$, not in terms of $J$ and $Y$. The spherical pair
\code{sph\_hankel\_h1(n,x)} $= j_n(x) + iy_n(x)$ does the same job in
three dimensions; it takes a \textbf{real} argument and returns a
complex value.

The \code{\_z} forms take a whole order, the \code{\_nu} forms any real
order, and \code{\_prime} gives the derivative, from
$C'_\nu(z) = C_{\nu-1}(z) - (\nu/z)C_\nu(z)$ (DLMF 10.6.2), which holds
for every cylinder function and so needs no separate algorithm.

\textbf{Why these are entry points at all.} Each is two calls to
functions the language already had, and an earlier version of this
manual said so and left it there. That was wrong, for the reason the
accuracy note below makes concrete: the assembly $J + iY$ silently
destroys most of its digits in half the plane, and a user writing it by
hand has no way to know. A named function is where that knowledge lives.

\begin{lstlisting}
In[1]:= hankel_h1_z(0, 3)
Out[1]= -0.26005195490193356 + 0.37685001001279045i
In[2]:= hankel_h1_prime_z(0, 3)
Out[2]= -0.3390589585259365 - 0.3246744247917999i
In[3]:= hankel_h1_z(1, 3)
Out[3]= 0.3390589585259365 + 0.3246744247917999i
In[4]:= abs(sph_hankel_h1(3, 800) * 800)
Out[4]= 1.0000046875439457
In[5]:= hankel_h1_z(0, 0 + 10i)
Out[5]= 0 - 0.00001131945100496523i
In[6]:= hankel_h2_z(0, 0 + 10i)
Out[6]= 5631.433256931902 + 0.00001131945100496523i
In[7]:= hankel_h1_nu(0.5, 2 - 0.5i)
Out[7]= 0.7802707009570485 + 0.4802074888838277i
In[8]:= sph_hankel_h1(0, 2.3)
Out[8]= 0.32421965746813924 + 0.2896852266434018i
In[9]:= sph_hankel_h2(0, 2.3)
Out[9]= 0.32421965746813924 - 0.2896852266434018i
\end{lstlisting}

\code{Out[2]} and \code{Out[3]} are exact negatives, which is
$H^{(1)\prime}_0 = -H^{(1)}_1$. \code{Out[4]} is the outgoing-wave
property $|x\,h^{(1)}_n(x)| \to 1$, and the gap from 1 is
$4.6875\times10^{-6}$, which is $n(n+1)/(4x^2)$ for $n=3$, $x=800$ to
every printed digit. \code{Out[8]} and \code{Out[9]} are conjugates, as
they must be for real $x$.

\textbf{\code{Out[5]} and \code{Out[6]} are the accuracy warning in one
pair of lines.} Both come from the same $J_0(10i)$ and $Y_0(10i)$, each
about $2800$ in magnitude. $H^{(2)}$ comes out as $5631$; $H^{(1)}$
comes out as $1.13\times10^{-5}$, which is what is left after those two
large numbers cancel. Above the real axis $H^{(1)}$ is the small
difference of large quantities and is only as accurate as that
subtraction allows:
\[
  \text{rel.\ error of } H^{(1)} \sim 10^{-16}e^{3\,\mathrm{Im}\,z}
  \quad(\mathrm{Im}\,z > 0),\qquad
  H^{(2)} \text{ likewise for } \mathrm{Im}\,z < 0 .
\]
So $H^{(1)}$ is good to $\mathrm{Im}\,z \approx 8$, has three digits
left at $10$, and is gone by $12$; below the axis it is exact, and
$H^{(2)}$ mirrors all of it. \textbf{Use whichever of the two is on its
good side.} Switching between the \code{\_z} and \code{\_nu} forms does
not help --- measured, they agree to within a factor of $1.5$, because
at whole order \code{\_nu} hands $Y$ to \code{\_z} and the two share the
dominant ingredient. Only a scaled formulation (returning
$e^{-iz}H^{(1)}$ and letting the caller supply the exponential, as AMOS
does) removes the problem, and that is \textbf{not} implemented.

That $e^{3\,\mathrm{Im}\,z}$ is the same law \code{bessel\_k\_z} obeys
on the real axis, which is no coincidence:
$K_\nu(y) = (\pi/2)i^{\nu+1}H^{(1)}_\nu(iy)$ (DLMF 10.27.8), so they are
literally the same computation.

The spherical pair has none of this trouble --- $j_n$ and $y_n$ are
recurrences on the real line, and \code{sph\_hankel\_h1} is accurate
wherever they are.

Measured, not asserted:
\begin{lstlisting}
cargo run -p special_functions --release --example hankel_accuracy
\end{lstlisting}

\subsubsection*{Scaled forms --- the exponential factored out}

The three sections above each end in the same place: a small quantity
computed from large ones, losing digits it cannot get back. $Y_0(40)$ is
wrong in its first digit, $K_0(15)$ is worthless, $H^{(1)}$ above the
real axis is gone by $\mathrm{Im}\,z = 12$. Each said the remedy was ``a
scaled formulation, which is not implemented''. It is now.

\begin{center}
\begin{tabular}{lp{6.2cm}}
\textbf{function} & \textbf{returns}\\
\hline
\code{bessel\_j\_scaled(nu,z)} & $e^{-|\mathrm{Im}\,z|}J_\nu(z)$\\
\code{bessel\_y\_scaled(nu,z)} & $e^{-|\mathrm{Im}\,z|}Y_\nu(z)$\\
\code{bessel\_i\_scaled(nu,z)} & $e^{-|\mathrm{Re}\,z|}I_\nu(z)$\\
\code{bessel\_k\_scaled(nu,z)} & $e^{z}K_\nu(z)$\\
\code{hankel\_h1\_scaled(nu,z)} & $e^{-iz}H^{(1)}_\nu(z)$\\
\code{hankel\_h2\_scaled(nu,z)} & $e^{iz}H^{(2)}_\nu(z)$\\
\end{tabular}
\end{center}

\textbf{Multiplying by the exponential afterwards would fix nothing} ---
the digits are gone before the multiplication. What makes these work is
a different algorithm: the asymptotic expansions of DLMF 10.17 and
10.40, plain series in $1/z$ with leading term 1 and \textbf{no
cancellation anywhere in them}:
\[
  e^{-iz}H^{(1)}_\nu(z) \sim \sqrt{\tfrac{2}{\pi z}}\,
    e^{-i(\nu\pi/2+\pi/4)}\,S(i),
  \qquad
  e^{z}K_\nu(z) \sim \sqrt{\tfrac{\pi}{2z}}\,S(1),
\]
\[
  S(c) = \sum_k \frac{c^k a_k(\nu)}{z^k},\qquad
  a_0 = 1,\quad a_k = a_{k-1}\frac{4\nu^2 - (2k-1)^2}{8k}.
\]
$J$ and $Y$ then come from the Hankel pair without the growing
exponential ever being formed, because $e^{\pm iz - |\mathrm{Im}\,z|}$
has modulus at most 1 by construction.

\begin{lstlisting}
In[1]:= bessel_y_scaled(0, 40)
Out[1]= 0.12593641705826092 + 0i
In[2]:= bessel_y_z(0, 40)
Out[2]= 0.09242859229058376 + 0i
In[3]:= bessel_k_scaled(0, 2000)
Out[3]= 0.02802320501460432 + 0i
In[4]:= sqrt(pi/4000)
Out[4]= 0.028024956081989644
In[5]:= bessel_i_scaled(0, 1000)
Out[5]= 0.012617240455891252 + 0i
In[6]:= hankel_h1_scaled(0.5, 3 + 700i)
Out[6]= -0.02127852045069675 - 0.02136990952385759i
In[7]:= bessel_k_scaled(40, 25)
Out[7]= 213915362812.9533 + 0i
In[8]:= bessel_y_scaled(0, 5000)
Out[8]= -0.009116740769643965 + -0i
In[9]:= bessel_k_scaled(0.5, 1e6)
Out[9]= 0.0012533141373154998 + 0i
In[10]:= sqrt(pi/2e6)
Out[10]= 0.0012533141373155003
\end{lstlisting}

On the real axis the $J$/$Y$ scaling is $e^0 = 1$, so \code{Out[1]}
\emph{is} $Y_0(40)$ and \code{Out[2]} is the same number computed the
old way. The true value is $0.12593641705826097$: the scaled form has
fifteen correct digits, the unscaled one none.

\code{Out[3]} and \code{Out[4]} are $e^xK_0(2000)$ and its leading
asymptotic $\sqrt{\pi/2x}$, agreeing to four digits with the gap equal
to $-1/(8x)$ to the digit. \textbf{Unscaled, that value does not
exist}: $K_0(2000)$ is about $10^{-870}$, below the smallest \code{f64}.
\code{Out[5]} is the same story upward, $I_0(1000) \approx e^{1000}$
being above the largest. \code{Out[6]} is $H^{(1)}$ seven hundred nepers
above the real axis, where $J$ and $Y$ are each about $e^{700}$.
\code{Out[9]}/\code{Out[10]} are the exact half-integer form
$e^{z}K_{1/2}(z) = \sqrt{\pi/2z}$ at $z = 10^{6}$.

\textbf{Order is handled by recurrence.} $K$, $Y$, $H^{(1)}$ and
$H^{(2)}$ all grow with order, so upward recurrence in $\nu$ is stable;
the expansion is used at a base order below 1 and stepped up.
\code{Out[7]} is order 40, where the vendored Cephes \code{kn} overflows
outright.

\textbf{Large order is a separate expansion.} A series in $1/z$ at fixed
order cannot reach $z$ below $\nu$, where $J$ is exponentially small and
was being built as the difference of two exponentially large Hankel
values --- at $\nu = 400.5$, $z = 240$ that gave an answer wrong by a
factor of $5\times10^{89}$. The remedy is a series in $1/\nu$: the
\textbf{Debye polynomials} and the uniform expansions of DLMF 10.19 and
10.41,
\[
  U_0(p) = 1,\qquad
  U_{k+1}(p) = \tfrac12 p^2(1-p^2)U_k'(p)
             + \tfrac18\int_0^p (1-5t^2)U_k(t)\,dt,
\]
which produce the small number directly. They are chosen automatically
by comparing error estimates, so nothing in the language changes; what
changes is that the answers are now right. Measured against Cephes over
orders 10 to 1000 and $z/\nu$ from 0.1 to 2, every value is within
$10^{-9}$ and most within $10^{-14}$ --- and at $\nu = 400.5$ it is
\emph{Cephes} that is looser, by $1.4\times10^{-9}$, which the
elementary $J$--$Y$ Wronskian adjudicates.

\textbf{When no method reaches a point these return an error} naming
both estimates. Since the large-order expansions arrived, most such
points instead fail for a different and better-stated reason: \emph{the
value is outside} \code{f64}. For $z$ well below $\nu$, $J$ is smaller
than the smallest double and $Y$ larger than the largest, and the error
says so and quotes the logarithm.

\textbf{The turning point has its own expansion.} Olver's uniform
Airy-type expansion (DLMF 10.20) replaces the elementary prefactor with
an Airy function, the exact local model of where the equation stops
oscillating:
\[
  J_\nu(\nu x) \sim \Big(\tfrac{4\zeta}{1-x^2}\Big)^{1/4}
    \Big[\tfrac{\mathrm{Ai}(\nu^{2/3}\zeta)}{\nu^{1/3}}
        \sum_k \tfrac{A_k(\zeta)}{\nu^{2k}}
      + \tfrac{\mathrm{Ai}'(\nu^{2/3}\zeta)}{\nu^{5/3}}
        \sum_k \tfrac{B_k(\zeta)}{\nu^{2k}}\Big].
\]
Its coefficients are built from the same Debye polynomials, and every
term in them is singular at $\zeta = 0$ with the singularities
cancelling exactly --- which is why near the turning point they come
from Taylor series in $1-x$ generated at 70 digits instead. Two of the
resulting constants are known independently and match:
$A_1(0) = -1/225$ and $\zeta'(1) = -2^{1/3}$.

An earlier version of this manual argued that 10.20 was unnecessary,
because across $z/\nu$ from 0.95 to 1.1 the other routes already gave
$10^{-14}$. \textbf{That measurement was too coarse and the conclusion
was wrong.} At $z/\nu = 0.85$ and $0.90$ --- not sampled --- the error
reached $1.4\times10^{-9}$. With 10.20 in the mix the whole band is
$10^{-14}$. Complex \emph{order} remains absent.

Measured, not asserted:
\begin{lstlisting}
cargo run -p special_functions --release --example bessel_scaled_accuracy
\end{lstlisting}

\subsubsection*{Complex order}

Every \code{\_nu} form above takes a \textbf{complex} order as well as a
real one. \code{bessel\_j\_nu(1 + 2i, 3)} is $J_{1+2i}(3)$; a real order
reaches exactly the routine it always did, so nothing that worked before
changes.

This was the last entry on the list of things chapter 10 did not cover,
and the one no expansion could have closed. The obstacle was not an
algorithm --- the ascending series is the same --- it was that
$1/\Gamma(\nu+k+1)$ had no meaning here for complex $\nu$. So the stage
is really a \textbf{complex gamma}, and the Bessel functions follow.

\begin{center}
\begin{tabular}{lp{7.4cm}}
\code{gamma\_z(z)} & $\Gamma$ at complex argument\\
\code{ln\_gamma\_z(z)} & still defined where $\Gamma$ has left \code{f64}\\
\code{rgamma\_z(z)} & $1/\Gamma$, \textbf{entire} --- exactly zero at the poles\\
\end{tabular}
\end{center}

\code{gamma\_z} is Stirling with argument shifting, not Lanczos. Lanczos
is a little faster and is rejected for the reason that has shaped this
crate: its coefficients are a \emph{table}, and the tables in
circulation are most often reproduced from \emph{Numerical Recipes},
whose licence this project will not inherit. Stirling needs no table ---
only the Bernoulli numbers, defined by a recurrence the crate states and
a test re-derives.

\begin{lstlisting}
In[1]:= gamma_z(0.5)
Out[1]= 1.7724538509055292 + 0i
In[2]:= sqrt(pi)
Out[2]= 1.7724538509055159
In[3]:= gamma_z(1 + 1i)
Out[3]= 0.4980156681183574 - 0.15494982830181092i
In[4]:= rgamma_z(-3)
Out[4]= 0 + 0i
In[5]:= ln_gamma_z(200)
Out[5]= 857.9336698258575 + 0i
In[6]:= bessel_j_nu(1 + 2i, 3)
Out[6]= 2.616967138257939 + 0.5245621513264704i
In[7]:= bessel_j_nu(0.5, 2)
Out[7]= 0.5130161365618277 + 0i
In[8]:= bessel_k_nu(1i, 2)
Out[8]= 0.09238545989039124 + 0i
In[9]:= bessel_k_nu(-1i, 2)
Out[9]= 0.09238545989039124 + 0i
\end{lstlisting}

\code{Out[1]} and \code{Out[2]} differ in the fourteenth digit, which is
the method's honest accuracy: reaching the Stirling regime from
$z = 0.5$ takes fourteen shifts and each rounds. \code{Out[4]} is the
reciprocal being \emph{entire}. \code{Out[5]} is a value
\code{gamma\_z(200)} refuses, $\Gamma(200)$ being about $10^{372}$.

\textbf{\code{Out[8]} and \code{Out[9]} are the result worth looking
at.} $K_{iy}(x)$ is real for real $y$ and real positive $x$ --- the
Macdonald function of imaginary order, which is why it appears as an
eigenfunction on the half line. Nothing in the implementation arranges
that: $K$ is built from two $I$s at $\pm\nu$ divided by $\sin\nu\pi$,
all three thoroughly complex. It comes out real, and even in the order,
because the mathematics says so.

\textbf{Two things change once the order is complex}, both because it
now sits in an exponent. $(z/2)^\nu = e^{\nu\ln(z/2)}$, so its
\emph{modulus} depends on $\arg z$ as well as $|z|$ --- a branch choice
is no longer a phase convention. And $\sin\nu\pi$ in the reflections
grows like $e^{\pi|\mathrm{Im}\,\nu|}$, so they get \textbf{better}
conditioned off the real axis: a whole-numbered real part is not a
special case unless the imaginary part is small too. The accuracy law
gains one term for the same reason:
\[
  \text{relative error} \sim 10^{-16}
    \exp\!\big(|z| - |\mathrm{Im}\,z|
      + \mathrm{Im}\,\nu \cdot \arg z\big),
\]
so complex order is free on the positive real axis and costs
$\mathrm{Im}\,\nu \cdot \arg z$ elsewhere.

\textbf{Beyond the series.} The $1/z$ asymptotics (DLMF 10.17, 10.40)
are expansions at \emph{fixed} order in which the order enters only as
$\mu = 4\nu^2$ --- polynomially --- so nothing in them assumes a real
order and they extend unchanged. The uniform expansions of DLMF 10.41
for $I$ and $K$ extend in the sector $|\arg\nu| < \pi/2$ (DLMF 10.41.5).
Both are offered alongside the series and chosen by comparing error
estimates, so the reach runs from $|z|\approx 25$ out past $|z| = 600$.

Two things the estimates could not see had to be measured. The $1/z$
expansion needs $|4\nu^2|$ small compared with $|z|$, not merely $|z|$
large --- at $\nu = 0.5+5i$ and $|z| = 22$ the actual error was
\textbf{2163 times} the truncation estimate --- so that is checked as a
validity condition rather than folded into a tolerance. And the term
DLMF 10.40.1 drops from $I$ is not of relative size $e^{-2\mathrm{Re}z}$
but $e^{\pi|\mathrm{Im}\nu| - 2\mathrm{Re}z}$: \textbf{complex order
makes it bigger}, by a factor this crate had no reason to carry until
now.

\textbf{The turning point is covered too.} Olver's Airy-type expansion
(DLMF 10.20) now runs at complex order, using the complex Airy functions
of \code{airy\_z}. It needed no new mathematics: near $x = 1$ every
ingredient --- $\zeta$, the prefactor, and the coefficients $A_k$, $B_k$
--- is already a Taylor series in $w = 1-x$, generated at 70 digits, and
a Taylor series does not care whether its variable is real. What did
\emph{not} continue to complex $x$ are the closed forms outside that
neighbourhood, where 10.20 is the wrong tool anyway.

\textbf{The Debye region is covered too}, on both sides of the turning
point and at complex order as well as real. With $t = \sqrt{1-x^2}$,
$\alpha = \ln((1+t)/x)$, $q = 1/t$ and $x = z/\nu$,
\[
  F_\pm(\nu,x) = \frac{e^{\pm\nu(t-\alpha)}}{\sqrt{2\pi\nu t}}
    \sum_k \frac{(\pm1)^k U_k(q)}{\nu^k}
\]
reads two ways. For $|x| < 1$ these are DLMF 10.19.3/4 directly,
$J = F_+$ and $Y = -2F_-$. For $|x| > 1$ --- the \textbf{oscillatory}
region, where $t$ turns imaginary --- they continue into the Hankel
functions instead, $H^{(1)} = 2F_+$ and $H^{(2)} = 2iF_-$, so
$J = F_+ + iF_-$ and $Y = -iF_+ - F_-$. That flip is the Stokes
phenomenon, and \textbf{the constants were identified by experiment
rather than transcribed}: continuing $F_+$ past $x = 1$ and dividing by
each of $J$, $Y$, $H^{(1)}$, $H^{(2)}$ in turn showed
$F_+/H^{(1)} = \tfrac12$.

That band had no method at all, at real order as well as complex.
\code{bessel\_y\_z(20, 60)} used to return $10^8$; it is now correct to
$2\times10^{-14}$, and across orders 8 to 150 and $z/\nu$ from 2 to 30
the whole band is $10^{-14}$ or better.

What remains is a sliver: $|\nu|$ between about 4 and 8 with $|z|$ a few
times larger, where the $1/z$ route is refused for $|4\nu^2|$ and the
Debye one for being a $1/\nu$ series at too small an order. Both
refusals are deliberate and measured. There is also a narrow ridge near
$z/\nu \approx 1.3$ where $Y$ reaches only about $10^{-7}$ at moderate
order.

\subsubsection*{Airy functions, complex argument}

\code{airy\_z(z)} returns all four at once --- \code{[Ai, Ai', Bi, Bi']}
--- because the routine computes them together and a caller who wants a
Wronskian or a turning-point boundary condition wants all four.

\begin{lstlisting}
In[1]:= airy_z(0)
Out[1]= [0.3550280538878172 + 0i, -0.2588194037928068 + 0i,
         0.6149266274460007 + 0i, 0.4482883573538264 + 0i]
In[2]:= airy_z(2 - 3i)
Out[2]= [0.008104457809530868 - 0.131178382604566i,
         0.09665817903311252 + 0.23198718538548577i,
         -0.3963682550403918 + 0.5697309129559497i,
         0.34945767192946653 + 1.105328588933856i]
\end{lstlisting}

\code{Out[1]} is the closed form at the origin:
$\mathrm{Ai}(0) = 3^{-2/3}/\Gamma(2/3)$,
$\mathrm{Ai}'(0) = -3^{-1/3}/\Gamma(1/3)$ and
$\mathrm{Bi}(0) = \sqrt3\,\mathrm{Ai}(0)$, the last of which you can read
straight off the numbers.

\textbf{Three regimes, one connection formula.} The ascending series
(DLMF 9.4) up to $|z|\approx 6$; the asymptotic expansions (DLMF 9.7.5,
9.7.6) in $\zeta = \tfrac23 z^{3/2}$ for $|\arg z| \le 2\pi/3$; and
nearer the negative real axis the connection
$\mathrm{Ai}(z) + \omega\mathrm{Ai}(\omega z)
 + \omega^2\mathrm{Ai}(\omega^2 z) = 0$ with $\omega = e^{2\pi i/3}$,
which rotates $\arg z \sim \pi$ to $\arg \sim \pm\pi/3$ where the
expansion is at its best --- the same trick the Bessel functions use for
their own cut, and it works for the same reason.

$\mathrm{Bi}$ needs no expansion of its own past the series: its
asymptotic is valid only for $|\arg z| < \pi/3$, and
$\mathrm{Bi}(z) = e^{i\pi/6}\mathrm{Ai}(ze^{2\pi i/3})
 + e^{-i\pi/6}\mathrm{Ai}(ze^{-2\pi i/3})$ gives it everywhere.

\textbf{How it is checked.} The Wronskian
$\mathrm{Ai}\,\mathrm{Bi}' - \mathrm{Ai}'\,\mathrm{Bi} = 1/\pi$ is
exact, elementary, and free of both the argument and any transcendental
function on the right. Across $|z|$ from 0.1 to 200 and the full range
of $\arg z$ the residual is \textbf{$10^{-11}$ or better}, except in one
band: between $|z| = 3$ and $10$ it is
$10^{-9}$ --- the crossover, where $\mathrm{Ai}$ is exponentially
recessive. The test states that band separately rather than loosening
the whole bound to accommodate it. Also checked: the connection
formula away from where it is used, conjugation symmetry, the defining
equation $\mathrm{Ai}'' = z\,\mathrm{Ai}$ by finite difference, and
Cephes on the real axis.

Past $|z|\approx 90$ off the real axis the \emph{dominant} solution
leaves \code{f64} --- it grows like $e^{2|z|^{3/2}/3}$ --- and
\code{airy\_z} says so rather than returning an infinity.

\textbf{A wrinkle worth knowing about lists.} The bracket literal is
overloaded: \code{[a,b,c]} is a \emph{vector}, \code{[a,b,c,d]} is a
\emph{quaternion}, and any other length is a list. That is convenient
for physics and would be a nuisance here, so all three shapes are
accepted anywhere a numeric list is expected --- a $4\times4$ matrix
whose rows are 4-element brackets works exactly as you would expect, and
\code{norm\_assoc\_legendre\_p} never lectures you about quaternions. A
$3\times3$ matrix value is likewise accepted directly wherever a matrix
argument is wanted.

\subsection{Comparisons}\label{sec:cmp}

\code{<}, \code{<=}, \code{>}, \code{>=}, \code{==}, \code{!=} compare
two numbers and yield \textbf{1 for true, 0 for false}. There is no
boolean type, and that is the point:

\begin{lstlisting}
In[1]:= 3 < 5
Out[1]= 1
In[2]:= (0.5 > 0) * (0.5 < 1)
Out[2]= 1
In[3]:= (2 > 0) * (2 < 1)
Out[3]= 0
\end{lstlisting}

\code{(x > a) * (x < b)} is an \textbf{indicator function} for the
interval $(a,b)$. Multiply it by a height and you have a rectangular
barrier; add two and you have a double barrier. This is how a piecewise
potential is written --- see \S\ref{sec:qm} and Example 17.

Comparisons sit at the \textbf{lowest} precedence, below \code{+} and
\code{-}, so \code{x + 1 > 2} groups as \code{(x + 1) > 2}. They are
left associative, so \code{a < b < c} means \code{(a < b) < c}: legal,
and almost certainly not what you meant.

\code{=} is still assignment; \code{==} is equality. Ordering is defined
for numbers only; \code{==} and \code{!=} also accept two vectors or two
quaternions, because asking whether two positions coincide is reasonable
while ordering them is not.

\textbf{NaN compares false to everything, including itself}, so
\code{x != x} is the idiomatic NaN test:

\begin{lstlisting}
In[4]:= 0/0 != 0/0
Out[4]= 1
\end{lstlisting}

That is not special-cased --- it falls out of IEEE-754.

\subsection{Complex numbers}\label{sec:complex}

A number with an \code{i} suffix is imaginary, so \code{2 + 3i} needs no
complex literal syntax of its own --- it is ordinary addition of a real
and an imaginary:

\begin{lstlisting}
In[1]:= 3i
Out[1]= 0 + 3i
In[2]:= 2 + 3i
Out[2]= 2 + 3i
In[3]:= (2 + 3i) * (2 - 3i)
Out[3]= 13 + 0i
In[4]:= 1 / (0 + 1i)
Out[4]= 0 - 1i
\end{lstlisting}

\code{+}, \code{-}, \code{*} and \code{/} all accept complex operands,
and reals promote automatically. Note \code{In[3]}: the result stays
typed as complex and displays \code{13 + 0i} rather than collapsing to
\code{13}. That is deliberate --- the type you get out should not depend
on whether a particular cancellation happened to be exact.

The suffix only binds when the \code{i} is not followed by another
identifier character, so \code{2intercept} still lexes exactly as it did
before complex numbers existed.

This is what makes the complex Crank--Nicolson solvers reachable.
\code{solve\_tridiag\_c} and \code{solve\_cyclic\_tridiag\_c} take the
same arguments as the real solver and accept a mix --- a real band with
a complex diagonal and a real right-hand side is fine, since reals
promote:

\begin{lstlisting}
In[5]:= solve_tridiag_c([0, 1, 1], [1i, 1i, 1i], [1, 1, 0], [1, 0, 0])
Out[5]= [0 - 0.6666666666666666i, 0.33333333333333337 + 0i, 0 + 0.3333333333333333i]
\end{lstlisting}

Over the machine bridge (JupyterLab), a complex value is reported as the
two-element array \code{[re, im]}, matching how \code{sph\_harm} reports
its value --- JSON has no complex type.

\section{Command semantics}

\subsection{\kw{new} --- create an object}\label{sec:new}

\begin{lstlisting}
NEW POINT    { field = expr, ... }
NEW SPHERE   { field = expr, ... }
NEW CUBOID   { field = expr, ... }        (alias: NEW CUBE)
NEW TORUS    { field = expr, ... }
NEW DISK     { field = expr, ... }        (alias: NEW DISC)
NEW CYLINDER { field = expr, ... }
NEW DUMBBELL { field = expr, ... }        (alias: NEW DUMBELL)

NEW <shape> AS <name> { field = expr, ... }   (register a user name)
\end{lstlisting}

Creates one \emph{physical object} and prints its handle (\code{obj0},
\code{obj1}, \ldots, numbered by position). Defaults: mass 1, at the
origin, at rest, no charge; sphere radius 1; cuboid half-extents
\code{[1,1,1]}; torus \code{ring\_radius} 1, \code{tube\_radius} 0.25;
disk radius 1; cylinder radius 0.5, \code{half\_height} 1; dumbbell
\code{m1} $= 1$, \code{m2} $= 1$, \code{m\_rod} $= 0.5$, \code{r1}
$= 0.25$, \code{r2} $= 0.25$, \code{rod\_radius} $= 0.1$,
\code{length} $= 1$ (total mass 2.5).

Initializer fields: \code{mass, charge, position, velocity, momentum,
orientation, angular\_velocity, angular\_momentum, radius} (spheres,
disks, cylinders), \code{half\_extents} (cuboids), \code{ring\_radius,
tube\_radius} (tori --- or the \code{inner\_radius} +
\code{outer\_radius} pair, see below), \code{height, half\_height}
(cylinders; \kw{height} is the full height $= 2\cdot$\code{half\_height}),
\code{m1, m2, m\_rod, r1, r2, rod\_radius, length} (dumbbells, see
below), \code{inertia\_tensor, inverse\_inertia\_tensor,
magnetic\_moment\_tensor, force, torque, id, inverse\_mass}.

A torus can be sized two ways: directly (\code{ring\_radius} $c$ = the
radius of the circle traced by the tube's center, \code{tube\_radius}
$a$ = the tube's own radius) or by the \textbf{\code{inner\_radius} +
\code{outer\_radius} pair} (the hole's radius and the outermost
radius: $\mathrm{inner} = c - a$, $\mathrm{outer} = c + a$). Inside a
\kw{new} initializer list the torus geometry is \textbf{deferred}: the
four radius fields are collected and resolved \emph{once}, at the end
of the list --- \code{ring\_radius}/\code{tube\_radius} are applied
first, then \code{inner\_radius}/\code{outer\_radius} override the
derived pair --- and the result is validated once, against the
\emph{final} values (a bad pair fails with \code{torus needs 0 <=
inner < outer (got inner = 2, outer = 1)}). Giving both members of the
pair is therefore \textbf{genuinely order-independent}:
\code{\{ inner\_radius = 1, outer\_radius = 2 \}} yields ring 1.5,
tube 0.5 whichever comes first, and a pair that shrinks the default
torus --- \code{\{ outer\_radius = 0.5, inner\_radius = 0.2 \}} ---
works in either order too (checking each write against the
half-updated default would have refused one of the two orders).
\code{inner\_radius = 0} is legal --- that is the \textbf{horn torus},
whose tube touches the axis --- on \kw{new} and on \kw{set} alike.

\textbf{The dumbbell} is ONE rigid body: two solid spheres joined by
a solid rod. Its seven part fields are \code{m1}, \code{m2},
\code{m\_rod} (the two sphere masses and the rod mass), \code{r1},
\code{r2} (the sphere radii), \code{rod\_radius} (alias
\code{rod\_r}) and \code{length} (alias \code{len} --- the distance
between the sphere centers). Like the torus radius pair, the part
fields are \textbf{deferred}: they are collected during the
initializer list and resolved \emph{once} at its end, so they are
order-independent and validated once against the final values. The
constructor places the local origin at the composite \textbf{center
of mass} --- the sphere centers sit at
$z_1 = -(m_2 + m_{\mathrm{rod}}/2)\,L/M$ and
$z_2 = +(m_1 + m_{\mathrm{rod}}/2)\,L/M$ on the body $z$-axis, which
makes $m_1 z_1 + m_2 z_2 + m_{\mathrm{rod}}(z_1+z_2)/2 = 0$ an
identity --- so \code{position} is the COM, exactly as for every
other shape. The inertia tensor is the exact composite (solid-sphere
$\tfrac25 m r^2$ terms with parallel-axis shifts, plus the rod), and
the surface is the exact union of the three parts. It is the
simulator's first non-centrally-symmetric shape; collisions against
every other shape decompose exactly over its parts (see
\code{collision\_detection.pdf}). The part fields remain readable
\emph{and writable} after creation (\S5.2).

\textbf{\kw{as} registers a user name for the object.} The name is a
string literal (\code{NEW SPHERE AS "ball"}) or a bare identifier
(\code{NEW SPHERE AS ball}); a bare identifier is first resolved
against a function parameter or \kw{let} variable holding a string
--- so a function can name the object it creates from its argument
(\S\ref{sec:userdefs}, Example 14) --- and otherwise the
identifier's own spelling is the name. The reply shows both handles:
\code{obj0 as ball}. Named paths then work everywhere a positional
path does --- \code{ball.mass}, \code{dumbell0.position.x}, in
\kw{set}/\kw{get} and in bare expressions. Names are
case-insensitive identifiers; \code{system}/\code{sys} and the
positional spellings are refused (\code{`system` is reserved},
\code{`obj7` is reserved for positional paths}), and so are
duplicates (\code{the name `d` already refers to obj0 --- DEL it or
pick another name}). The registry follows every renumbering:
\kw{del} removes the deleted object's name and shifts the names of
higher-numbered objects down with their objects, and the \kw{box}
commands do the same when walls come and go.

Four guarantees make initializers forgiving:
\begin{enumerate}
\item \textbf{Order does not matter for velocities.} \code{velocity} and
      \code{angular\_velocity} are applied \emph{after} mass and inertia
      are final, so \code{\{ velocity = [1,0,0], mass = 2 \}} still
      yields momentum \code{[2,0,0]}.
\item \textbf{Inertia is computed for you.} For every extended shape
      the inertia tensor is recomputed from the final mass and shape
      (solid sphere $\tfrac{2}{5}mr^2$; cuboid $\tfrac{m}{3}(h_y^2+h_z^2)$
      etc.; torus $I_z = m(c^2+\tfrac34 a^2)$,
      $I_{xy} = m(\tfrac12 c^2+\tfrac58 a^2)$; disk
      $I_z = \tfrac12 ma^2$, $I_{xy} = \tfrac14 ma^2$; cylinder
      $I_z = \tfrac12 mr^2$, $I_{xy} = m(3r^2+4h^2)/12$ with $h$ the
      half-height) --- \emph{unless} you supplied
      \code{inertia\_tensor} yourself, in which case yours is kept.
\item \textbf{Coupled fields stay consistent} (see \kw{set} below).
\item \textbf{\kw{new} is transactional.} If any initializer --- or
      the final geometry validation --- fails, the half-built object
      is removed before the error is reported: a failing \kw{new}
      never leaves a ghost behind, and \code{system.count} and the
      \code{objN} numbering are exactly what they were before the
      command.
\end{enumerate}

\subsection{\kw{set} / \kw{get} --- write and read any field}

A \emph{path} is \code{objN.field}, \code{system.field},
\code{contactK.field} (\S\ref{sec:collide}) or \code{name.field} for
a name registered with \kw{new} \ldots{} \kw{as}
(\S\ref{sec:new}), optionally followed by a component
\code{.x .y .z} (vectors) or \code{.w .x .y .z} (quaternions).
Component writes are safe read-modify-write operations through the
full field's get/set pair. Every object also has six
\textbf{component shorthands}: \code{.x .y .z} read and write the
position components and \code{.vx .vy .vz} the velocity components
--- \code{ball.x} is \code{ball.position.x}, \code{SET ball.vx = 2}
is a velocity write.

\textbf{Object fields} (R = readable, W = writable):

\begin{longtable}{lllp{6.4cm}}
\textbf{field} & \textbf{type} & \textbf{R/W} & \textbf{meaning}\\\hline
\code{id} & number & RW & user label (no effect on physics)\\
\code{mass} & number & RW & writing also updates \code{inverse\_mass} ($m\le 0 \Rightarrow 1/m := 0$)\\
\code{inverse\_mass} & number & RW & writing back-computes \code{mass}; $0$ makes the body \textbf{static}\\
\code{charge} & number & RW & Coulombs\\
\code{position}, \code{pos} & vec3 & RW & world position\\
\code{velocity}, \code{vel} & vec3 & RW & \emph{derived}: reads $p/m$, writes $p := m v$\\
\code{momentum} & vec3 & RW & canonical stored linear state\\
\code{orientation} & quat & RW & renormalized on write\\
\code{angular\_velocity} & vec3 & RW & derived: $\omega = R\,I^{-1}R^{\mathsf T}L$; writes $L := R\,I\,R^{\mathsf T}\omega$\\
\code{angular\_momentum} & vec3 & RW & canonical stored angular state (spin)\\
\code{inertia\_tensor} & mat3 & RW & body frame; write updates the inverse (singular $\to$ zero inverse = cannot rotate)\\
\code{inverse\_inertia\_tensor} & mat3 & RW & write back-computes the tensor\\
\code{magnetic\_moment\_tensor} & mat3 & RW & torque $\tau = (R M R^{\mathsf T})B$\\
\code{radius} & number & RW & spheres, disks, cylinders; write recomputes inertia and \textbf{keeps the shape family} (a disk stays a disk, a cylinder keeps its height); a torus is \textbf{refused} --- \code{obj6 is a torus \textemdash{} set ring\_radius/tube\_radius or inner\_radius/outer\_radius instead of radius} --- rather than silently becoming a sphere\\
\code{half\_extents} & vec3 & RW & cuboids only; write recomputes inertia\\
\code{ring\_radius}, \code{tube\_radius} & number & RW & tori only; write recomputes inertia\\
\code{inner\_radius}, \code{outer\_radius} & number & RW & tori only, derived ($\mathrm{inner} = c-a$, $\mathrm{outer} = c+a$); writing one holds the other fixed; \code{inner\_radius} accepts $\ge 0$ (0 = the horn torus; the other radii must be $> 0$)\\
\code{height}, \code{half\_height} & number & RW & cylinders only; \code{height} is the full height ($= 2\cdot$\code{half\_height}); write recomputes inertia\\
\code{m1}, \code{m2}, \code{m\_rod} & number & RW & dumbbells only: the two sphere masses and the rod mass; writing rebuilds the body --- total mass, the COM offsets and the inertia tensor all follow (position is untouched; the stored momenta are preserved --- momentum-canonical, like every mass write --- so the velocities rescale with the new mass and inertia)\\
\code{r1}, \code{r2}, \code{rod\_radius} & number & RW & dumbbells only: the sphere radii and the rod radius (alias \code{rod\_r}); same rebuild on write\\
\code{length} & number & RW & dumbbells only: distance between the sphere centers (alias \code{len}); same rebuild on write\\
\code{x}, \code{y}, \code{z} & number & RW & every object: shorthands for \code{position.x/.y/.z}\\
\code{vx}, \code{vy}, \code{vz} & number & RW & every object: shorthands for \code{velocity.x/.y/.z}\\
\code{boundary}, \code{shape} & string & R & description of the shape\\
\code{kinetic\_energy}, \code{energy} & number & R & $\tfrac12 m|v|^2 + \tfrac12\,\omega\!\cdot\!L$\\
\code{force} & vec3 & RW & constant external force during \kw{step}/\kw{run}\\
\code{torque} & vec3 & RW & constant external torque\\
\code{restitution} & number & RW & collision bounciness $e \in [0,1]$ (default 1 = elastic; a pair uses $\min(e_i, e_j)$)\\
\end{longtable}

\textbf{System fields}:

\begin{center}
\begin{tabular}{lllp{6.2cm}}
\textbf{field} & \textbf{type} & \textbf{R/W} & \textbf{meaning}\\\hline
\code{g\_constant}, \code{g} & number & RW & Newton's $G$ (default 1)\\
\code{softening} & number & RW & Plummer length $\varepsilon$ (default $10^{-6}$); forces use $(r^2+\varepsilon^2)^{3/2}$\\
\code{uniform\_gravity}, \code{gravity} & vec3 & RW & constant acceleration field\\
\code{e\_field} & vec3 & RW & uniform $E$; force $qE$\\
\code{b\_field} & vec3 & RW & uniform $B$; force $q\,v\times B$, torque $(RMR^{\mathsf T})B$\\
\code{rtol}, \code{atol} & number & RW & CVODE tolerances ($10^{-10}$, $10^{-12}$)\\
\code{time}, \code{t} & number & RW & current simulation time\\
\code{method} & string & R & current integrator\\
\code{count}, \code{n} & number & R & number of objects\\
\code{collide} & number & R & 1 when collision detection is on (switch with \kw{collide on/off})\\
\code{contacts} & number & R & contacts recorded by the last \kw{step}/\kw{run}\\
\code{collisions} & number & R & running total of resolved impulses\\
\code{restitution\_threshold} & number & RW & approach speeds below this bounce with $e=0$ (default $10^{-3}$)\\
\code{contact\_slop} & number & RW & tolerated overlap before positional projection (default $10^{-9}$)\\
\code{box} & number & R & inner side length of the rigid bounding box (\S\ref{sec:box}); 0 = none. Create/remove with \kw{box}\\
\end{tabular}
\end{center}

\subsection{\kw{step} and \kw{run} --- integrate (always via SUNDIALS)}

\code{STEP dt} advances time by \code{dt}. \code{RUN t STEPS n} advances
by \code{t}, stopping at \code{n} evenly spaced output points (default
10). Both accept full expressions (\code{run 2 * pi steps 8} is legal).
The reply summarizes the run; \code{|dE/E|} is the relative change in
total energy --- your built-in sanity check. \textbf{All integration is
performed by the pure-Rust SUNDIALS solvers}; there is no hand-rolled
stepper.

\subsection{\kw{method} --- choose the integrator}

\begin{lstlisting}
METHOD ADAMS               (default; CVODE Adams-Moulton, adaptive)
METHOD BDF                 (CVODE BDF -- stiff problems, fast gyration)
METHOD SPRK <table> [dt]   (ARKODE symplectic, fixed step; default dt 0.01)
\end{lstlisting}

Table names may be abbreviated: \code{leapfrog\_2\_2} becomes
\code{ARKODE\_SPRK\_LEAPFROG\_2\_2}. Useful tables: \code{EULER\_1\_1},
\code{LEAPFROG\_2\_2} ($\equiv$ classic velocity Verlet),
\code{MCLACHLAN\_2\_2/3\_3/4\_4/5\_6}, \code{RUTH\_3\_3},
\code{YOSHIDA\_6\_8}.

SPRK requires a \emph{separable} system: point-like translation only.
If anything velocity-dependent or rotational is active (a
\code{b\_field}, a magnetic tensor, an external torque, a spinning
rigid body), \kw{run} refuses with an error \emph{naming the offending
feature}.

\subsection{Observables, bookkeeping, help}

\begin{center}
\begin{tabular}{lp{10cm}}
\textbf{command} & \textbf{prints}\\\hline
\kw{energy} & kinetic + softened pairwise potential + uniform-field potentials\\
\kw{com} & system center of mass\\
\kw{momentum} & total linear momentum\\
\kw{angmom} & total angular momentum about the origin (orbital + spin)\\
\kw{laplace} $n$ & Laplace--Runge--Lenz vector of object $n$ about the center of mass, $k = G\,M_{\mathrm{total}}$\\
\kw{list} & one line per object\\
\kw{del} $n$ & removes object $n$ (\textbf{later objects renumber!})\\
\kw{reset} & wipes everything (not to be confused with \code{SCENE RESET}, which re-initializes only the scene window's playback copy --- see \code{scene\_info.pdf})\\
\kw{help} & the quick-reference card\\
\end{tabular}
\end{center}

\subsection{\kw{scene} --- the graphical scene window}\label{sec:scene}

\kw{scene create} starts a tiny web server \emph{inside} posim (pure
Rust standard library --- no external dependencies) and opens a page in
your web browser. That page \textbf{is} the scene window: it draws
every simulator entity on a 3-D canvas, has a \textbf{toolbar} (Start,
Pause, Stop, Reverse, a permanent Reset button, single-step, a dt box,
zoom, view reset, grid/trails/labels toggles, help) and a
\textbf{status bar} (connection light, playback mode, simulation time,
dt, total energy, body count, hidden count, history depth, camera
readout, frames per second). Inside the window:

\begin{center}
\begin{tabular}{lp{8.4cm}}
\textbf{gesture} & \textbf{effect}\\\hline
$\leftarrow$ $\rightarrow$ (arrow keys) & translate the view left / right\\
$\uparrow$ $\downarrow$ & translate the view up / down\\
left-click + drag & rotate (orbit) around the scene\\
mouse wheel & zoom in / out\\
\code{+} / \code{-} & zoom in / out from the keyboard\\
shift-drag or right-drag & translate (pan) with the mouse\\
Space & start / pause playback\\
\code{R}, \code{G}, \code{T}, \code{L}, \code{H} & reset view, toggle grid / trails / labels, help\\
\end{tabular}
\end{center}

The same controls are scriptable from the notebook:

\begin{longtable}{lp{9.4cm}}
\textbf{command} & \textbf{effect}\\\hline
\endhead
\kw{scene create} \code{[port]} & open the window (all entities shown). Default port: one chosen by the OS; give a number (0--65535) to pin it. Prints the URL. A second \kw{create} is harmless --- it just reminds you of the URL.\\
\kw{scene close} (= \kw{destroy}) & shut the server down; every window disconnects\\
\kw{scene translate} \code{dx dy [dz]} & move the camera's look-at point by (dx, dy, dz) world units (\code{dz} defaults to 0)\\
\kw{scene rotate} \code{dyaw dpitch} & orbit the camera: yaw (azimuth) and pitch (elevation) in \textbf{degrees}; pitch is clamped to $\pm89^\circ$\\
\kw{scene zoom in} / \kw{out} / \code{f} & zoom by $1.25\times$, by $1/1.25$, or by any factor $f > 0$ ($f > 1$ zooms in)\\
\kw{scene hide} \code{[n\textbar ALL]} & hide object n, or everything (bare \kw{hide} = \kw{hide all})\\
\kw{scene show} \code{[n\textbar ALL]} & undo \kw{hide}\\
\kw{scene refresh} & copy the notebook's current system into the window (see below) and clear playback history\\
\kw{scene redraw} & re-send the complete scene description to every window (forces a full redraw)\\
\kw{scene start} & begin time-stepped evolution, forward\\
\kw{scene pause} & freeze; \kw{start} resumes, history is kept\\
\kw{scene stop} & halt \textbf{and clear the recorded history}\\
\kw{scene reverse} & play \textbf{backward in time} through the recorded history\\
\kw{scene reset} & \textbf{re-initialize the playback}: every mutable value and the time return to their initial values --- the state last synced at \kw{create}/\kw{refresh}, restored bit-identically --- history and the step counter clear, and the mode returns to \emph{stopped}; \kw{start} then re-runs the simulation from the beginning. The window's permanent toolbar \textbf{Reset} button does exactly the same (both call one primitive)\\
\kw{scene set\_time\_step} \code{dt} & set the playback time step (must be positive and finite)\\
\kw{scene status} & a four-line report: URL + connected windows; mode/t/dt/steps/history; entities + hidden list; camera\\
\kw{scene events} & print (and clear) the asynchronous messages the window has sent: errors, connect/disconnect notices, toolbar actions, data requests\\
\end{longtable}

\textbf{The playback copy.} The window animates its \emph{own
synchronized copy} of your system, evolved by a background thread ---
your notebook state does \textbf{not} move while the animation runs, so
\code{GET obj0.position} still answers instantly and exactly. The copy
is taken at \kw{create} and again at every \kw{refresh}. All forward
stepping goes through the same SUNDIALS integrators as
\kw{step}/\kw{run} --- there is no separate physics engine in the
window.

\textbf{Playback is a four-state machine} (\emph{stopped} $\to$
\emph{running} $\rightleftarrows$ \emph{paused}, plus
\emph{reversing}): \kw{stop} clears history and a later \kw{start}
begins fresh; \kw{pause} keeps history so both \kw{start} (forward) and
\kw{reverse} (backward) can continue from the freeze point.
\kw{reset} is stronger than \kw{stop}: besides clearing history it puts
every value back to its initial state, so the next \kw{start} replays
the simulation from the very beginning rather than continuing from
wherever playback halted. The reply confirms it:

\begin{lstlisting}
In[6]:= scene reset
Out[6]= scene playback reset to its initial state (t = 0, 1 entity); Start runs the simulation again from the beginning
\end{lstlisting}

\textbf{How REVERSE works.} While running forward, the playback thread
records a snapshot of the whole system before every step (a ring
buffer, at most 20\,000 frames). \kw{reverse} replays those snapshots
newest-first, which is an \emph{exact} rewind --- bit-for-bit the
states you already visited. When the buffer runs out it pauses and
sends an event (\code{reverse: reached the beginning of recorded
history --- paused}). Reversing with no history at all is refused with
an error.

\textbf{Asynchronous events.} The window can talk back at any time ---
it reports JavaScript errors, connections, disconnections, and every
toolbar action. These messages queue up inside posim; \kw{scene events}
drains the queue (up to 1000 messages are kept). In JupyterLab they
also appear \emph{by themselves}, without you asking --- see
\S\ref{sec:jlab}.

\subsection{\kw{collide} and \kw{contacts} --- rigid-body collisions}
\label{sec:collide}

\begin{lstlisting}
COLLIDE            (report: on/off, collidable pairs, impulses so far)
COLLIDE ON         (default - collisions are detected in EVERY scene)
COLLIDE OFF        (pure point-mass/gravity studies)
CONTACTS           (list every contact of the last STEP/RUN)
\end{lstlisting}

When collisions are on (the default) and the system contains at least
one collidable pair (spheres, cuboids, or a point against either ---
two points cannot collide), \kw{step} and \kw{run} detect impacts
\textbf{during} the time step by SUNDIALS \emph{event rootfinding}: the
integrator itself lands on the instant where a pair's signed separation
crosses zero, interpolated to solver precision. At that instant an
impulse acts along the \textbf{contact normal} --- the line of the
action--reaction force pair --- and integration continues. Nothing
tunnels, and the reply reports what happened.

Every contact of the last \kw{step}/\kw{run} is recorded and exposed to
the rest of the simulation through read-only \code{contactK} paths:

\begin{center}\small
\begin{tabular}{llp{8.0cm}}
\textbf{field} & \textbf{type} & \textbf{meaning}\\\hline
\code{contactK.i}, \code{contactK.j} & number & the colliding pair (indices, $i < j$)\\
\code{contactK.t} & number & the event time (the root the solver landed on)\\
\code{contactK.point} & vec3 & world-space contact point\\
\code{contactK.normal} & vec3 & \textbf{unit normal, from \code{obji} toward \code{objj}} --- the action--reaction line; \code{objj} receives $+J\hat n$, \code{obji} receives $-J\hat n$\\
\code{contactK.depth} & number & penetration depth at the event ($\approx 0$)\\
\code{contactK.rel\_vel\_n} & number & approach speed along the normal (negative = approaching)\\
\code{contactK.impulse} & number & scalar impulse magnitude $J$\\
\end{tabular}
\end{center}

Response physics (full derivation in \code{collision\_detection.pdf}):
per-object \code{restitution} $e \in [0,1]$ (1 = elastic default, 0 =
perfectly plastic), pair-combined as $\min(e_i, e_j)$; impulse
$J = -(1+e)\,(\mathbf v_{\mathrm{rel}}\cdot\hat n)\,/\,(\hat n^{\!\top}
K \hat n)$ with the effective-mass matrix $K$ including the angular
terms, applied through the canonical momentum/angular-momentum setters;
static bodies (\code{inverse\_mass} $= 0$) are immovable walls;
approach speeds below \code{system.restitution\_threshold} settle
without bounce; leftover overlap beyond \code{system.contact\_slop} is
projected out. \kw{scene} windows draw each contact normal as an arrow
at the contact point (\textbf{Contacts} toolbar button or the \code{C}
key).

\subsection{\kw{box} --- the rigid bounding box}
\label{sec:box}

\begin{lstlisting}
BOX <size>         (create: six static walls enclosing an inner
                    size x size x size cube centered on the origin)
BOX OFF            (remove the box - its six walls are deleted)
BOX                (bare: report status)
\end{lstlisting}

\kw{box} \code{<size>} builds a closed rigid room out of \textbf{six
ordinary cuboid objects} --- wall slabs behind the planes
$x = \pm\,\mathrm{size}/2$, $y = \pm\,\mathrm{size}/2$,
$z = \pm\,\mathrm{size}/2$ (half-thickness $\mathrm{size}/4$, centers
at $\pm 3\cdot\mathrm{size}/4$, cross-sections wide enough to cover
the corners) --- and prints their handles:

\begin{lstlisting}
In[2]:= box 4
Out[2]= box: inner size 4 x 4 x 4 — six static walls obj0, obj1, obj2, obj3, obj4, obj5 with inverse_mass = 0 (infinitely massive); objects collide elastically off the inside faces
\end{lstlisting}

\textbf{Infinite mass is exact, not approximate.} The equations of
motion never use the mass itself --- only the \emph{inverse} mass:
velocity is $v = p\,m^{-1}$, and the collision impulse divides by
$\hat n^{\!\top}K\hat n = m_i^{-1} + m_j^{-1} + (\text{angular terms})$
(\S\ref{sec:collide}). Each wall is created with
\code{inverse\_mass} $= 0$ and a zero inverse inertia tensor, so it
contributes exactly 0 to every impulse denominator and receives no
state writes: bodies bounce off the inside faces elastically while the
walls stay \textbf{bit-identically at rest}. One measurable
consequence: system \textbf{momentum is not conserved} inside a box
--- the infinitely massive walls absorb it without moving
(Example 13). \kw{list} tags every wall
\code{[wall: static, inverse\_mass=0]} (and shows \code{mass=0} ---
the canonical stored quantity for a static body is the inverse).

Bookkeeping:
\begin{itemize}
\item \code{GET system.box} reads the inner side length; 0 means no
      box. The path is read-only --- the box is created and removed by
      the command.
\item A second \kw{box} \code{<size>} \textbf{replaces} the existing
      box: the old walls are removed first and the reply notes
      \code{(replacing the previous box)}.
\item The walls are ordinary objects with ordinary indices: \kw{del}
      on a wall deletes it like any other object \textbf{and dissolves
      the box} --- \code{system.box} drops to 0 (five walls no longer
      enclose anything) --- but the surviving slabs \textbf{stay
      tracked}: \kw{list} keeps their
      \code{[wall: static, inverse\_mass=0]} tag, and bare \kw{box}
      reports \code{box: dissolved (a wall was deleted; 5 tracked
      slab(s) remain \textemdash{} BOX <size> replaces them, BOX OFF
      removes them)}. The next \kw{box} \code{<size>} removes the
      survivors \emph{before} building the new box (the reply notes
      \code{(replacing the previous box)}) and \kw{box off} simply
      deletes them --- either way no orphan slab leaks. Wall indices
      are tracked through every \kw{del} renumbering, so deleting a
      non-wall object never confuses the box.
\item \kw{box off} replies \code{box removed (6 wall(s) deleted,
      indices renumbered)} --- or \code{box: none} if there was
      nothing to remove; bare \kw{box} reports the size and wall
      handles, or \code{box: none (BOX <size> creates one)}.
\item A \kw{scene} window draws the box as a \textbf{dashed interior
      wireframe}; the wall slabs themselves are not drawn as bodies.
      Creating a box while a window is open appends a reminder to the
      reply: \code{(scene window open: SCENE REFRESH shows the box)}.
      \kw{reset} re-syncs an open window to the now-empty system
      \textbf{and clears its box wireframe and wall flags} --- no
      stale overlay survives the wipe.
\end{itemize}

\subsection{User definitions: \kw{def}, \kw{let}, \kw{funcs}, \kw{show}}
\label{sec:userdefs}

\begin{lstlisting}
DEF name(param [= default], ...) { body }     (define a function)
name(arg, ...)                                (call it)
LET name = <expr>                             (session variable)
FUNCS                                         (list signatures; alias: FUNCTIONS)
SHOW name                                     (print a function's source)
\end{lstlisting}

\textbf{\kw{def} is a line form, not a grammar production.} A line
that starts with \code{DEF } is recognized \emph{before} the ordinary
grammar (\S3): the notebook captures everything up to the closing
\code{\}} --- interactively it keeps prompting \code{...:=} for
continuation lines until the brace closes; scripts and JupyterLab
cells just keep reading --- and installs the function. The
\textbf{body} is a sequence of ordinary commands, separated by
newlines or \code{;}, in which the parameters act as variables.
\textbf{Every body line is syntax-checked at definition time} (a typo
fails the \kw{def} immediately, naming the body line), and
\textbf{defaults are ordinary expressions evaluated once, at
definition} (\kw{let} variables are visible in them).

\textbf{Calling} uses the ordinary call syntax \code{name(arg, ...)}
--- the same atom as \code{sqrt(2)}; a name that is not a builtin is
looked up among your functions. Missing trailing arguments take their
defaults (an argument with no default must be supplied). Each call
pushes a \textbf{call frame} binding the parameters (depth cap: 32
--- the language has no conditionals, so recursion could never
terminate anyway) and runs the body lines through the same
compile-and-execute pipeline as typed input. The call \textbf{returns
the last body line's value} --- if that line is a \kw{set}, the call
prints no \code{Out[n]}, exactly like \kw{set} itself. A failing body
line aborts the call and the error names the function \emph{and} the
line; lines that already ran keep their effects (each command's own
guarantees still hold per line --- e.g.\ a failing \kw{new} inside a
body is still transactional), so a partly-run call leaves no
half-built object, only completed commands.

\textbf{Editing is \kw{show} + re-\kw{def}.} \code{SHOW name} prints
the definition verbatim, exactly as you typed it; redefining the same
name replaces the old version and the reply appends
\code{(redefined)}. \kw{funcs} lists every signature with its
defaults.

\textbf{\code{LET name = expr}} binds a session variable (reply:
\code{name set}). It is visible in bare expressions, in function
bodies, in \kw{def} defaults --- and, holding a string, it can
\emph{name} things: \kw{new} \ldots{} \kw{as} and named paths resolve
a bare identifier through the parameter/\kw{let} bindings first
(\S\ref{sec:new}). Note \kw{get} takes only a \emph{path} (\S3), so a
variable is read back as a bare expression: \code{g0}, not
\code{GET g0}.

A complete session (genuine output):

\begin{lstlisting}
In[1]:= let g0 = 9.81
Out[1]= g0 set
In[2]:= def drop(name, m = 1, h = 10) {
  new sphere as name { mass = m, position = [0, h, 0] }
  set system.gravity = [0, -g0, 0]
}
Out[2]= function drop(3 parameter(s)) defined — 2 body line(s)
In[3]:= drop("ball")
In[4]:= get ball.position.y
Out[4]= 10
In[5]:= drop("pebble", 0.1, 2)
In[6]:= get pebble.mass
Out[6]= 0.1
In[7]:= funcs
Out[7]= drop(name, m = 1, h = 10) — 2 body line(s); SHOW drop prints it
In[8]:= show drop
Out[8]= def drop(name, m = 1, h = 10) {
  new sphere as name { mass = m, position = [0, h, 0] }
  set system.gravity = [0, -g0, 0]
}
In[9]:= def drop(name, m = 1, h = 100) { new sphere as name { mass = m, position = [0, h, 0] }; set system.gravity = [0, -g0, 0] }
Out[9]= function drop(3 parameter(s)) defined — 2 body line(s) (redefined)
In[10]:= drop("skydiver")
In[11]:= skydiver.y
Out[11]= 100
In[12]:= speed
Err[12]: unknown name `speed` (define it with LET, pass it as a function parameter, or use `speed.field` for a registered object)
\end{lstlisting}

\emph{What to notice.} The function's \code{name} parameter holds a
string, and \code{new sphere as name} inside the body names each
created object from it --- \code{ball}, \code{pebble},
\code{skydiver} --- so one definition mints a family of named
objects. \code{drop("ball")} used both defaults; cell 9 is the edit
loop (copy cell 8's \kw{show} output, change \code{h}, re-\kw{def}
--- note \code{;} separating body lines works as well as newlines);
calls 3, 5 and 10 print no \code{Out} because the last body line is a
\kw{set}. Cell 12 is the bare-identifier error: execution-time
resolution means the message can list every way to bind a name.

Definition-time and call-time errors are specific: a reserved or
builtin name is refused (\code{DEF: `norm` is a builtin function and
cannot be redefined}), nested \kw{def} is refused, an empty body is
refused, a missing non-default argument fails as
\code{name(): missing argument `p` (it has no default)}, too many
arguments fail naming the signature, and the depth cap fails as
\code{function call depth limit (32) exceeded}.

\subsection{\kw{qm} --- one-dimensional quantum mechanics}\label{sec:qm}

The \kw{qm} family solves the two problems a 1-D quantum solver is for:
\textbf{bound states} in an arbitrary potential, and \textbf{time
evolution} of a wavepacket. A session carries one quantum problem, built
up command by command.

\begin{center}
\begin{tabular}{lp{7.4cm}}
\textbf{command} & \textbf{meaning}\\
\hline
\code{QM} & report what is currently set up\\
\code{QM GRID <x\_min> <x\_max> <n>} & the domain and its $n$ interior points\\
\code{QM POTENTIAL ZERO} & free particle\\
\code{QM POTENTIAL BARRIER <v0>, <x1>, <x2>} & \code{v0} on $[x_1,x_2]$, zero elsewhere\\
\code{QM POTENTIAL WELL <depth>, <x1>, <x2>} & $-$\code{depth} on $[x_1,x_2]$\\
\code{QM POTENTIAL <function>} & sample a \kw{def}ined $V(x)$ onto the grid\\
\code{QM MASS <m>} / \code{QM HBAR <h>} & both default to 1\\
\code{QM METHOD CAYLEY} & Crank--Nicolson on the Dirichlet grid --- the default\\
\code{QM METHOD NASH [LIE|STRANG]} & the Bessel-stencil split-operator scheme, \textbf{periodic}\\
\code{QM STATES <k>} & the $k$ lowest bound-state energies\\
\code{QM STATE <n>} & load bound state $n$ as $\psi$\\
\code{QM PACKET <x0> <sigma> <k0>} & a normalised Gaussian wavepacket\\
\code{QM STEP <dt>} / \code{QM RUN <t> [STEPS <n>]} & propagate with the current \code{QM METHOD}\\
\code{QM TRANSMISSION <e>} & $T(E)$ and $R(E)$ at one energy, by transfer matrix\\
\code{QM SCAN <e1> <e2> <n>} & scan $T(E)$, report resonances, push the list\\
\code{QM NORM} / \code{QM ENERGY} / \code{QM POSITION} / \code{QM MOMENTUM} & observables\\
\code{QM PROB <a> <b>} & probability of being found in $[a,b]$\\
\code{QM DENSITY} & $|\psi|^2$ as a list\\
\code{QM DRIVE <shape> <modulation>} & time-dependent $V(x,t) \mathrel{+}= f(t)g(x)$\\
\code{QM DRIVE OFF} & back to a static potential\\
\code{QM ABSORB <width> <strength> [<power>]} & absorbing edges; \code{power} defaults to 2\\
\code{QM ABSORB OFF} & back to reflecting walls\\
\code{QM ANIMATE "<file>" <t> [FRAMES <n>]} & write a self-contained HTML animation\\
\code{QM RESET} & forget the quantum problem\\
\end{tabular}
\end{center}

Five things are worth knowing before you start, because each will
otherwise cost you an afternoon.

\textbf{\code{QM METHOD} changes the boundary condition, not just the
algorithm.} \code{CAYLEY} is Crank--Nicolson on the Dirichlet grid: the
walls reflect. \code{NASH} is the Bessel-stencil split-operator scheme
ported from the original C++, and it is \textbf{periodic} --- a packet
leaving the right edge re-enters at the left. The grid points and the
potential samples are identical either way, so switching moves nothing;
only the two ends change meaning. The status line always states which
is in force.

Because bound states are computed with Dirichlet walls, \code{QM STATES}
and \code{QM STATE} are \textbf{refused} while the method is
\code{NASH} rather than quietly handing you eigenstates of a different
problem. \code{QM ABSORB} and \code{QM DRIVE} are likewise refused
under \code{NASH}: the propagator takes a real, static potential, and an
absorber is a complex one. Each refusal names the way out.

\code{NASH} alone is first order in $dt$ --- that is what the original
does, so it is the default. \code{NASH STRANG} is second order for
essentially the same cost, and is the better choice unless you are
reproducing SolveIt output.

\textbf{The walls are infinite and they reflect} --- under
\code{CAYLEY}, which is the default. \code{QM GRID} pins
$\psi$ to zero just outside the domain, which is exactly right for bound
states and a trap for scattering: a packet that reaches a wall bounces
back and corrupts your transmission number. \code{QM PACKET} and
\code{QM RUN} therefore \emph{watch for it} and print a warning when
probability accumulates within 5\,\% of either wall. Take the warning
seriously --- widen the domain.

\textbf{The potential is sampled when the command runs.}
\code{QM POTENTIAL v} evaluates $v(x)$ at each grid point once and
stores the numbers. Editing \code{v} afterwards changes nothing until
you re-issue the command. The status line says so.

\textbf{Separate negative arguments with commas.}
\code{QM POTENTIAL WELL 5 -2 2} reads \code{5 - 2} as subtraction and
then finds only two arguments where three were wanted. Write
\code{5, -2, 2} --- or \code{5 (-2) 2}. Space separation is fine when
everything is positive, which is the usual case.

\textbf{Why \code{BARRIER} and \code{WELL} are built in.} This language
has no comparison operators, so a piecewise potential cannot be written
as a \kw{def} at all --- and a square barrier is \emph{the} canonical
1-D problem. They are a deliberate workaround for a language limitation,
not a preference for built-ins. Any smooth potential should be a
\kw{def}.

\subsubsection*{Time-dependent potentials}

\code{QM DRIVE} makes the potential depend on time, as
$V(x,t) = V_0(x) + f(t)\,g(x)$, built from two \kw{def}ined functions.

\textbf{Why it is factorised rather than a general $V(x,t)$.} A general
form would need re-sampling at every grid point on every step --- for a
user-supplied function, thousands of VM calls per step, costing more
than the solve they feed. Factorising costs one evaluation per step and
covers the physically important cases exactly: a dipole drive $f(t)x$, a
shaken trap, a pulse envelope, an adiabatic ramp. A drive whose spatial
\emph{profile} changes shape with time is not expressible this way, and
that is a real limit rather than an oversight.

The modulation is sampled at the \textbf{midpoint} of each step, keeping
the scheme second order; sampling at the start would quietly halve it.

Two consequences:

\begin{itemize}
\item \textbf{Energy is no longer conserved.} A driven system exchanges
      energy with whatever drives it --- physics, not error. \code{QM RUN}
      says so, and the $\langle E\rangle$ it reports is that of the
      \emph{static} potential.
\item \textbf{Propagation is still unitary.} $H(t)$ is Hermitian at every
      instant, so each step remains an exact Cayley transform.
\end{itemize}

For a quadratic potential with a linear drive Ehrenfest's theorem is
\textbf{exact}: $\langle x\rangle$ obeys the classical equation of motion
with no approximation. Driving the oscillator ground state with
$f(t)=F_0\cos\omega t$ and $g(x)=x$ gives
\[
  x(t) = -\frac{F_0}{1-\omega^2}\bigl[\cos\omega t - \cos t\bigr],
\]
and the simulation reproduces it --- at $F_0=0.3$, $\omega=0.7$, $t=5$
the notebook gives $\langle x\rangle = 0.716918$ against an analytic
$0.717840$, with the norm at $1.000000000000146$.

\subsubsection*{Absorbing edges}

The reflecting walls are the main practical limit on a scattering run.
\code{QM ABSORB} replaces them with a \textbf{complex absorbing
potential}: the Hamiltonian becomes $H - iW(x)$ with $W\ge0$ ramping up
smoothly over \code{width} at each edge, so probability arriving there is
\emph{drained} rather than bounced.

Three consequences, all of them reported rather than left to be
discovered:

\begin{itemize}
\item \textbf{Propagation is no longer unitary.} The norm decays. That is
      the absorber working, so the drift figure printed by \code{QM RUN}
      stops being a health check while it is on.
\item \textbf{\code{QM STATES} is refused.} $H-iW$ is not Hermitian, and a
      symmetric eigensolver fed a non-Hermitian matrix returns confident
      nonsense.
\item \textbf{The tuning is real.} Too weak and the packet reaches the
      wall and reflects off \emph{that}; too strong and it reflects off
      the absorber's own leading edge, because a steep change in the
      potential is a mirror whether it is real or imaginary. The useful
      window spans more than an order of magnitude in strength.
      \code{cargo run -p quantum -{}-release -{}-example absorber\_tuning}
      measures the whole surface; for $k_0\approx3$ the optimum is near
      width 18, strength 3, giving about $10^{-9}$ reflection against a
      plain wall's $1.0$.
\end{itemize}

A worked comparison: the barrier problem gives $T = 0.327563$ on a
200-wide reflecting domain, and $T = 0.327690$ on a \textbf{90-wide}
domain with \code{QM ABSORB 15 3} --- agreement to 4 parts in 10\,000 for
less than half the grid.

Bound states are found by diagonalising a dense $n\times n$ matrix with
the cyclic Jacobi method, which is $O(n^3)$: comfortable to a few
hundred points, slow past about a thousand. Propagation is $O(n)$ per
step and has no such limit, so a scattering run can afford a much finer
grid than a bound-state calculation.

\subsection{\kw{qm2} --- two dimensions, by ADI}\label{sec:qm2}

A \emph{separate} family from \kw{qm}, not a mode on it: a 2-D problem
differs in almost every argument list, and a hidden mode that silently
reinterprets your commands is worse than a second word.

\begin{center}
\begin{tabular}{lp{7.2cm}}
\textbf{command} & \textbf{meaning}\\
\hline
\code{QM2} & report the current 2-D setup\\
\code{QM2 GRID <x0> <x1> <nx>, <y0> <y1> <ny>} & domain and resolution per axis\\
\code{QM2 POTENTIAL ZERO | <function>} & a \kw{def}ined $V(x,y)$\\
\code{QM2 PACKET <x0> <y0>, <sx> <sy>, <kx> <ky>} & 2-D Gaussian packet\\
\code{QM2 STATES <k>} & the $k$ lowest bound-state energies\\
\code{QM2 STATE <n>} & load bound state $n$ as $\psi$\\
\code{QM2 STEP <dt>} / \code{QM2 RUN <t> [STEPS <n>]} & ADI propagation\\
\code{QM2 NORM | QM2 ENERGY | QM2 CENTROID} & observables (\code{CENTROID} also spells \code{POSITION})\\
\code{QM2 PROB <xa> <xb>, <ya> <yb>} & probability in a rectangle\\
\code{QM2 DRIVE <shape> <modulation> | OFF} & time-dependent $V(x,y,t)$\\
\code{QM2 ABSORB <w> <s> [<p>] | OFF} & absorbing edges, all four walls\\
\code{QM2 ANIMATE "<file>" <t> [FRAMES <n>]} & heat-map animation\\
\code{QM2 RESET} & forget the 2-D problem\\
\end{tabular}
\end{center}

\subsubsection*{Why ADI, and why not the textbook ADI}

In 1-D the Crank--Nicolson operator is tridiagonal, so one solve per step
gives an exactly unitary propagator. \textbf{In 2-D it is not}: the
five-point Laplacian couples each point to neighbours in both directions
and the matrix is block-tridiagonal with bandwidth $n_x$, costing
$O(n_x^3 n_y)$ per step to solve directly.

ADI splits the step by direction, so each half-step is a \emph{set of
independent tridiagonal solves} --- $n_y$ along $x$, then $n_x$ along
$y$ --- and the cost falls to $O(n_x n_y)$.

The textbook scheme is \textbf{Peaceman--Rachford}, implicit in $x$
against an explicit $y$ and then swapped. It is second-order and
unconditionally stable, but \textbf{not unitary}: the two half-steps use
different operators, so the norm drifts at $O(dt^2)$ per step.

This implementation applies a \textbf{Cayley transform in each direction
separately}, composed by Strang splitting:
\[
  \psi(t+dt) = U_x(dt/2)\,U_y(dt)\,U_x(dt/2)\,\psi(t).
\]
Each $A_d = T_d + V/2$ is Hermitian, so each $U_d$ is exactly unitary and
so is their product. \textbf{The norm is conserved to machine precision
for any $dt$}, while the splitting error stays $O(dt^2)$ in the
\emph{dynamics}. Norm conservation therefore remains a sharp check on the
linear algebra, independent of the accuracy question.

One subtlety: with $A_x=T_x+V/2$ and $A_y=T_y+V/2$ the commutator
$[A_x,A_y]$ is non-zero \emph{even for a separable potential}, so a
product eigenstate is not perfectly stationary. The drift is second order
in $dt$ and the tests assert that scaling rather than a magnitude.

\subsubsection*{Bound states, and why they need a different eigensolver}

On an $n_x\times n_y$ grid the Hamiltonian is $(n_x n_y)^2$: a modest
$200\times200$ grid gives a $40\,000\times40\,000$ matrix, about 12\,GB
dense before any arithmetic. The Jacobi solver used in 1-D does not
apply.

\code{QM2 STATES} uses \textbf{Lanczos}, which never forms the matrix ---
it needs only $Hv$, the five-point stencil at $O(n_x n_y)$. It builds a
Krylov basis, projects $H$ onto it as a small tridiagonal, and those
eigenvalues converge fastest to the \emph{extremes} of the spectrum,
which is where bound states are.

Two separate things must go right for degeneracies, and they are easy to
confuse:

\begin{itemize}
\item \textbf{Ghosts.} Lanczos vectors lose orthogonality in floating
      point as soon as a value converges, and the method rediscovers
      eigenvalues it already has. Those spurious copies look exactly like
      degeneracy. Cured by \textbf{full reorthogonalisation}.
\item \textbf{Genuine multiplicity.} Full reorthogonalisation does
      \emph{not} help. A Krylov space built from one starting vector
      contains exactly \textbf{one} direction from each degenerate
      eigenspace, so plain Lanczos finds each distinct eigenvalue once
      however long it runs; no tolerance fixes it. Cured by
      \textbf{deflation}: each converged state is shifted to the top of
      the spectrum and the solver re-run \emph{from a different starting
      vector} --- reusing the same start would re-derive the direction
      just removed.
\end{itemize}

The 2-D isotropic oscillator on a $70\times70$ grid:

\begin{lstlisting}
In[4]:= qm2 states 6
Out[4]= 6 lowest bound state(s) -- Lanczos, 620 iterations:
  E[0] = 0.9975639778   residual 3.86e-8
  E[1] = 1.9926799032   residual 1.92e-8
  E[2] = 1.9926799032   residual 2.89e-8
  E[3] = 2.9828813394   residual 3.03e-8
  E[4] = 2.9828813394   residual 5.15e-8
  E[5] = 2.9877958286   residual 5.21e-8
\end{lstlisting}

Against the exact $1,2,2,3,3,3$. Note what the grid does to the third
level: $E[3]$ and $E[4]$ agree to \textbf{ten digits}, because the square
grid keeps the $x\leftrightarrow y$ symmetry relating the $(2,0)$ and
$(0,2)$ states --- while $E[5]$, the $(1,1)$ state, sits $0.005$ away
because the continuum \emph{rotational} symmetry that would complete the
degeneracy is not a symmetry of the grid. The splitting is
discretisation, not solver error.

\textbf{Residuals are printed as part of the answer}: an iterative
eigensolver has no exact stopping point, and a caller who cannot see how
well a state converged cannot know whether to trust it.

Propagation has no comparable limit.

\subsection{\kw{qm3} --- three dimensions}\label{sec:qm3}

A third family, for the same reason \kw{qm2} is separate from \kw{qm}.

\begin{center}
\begin{tabular}{lp{6.6cm}}
\textbf{command} & \textbf{meaning}\\
\hline
\code{QM3 GRID <x0> <x1> <nx>, ..., <z0> <z1> <nz>} & domain and resolution per axis\\
\code{QM3 POTENTIAL ZERO | <function>} & a \kw{def}ined $V(x,y,z)$\\
\code{QM3 PACKET <x0> <y0> <z0>, ...} & 3-D Gaussian packet\\
\code{QM3 STATES <k> | QM3 STATE <n>} & bound states\\
\code{QM3 STEP <dt> | QM3 RUN <t> [STEPS <n>]} & ADI propagation\\
\code{QM3 NORM | QM3 ENERGY | QM3 CENTROID} & observables (\code{CENTROID} also spells \code{POSITION})\\
\code{QM3 PROB <xa> <xb>, ..., <za> <zb>} & probability in a box\\
\code{QM3 DRIVE <shape> <modulation> | OFF} & time-dependent $V(x,y,z,t)$\\
\code{QM3 ABSORB <w> <s> [<p>] | OFF} & absorbing faces, all six sides\\
\code{QM3 ANIMATE "<file>" <t> [FRAMES <n>]} & three marginal densities, animated\\
\code{QM3 ISO "<file>" <t> [FRAMES <n>] [LEVEL <f>]} & rotatable isosurface\\
\code{QM3 RESET} & forget the 3-D problem\\
\end{tabular}
\end{center}

The scheme is the 2-D one with a third direction, Strang-composed so
every factor stays an exact Cayley transform:
\[
  \psi(t+dt) = U_x(dt/2)\,U_y(dt/2)\,U_z(dt)\,U_y(dt/2)\,U_x(dt/2)\,\psi(t),
\]
with $A_d = T_d + V/3$. Five directional sweeps per step instead of
three, each a set of independent tridiagonal solves, so a step is still
$O(n_x n_y n_z)$ and \textbf{the norm is conserved to machine precision
for any $dt$}.

\subsubsection*{What changes in three dimensions is memory}

A $100^3$ grid is a million points. One complex wavefunction is 16\,MB,
which is fine --- but the 2-D propagator precomputes full-size band
arrays, three per direction, and doing that here would cost about
140\,MB before any work began. \kw{qm3} builds its bands \textbf{per
line, on the fly}, so the only extra storage is
$O(\max(n_x,n_y,n_z))$.

\subsubsection*{The eigensolver has a ceiling, and says so}

Propagation is comfortable past $64^3$. \code{QM3 STATES} is not:
Lanczos reorthogonalises fully and stores its whole Krylov basis,
costing $O(m^2 n)$ in time and $O(mn)$ in memory. The practical ceiling
is about $\mathbf{40^3}$, and beyond it the command \textbf{refuses up
front} rather than exhausting memory. Propagation on the same grid still
works.

\subsubsection*{Looking at a volume}

A volume cannot be drawn on a flat canvas without an isosurface mesh or
a ray-caster, either of which would mean shipping a WebGL pipeline
inside a file that must work from \code{file://}. \code{QM3 ANIMATE}
shows the three \textbf{marginal densities} instead:
\[
  P(x,y)=\!\int\!|\psi|^2 dz,\quad
  P(x,z)=\!\int\!|\psi|^2 dy,\quad
  P(y,z)=\!\int\!|\psi|^2 dx .
\]
Each is a genuine observable --- the probability of finding the particle
at those two coordinates whatever the third --- not a rendering
convention, and each integrates to the total norm, which the page prints
under every panel.

What they cannot show is correlation between axes: a state concentrated
on a diagonal shell and one spread over a box can share all three
marginals. That is a real limit of the representation, and the reason a
true isosurface view is listed as unfinished rather than unnecessary.

Verified in a browser on a driven $34^3$ run: all three marginals
integrate to $1.000000$ in every sampled frame, the $P(x,y)$ peak travels
along $x$ while its $y$ coordinate does not move, and $P(y,z)$ does not
move at all --- exactly right for a drive along $x$.

\subsubsection*{Isosurfaces}

Marginals cannot show correlation between axes; an isosurface can.
\code{QM3 ISO} extracts the surface where $|\psi|^2$ equals a chosen
fraction of its peak.

\textbf{Marching tetrahedra, not marching cubes.} Marching cubes needs a
256-entry table mapping corner-sign patterns to triangle lists --- and
that table \emph{is} the algorithm, so one wrong entry gives a surface
with a hole that looks fine from most angles. Marching tetrahedra needs
no table: six tetrahedra per cube, $2^4=16$ sign patterns, three cases up
to symmetry. The classic ambiguity where two cubes triangulate a shared
face incompatibly cannot arise, because a tetrahedron's faces are
triangles.

The extractor is a library function with quantitative tests: for a sphere
and an ellipsoid it checks enclosed volume by the divergence theorem,
surface area, convergence under refinement, and that the mesh is
\textbf{watertight and consistently oriented} --- every directed edge
exactly once, its reverse exactly once.

Two defects that testing found rather than inspection: samples lying
\emph{exactly} on the isolevel produced coincident vertices and pinholes
(fixed by nudging the level, not the samples); and orienting a split
quad's two triangles independently let one flip and open the shared
diagonal (they now share one winding decision).

Rendering is a \textbf{software rasteriser on a 2-D canvas}, not WebGL:
WebGL can fail silently without a GPU and then the page shows nothing,
whereas a painter's-algorithm rasteriser works everywhere and can be
verified by reading pixels back.

\subsubsection*{The 3-D oscillator}

\begin{lstlisting}
In[4]:= qm3 states 4
Out[4]= 4 lowest bound state(s) -- Lanczos, 245 iterations:
  E[0] = 1.4738115042   residual 5.71e-8
  E[1] = 2.4382167269   residual 2.51e-8
  E[2] = 2.4382167269   residual 5.03e-8
  E[3] = 2.4382167269   residual 5.69e-8
\end{lstlisting}

Against the exact $E=n_x+n_y+n_z+3/2$: $3/2$, then $5/2$ \textbf{three
times}. The three excited values agree to ten digits --- that degeneracy
is what deflation exists to resolve --- while both levels sit about 0.03
below the continuum, which is second-order discretisation error at
$h=0.52$.

\section{The notebook (cells, editing, magics)}

\subsection{Cells}
Start \code{posim} with no arguments. You get numbered prompts, exactly
like Jupyter or Mathematica:
\begin{lstlisting}
In[1]:= new sphere { mass = 2, radius = 0.5 }
Out[1]= obj0
In[2]:= get obj0.inertia_tensor
Out[2]= [[0.2, 0, 0], [0, 0.2, 0], [0, 0, 0.2]]
\end{lstlisting}
Pressing \textbf{Enter executes the cell} --- in a plain terminal, Enter
\emph{is} shift-enter. Commands with no interesting result (like
\kw{set}) print no \code{Out[n]}. Errors print \code{Err[n]:} and the
session continues.

\subsection{State is cumulative}
Like Jupyter, the notebook has one live simulator state that every cell
mutates in order. Re-running an old \kw{new} cell creates a \emph{new}
object; it does not replace the old one. To rebuild cleanly,
\code{\%reset} and replay.

\subsection{Magics}\label{sec:magics}
\begin{center}
\begin{tabular}{lp{9.8cm}}
\textbf{magic} & \textbf{effect}\\\hline
\code{\%history} & lists every cell with input and output; failed cells marked \code{!}\\
\code{\%edit n <text>} & replaces cell $n$'s input and executes it as the next cell\\
\code{\%rerun n} & executes cell $n$'s input again as a new cell\\
\code{\%save <file>} & writes all successful non-magic inputs to a replayable script\\
\code{\%load <file>} & replays a script file cell by cell\\
\code{\%reset} & clears the simulator (history kept)\\
\code{\%quit}, \code{\%exit} & leave\\
\end{tabular}
\end{center}

A plain terminal has no cursor-addressable cells (posim uses only the
Rust standard library --- no raw terminal mode), so backward/forward
movement is by these magics. \textbf{For true click-and-edit cells with
shift-enter, use the JupyterLab front end} (\S\ref{sec:jlab}).

\subsection{Scripts}
\code{posim --script file} executes a file of command lines (blank lines
and \code{\#} comments skipped), echoing \code{In[n]}/\code{Out[n]}, and
exits nonzero if any cell failed --- good for CI.

\code{posim --notebook file} executes the same file format but then
\textbf{stays in the interactive notebook} instead of exiting: the
loaded cells keep their \code{In[n]} numbers, the next command
continues the numbering, and a scene window the file opened
(\code{SCENE CREATE}) stays alive --- the launcher for the
\emph{dynamic notebooks} in \code{dynamic\_notebooks/}, files that
build a simulation and open the GUI ready for Start. A failing cell
aborts the load with a nonzero exit instead of entering the notebook.

\section{JupyterLab in one paragraph}\label{sec:jlab}

The \code{jupyter/} directory ships a small \emph{wrapper kernel}:
JupyterLab starts a Python shim, the shim starts \code{posim --machine}
(a JSON request/reply protocol over stdin/stdout), and each notebook
cell you shift-enter is forwarded line-by-line as
\code{\{"op":"exec","code":...\}}. Setup and tests:
\code{jupyter/README.md}. Everything in \S\S2--5 applies verbatim inside
JupyterLab cells; multi-line cells run one line at a time, stopping at
the first error.

\section{The stack machine (what runs your line)}

The parser emits \emph{postfix instructions}; e.g.\
\code{set obj0.mass = 2 + 3 * 4} compiles to
\begin{lstlisting}
Push 2
Push 3
Push 4
Mul               <- pops 3,4  pushes 12
Add               <- pops 2,12 pushes 14
Store obj0.mass   <- pops 14, calls set_mass(14)
\end{lstlisting}
The machine's whole memory is a stack of typed values. Instructions:
\code{Push}, \code{Load}, \code{Store}, \code{LoadIdent} (a bare
identifier --- parameter or \kw{let} variable, resolved at execution,
\S\ref{sec:userdefs}), \code{StoreGlobal} (\kw{let}),
\code{Add Sub Mul Div Neg},
\code{PackList} (literals), \code{Call} (builtin or user function ---
a user call runs each body line through this same compile-and-execute
pipeline inside a call frame),
\code{NewObject}/\code{InitField}/\code{FinishNew} (which also
registers the \kw{as} name), \code{Delete},
\code{ListObjects}, \code{ListFns}/\code{ShowFn}
(\kw{funcs}/\kw{show}), \code{Step}, \code{Run}, \code{SetMethod},
\code{Energy}, \code{CenterOfMass}, \code{TotalMomentum},
\code{TotalAngularMomentum}, \code{Laplace}, \code{Reset}, \code{Help}.
\code{Load}/\code{Store} go \emph{only} through the
\code{physical\_object} get/set API --- the machine cannot corrupt
simulator state, and every coupled invariant (mass$\leftrightarrow$inverse,
inertia$\leftrightarrow$inverse, unit quaternions) is enforced on every
write. When the program ends, the top of the stack becomes \code{Out[n]}.

\section{Eighteen worked examples}

All transcripts are genuine program output.

\subsection{Example 1 --- an eccentric Kepler orbit and the Runge--Lenz compass}

The Laplace--Runge--Lenz vector $\mathbf A$ points along a Kepler
orbit's major axis and is conserved \emph{only} for a perfect $1/r^2$
force --- the most delicate conservation test available. We build a
``sun'' so heavy the reduced problem is exact and integrate two orbits
with a 4th-order \emph{symplectic} method.

\begin{lstlisting}
In[1]:= new point { mass = 1e9, position = [0, 0, 0] }
Out[1]= obj0
In[2]:= new point { mass = 1, position = [0.4, 0, 0], velocity = [0, 2, 0] }
Out[2]= obj1
In[3]:= set system.g_constant = 1e-9
In[4]:= set system.softening = 0
In[5]:= laplace 1
Out[5]= [0.5999999974000001, 0, 0]
In[6]:= method sprk mclachlan_4_4 0.001
Out[6]= method = ARKODE SPRK ARKODE_SPRK_MCLACHLAN_4_4, fixed dt = 0.001
In[7]:= run 12.6 steps 2
Out[7]= t = 12.6 (12600 solver steps, 2 snapshots, |dE/E| = 9.237e-14)
In[8]:= laplace 1
Out[8]= [0.5999999974140128, -0.00000000034051644837163053, 0]
In[9]:= get obj1.position
Out[9]= [0.3964802853115723, 0.06706198328283366, 0]
\end{lstlisting}

\emph{What to notice.} \code{g\_constant = 1e-9} makes
$G M_{\text{total}} = 1$ so the orbit has period $2\pi$. \kw{laplace}~1
divided by $mk$ is the \emph{eccentricity vector}: $e = 0.6$ is read
directly off \code{Out[5]}. After $\sim$2 orbits the energy is conserved
to $9\times10^{-14}$ and $\mathbf A$ has rotated by only
$3\times10^{-10}$ --- no spurious perihelion precession. Softening was
zeroed because any $\varepsilon\ne 0$ slightly breaks the $1/r^2$ law
and \emph{would} precess the axis.

\subsection{Example 2 --- a thrown ball vs.\ the textbook formula}

A projectile has the closed form $x(t)=v_x t$,
$y(t)=y_0+v_y t-\tfrac12 g t^2$. We fly a baseball for 6.12\,s and
subtract the formula \emph{inside the notebook} using bare expressions.

\begin{lstlisting}
In[1]:= new point { mass = 0.145, position = [0, 1, 0], velocity = [30, 30, 0] }
Out[1]= obj0
In[2]:= set system.gravity = [0, -9.81, 0]
In[3]:= run 6.12 steps 3
Out[3]= t = 6.12 (13 solver steps, 3 snapshots, |dE/E| = 4.740e-15)
In[4]:= get obj0.position.x
Out[4]= 183.5999999999999
In[5]:= 30 * 6.12
Out[5]= 183.6
In[6]:= obj0.position.x - 30 * 6.12
Out[6]= -0.00000000000008526512829121202
In[7]:= obj0.position.y - (1 + 30 * 6.12 - 9.81 / 2 * 6.12 * 6.12)
Out[7]= -0.0000000000004263256414560601
\end{lstlisting}

\emph{What to notice.} Cells 6--7 are \emph{bare expressions} mixing
paths and arithmetic (\kw{get} itself would refuse the arithmetic). The
trajectory matches the parabola to $\sim\!10^{-13}$ --- and the adaptive
Adams integrator needed only \textbf{13 internal steps}, because for a
polynomial solution its error estimator lets it take huge strides.

\subsection{Example 3 --- cyclotron motion: expressions as arguments}

A charge in a uniform magnetic field circles with period
$T = 2\pi m/(|q|B)$. We \emph{compute $T$ inside the \kw{run} command}.

\begin{lstlisting}
In[1]:= new sphere { mass = 2, radius = 0.1, charge = -1.5, velocity = [3, 0, 0] }
Out[1]= obj0
In[2]:= set system.b_field = [0, 0, 4]
In[3]:= method bdf
Out[3]= method = CVODE BDF
In[4]:= 2 * pi * 2 / (1.5 * 4)
Out[4]= 2.0943951023931953
In[5]:= run 2 * pi * 2 / (1.5 * 4) steps 8
Out[5]= t = 2.0943951023931953 (318 solver steps, 8 snapshots, |dE/E| = 1.444e-8)
In[6]:= get obj0.position
Out[6]= [0.00000000043106769672882073, 0.000000007220835657713453, 0]
In[7]:= norm(obj0.velocity)
Out[7]= 2.9999999783374913
\end{lstlisting}

\emph{What to notice.} After one analytic period the particle is back at
the origin to $7\times10^{-9}$, and \code{norm()} shows the speed
unchanged --- the magnetic force does no work. The Lorentz force
$q\,v\times B$ is velocity-dependent, so the symplectic path is
\emph{not allowed} here; \kw{bdf} is the right tool for fast gyration.

\subsection{Example 4 --- the tennis-racket theorem (Dzhanibekov effect)}

A rigid body spun about its \emph{intermediate} principal axis is
unstable: a tiny wobble grows into a dramatic flip, yet energy and
angular momentum stay conserved. Half-extents \code{[0.5, 1, 2]} give
moments $[5, 4.25, 1.25]$; spinning about $y$ (the middle value)
triggers it.

\begin{lstlisting}
In[1]:= new cuboid { mass = 3, half_extents = [0.5, 1, 2], angular_velocity = [0.01, 3, 0.01] }
Out[1]= obj0
In[2]:= get obj0.inertia_tensor
Out[2]= [[5, 0, 0], [0, 4.25, 0], [0, 0, 1.25]]
In[3]:= get obj0.angular_velocity
Out[3]= [0.010000000000000002, 3, 0.010000000000000002]
In[4]:= run 40 steps 4
Out[4]= t = 40 (2361 solver steps, 4 snapshots, |dE/E| = 5.605e-9)
In[5]:= get obj0.angular_velocity
Out[5]= [1.5428438559953521, 2.9947703802651175, -0.7871804456905063]
In[6]:= run 40 steps 4
Out[6]= t = 80 (2334 solver steps, 4 snapshots, |dE/E| = 5.332e-9)
In[7]:= get obj0.angular_velocity
Out[7]= [0.020666590902422656, 2.9999548562146483, 0.013346831306280834]
In[8]:= get obj0.angular_momentum
Out[8]= [0.05, 12.75, 0.0125]
\end{lstlisting}

\emph{What to notice.} These two samples are snapshots of a
\emph{recurring} motion, not the start and end of one flip. The
turn-over is a \emph{body-frame} event: the angular momentum is fixed in
space, so the world-frame $\omega_y$ printed above stays near $+3$ for
the whole run and never reverses at all. Extract
$\omega_{\text{body}} = \bar{q}\,\omega_{\text{world}}\,q$ and its
$y$-component reverses at roughly $t = 4, 18, 28, 40, 50, 62, 74$ ---
about every 12--13 units. By $t=80$ the body has turned over some
\emph{seven} times; $t=40$ merely lands mid-reversal and $t=80$ between
reversals. Throughout, the world-frame angular momentum
(\code{Out[8]}) is exactly the initial
$I\omega = [0.05, 12.75, 0.0125]$ --- the solver integrates the full
quaternion + angular-momentum state with energy error
$\sim 5\times10^{-9}$. (With these moments a symmetric wobble
$(d, 3, d)$ sits exactly on the separatrix $G^2 = BT$, since
$G^2 - BT = d^2[A(A-B) + C(C-B)] = d^2[5(0.75) + 1.25(-3)] = 0$ for
every $d$; the repeated reversal is robust, but the exact times shift
with the tolerances. \code{dynamic\_notebooks/tumbling\_body.posim}
shows the body-frame component directly.)

\subsection{Example 5 --- three bodies, then surgery with \kw{del}}

\begin{lstlisting}
In[1]:= new point { mass = 1,   position = [1, 0, 0],  velocity = [0, 1, 0] }
Out[1]= obj0
In[2]:= new point { mass = 4,   position = [-1, 0, 0], velocity = [0, -0.25, 0] }
Out[2]= obj1
In[3]:= new point { mass = 256, position = [0, 8, 0],  velocity = [0, 0, 0] }
Out[3]= obj2
In[4]:= set system.g_constant = 0.001
In[5]:= momentum
Out[5]= [0, 0, 0]
In[6]:= com
Out[6]= [-0.011494252873563218, 7.846743295019157, 0]
In[7]:= run 3 steps 3
Out[7]= t = 3 (84 solver steps, 3 snapshots, |dE/E| = 3.141e-10)
In[8]:= list
Out[8]= obj0: point, mass=1, charge=0, pos=[0.9936356911262184, 3.022315903943394, 0]
obj1: point, mass=4, charge=0, pos=[-0.9972663866755148, -0.7330950497234504, 0]
obj2: point, mass=256, charge=0, pos=[-0.000017852126656859762, 7.999648688652149, 0]
In[9]:= del 2
Out[9]= deleted obj2; 2 object(s) remain (indices renumbered)
In[10]:= list
Out[10]= obj0: point, mass=1, charge=0, pos=[0.9936356911262184, 3.022315903943394, 0]
obj1: point, mass=4, charge=0, pos=[-0.9972663866755148, -0.7330950497234504, 0]
In[11]:= momentum
Out[11]= [0.0021430470372574067, 0.061528401914275554, 0]
\end{lstlisting}

\emph{What to notice.} Cells 1--3 chose velocities so the initial total
momentum is exactly zero ($1\cdot1 + 4\cdot(-0.25) = 0$), and it stays
zero through the run. Deleting the heavy body removes its momentum: the
remainder (\code{Out[11]}) is what the surviving pair carries ---
conservation applies to the system you keep. \kw{del} renumbers later
objects.

\subsection{Example 6 --- magnetic torque (and why \kw{energy} grows)}

$\code{magnetic\_moment\_tensor}$ $M$ couples the world $B$ field to
torque $\tau = (RMR^{\mathsf T})B$. Unlike everything else, this
coupling is \emph{not derived from a potential}, so total energy is
\emph{expected} to change.

\begin{lstlisting}
In[1]:= new sphere { mass = 1, radius = 0.5, magnetic_moment_tensor = [[0.2, 0, 0], [0, 0.2, 0], [0, 0, 0.2]] }
Out[1]= obj0
In[2]:= set system.b_field = [0, 0.5, 0]
In[3]:= get obj0.angular_momentum
Out[3]= [0, 0, 0]
In[4]:= energy
Out[4]= 0
In[5]:= run 4 steps 4
Out[5]= t = 4 (214 solver steps, 4 snapshots, |dE/E| = 8.000e-1)
In[6]:= get obj0.angular_momentum
Out[6]= [0, 0.40000000000000085, 0]
In[7]:= get obj0.angular_velocity
Out[7]= [0, 4.000000000000009, 0]
In[8]:= energy
Out[8]= 0.8000000000000035
\end{lstlisting}

\emph{What to notice.} Constant torque $\tau = MB = 0.2\times0.5 = 0.1$
about $y$ gives $L(t) = 0.1t$: after 4\,s, $L = 0.4$ exactly (the ODE is
linear, so the solver nails it). Sphere inertia
$= 0.4\cdot1\cdot0.25 = 0.1$, so $\omega = L/I = 4$ and the rotational
energy $\tfrac12\omega L = 0.8$ matches \code{Out[8]}. The reported
\code{|dE/E| = 8e-1} is not an error --- it honestly reports that the
torque pumped energy in.

\subsection{Example 7 --- fixing a typo in an old cell with \code{\%edit}}

\begin{lstlisting}
In[1]:= new sphere { mass = 20, radius = 0.5 }
Out[1]= obj0
In[2]:= set obj0.velocity = [1, 0, 0]
In[3]:= get obj0.momentum
Out[3]= [20, 0, 0]
In[4]:= %edit 1 new sphere { mass = 2, radius = 0.5 }
In[4]:= new sphere { mass = 2, radius = 0.5 }
Out[4]= obj1
In[5]:= %history
 In[1]:= new sphere { mass = 2, radius = 0.5 }
  Out[1]= obj0
 In[2]:= set obj0.velocity = [1, 0, 0]
 In[3]:= get obj0.momentum
  Out[3]= [20, 0, 0]
 In[4]:= new sphere { mass = 2, radius = 0.5 }
  Out[4]= obj1
\end{lstlisting}

\emph{What to notice.} \code{\%edit 1 ...} rewrites cell 1's stored
input \textbf{and re-executes it as the next cell} --- like editing a
Jupyter cell and shift-entering again. Because state is cumulative, the
re-executed \kw{new} made a \emph{second} object; the fat \code{obj0} is
still there. For a clean rebuild use \code{\%reset} + \code{\%rerun} or
\code{\%load}; when a field tweak suffices, \code{set obj0.mass = 2}.
\code{\%history} shows the \emph{edited} input but the \emph{original}
outputs --- an audit trail of what actually ran.

\subsection{Example 8 --- checking the simulator against itself}

\begin{lstlisting}
In[1]:= new point { mass = 2, position = [1, 0, 0], velocity = [0, 3, 0] }
Out[1]= obj0
In[2]:= cross(obj0.position, obj0.momentum)
Out[2]= [0, 0, 6]
In[3]:= angmom
Out[3]= [0, 0, 6]
In[4]:= dot(obj0.position, obj0.velocity)
Out[4]= 0
In[5]:= norm(cross(obj0.position, obj0.momentum)) / (norm(obj0.position) * norm(obj0.momentum))
Out[5]= 1
\end{lstlisting}

\emph{What to notice.} Cell 2 computes $L = r\times p$ by hand and cell
3 confirms \kw{angmom} agrees. Cell 4 proves $r\perp v$; cell 5 computes
$|r\times p|/(|r||p|) = \sin\theta = 1$ --- the angle between $r$ and
$p$ is $90^\circ$ --- inside one nested expression.

\subsection{Example 9 --- hand-built tensors and a quaternion orientation}

\begin{lstlisting}
In[1]:= new cuboid { mass = 1, inertia_tensor = [[2, 0, 0], [0, 3, 0], [0, 0, 4]], orientation = [0.7071067811865476, 0, 0, 0.7071067811865476] }
Out[1]= obj0
In[2]:= get obj0.inverse_inertia_tensor
Out[2]= [[0.5, 0, 0], [0, 0.3333333333333333, 0], [0, 0, 0.25]]
In[3]:= set obj0.angular_velocity = [0, 2, 0]
In[4]:= get obj0.angular_momentum
Out[4]= [0.0000000000000004440892098500628, 4.000000000000002, 0]
In[5]:= get obj0.orientation
Out[5]= quat[w=0.7071067811865476, x=0, y=0, z=0.7071067811865476]
\end{lstlisting}

\emph{What to notice.} Because \code{inertia\_tensor} appears in the
initializer, the analytic inertia is \textbf{not} recomputed --- your
$\mathrm{diag}(2,3,4)$ is kept, with the inverse maintained
automatically. The quaternion $[w,x,y,z] = [\sqrt{1/2},0,0,\sqrt{1/2}]$
is a $90^\circ$ rotation about $z$, so setting the \emph{world} angular
velocity $[0,2,0]$ exercises the full $L = (R I R^{\mathsf T})\omega$
transformation. By hand: rotating $\mathrm{diag}(2,3,4)$ by $90^\circ$
about $z$ swaps the $x$ and $y$ moments, so
$R I R^{\mathsf T} = \mathrm{diag}(3,2,4)$, the world-$y$ moment is 2,
and $L = [0, 4, 0]$ --- exactly \code{Out[4]} (up to $4\times10^{-16}$
of floating-point dust).

\subsection{Example 10 --- reproducibility: \code{\%save}, \code{\%reset}, \code{\%load}}

\begin{lstlisting}
In[1]:= new sphere { mass = 2, radius = 0.5, velocity = [1, 0, 0] }
Out[1]= obj0
In[2]:= set system.gravity = [0, -9.81, 0]
In[3]:= step 0.5
Out[3]= t = 0.5 (advanced by 0.5, 12 solver steps)
In[4]:= %save session.posim
saved 3 cell(s) to session.posim
In[4]:= %reset
system reset
In[4]:= list
Out[4]= (no objects)
In[5]:= %load session.posim
In[5]:= new sphere { mass = 2, radius = 0.5, velocity = [1, 0, 0] }
Out[5]= obj0
In[6]:= set system.gravity = [0, -9.81, 0]
In[7]:= step 0.5
Out[7]= t = 0.5 (advanced by 0.5, 12 solver steps)
In[8]:= get obj0.velocity
Out[8]= [1, -4.905000000000001, 0]
\end{lstlisting}

\emph{What to notice.} \code{\%save} writes only the successful,
non-magic inputs --- a clean, replayable script. After \code{\%reset},
\code{\%load} replays it and the state is bit-identical
($v_y = -9.81\cdot0.5 = -4.905$). The same file runs headlessly via
\code{posim --script}.

Notice too that \textbf{magics do not consume cell numbers}:
\code{\%save} and \code{\%reset} print their message and leave the
counter alone, which is why the prompt shows \code{In[4]} three times
in a row before \kw{list} finally becomes cell 4. Only executed
commands become numbered cells --- so the \code{In[n]} numbering always
matches what \code{\%history} will replay.

The machine mode speaks the same language over JSON:

\begin{lstlisting}
$ printf '%s\n' '{"op":"exec","code":"new point { mass = 1, velocity = [0, 1, 0] }"}' \
                '{"op":"get","path":"obj0.momentum"}' \
                '{"op":"set","path":"obj0.mass","value":5}' \
                '{"op":"get","path":"obj0.momentum"}' | posim --machine
{"display":"obj0","ok":true,"result":"obj0"}
{"display":"[0, 1, 0]","ok":true,"result":[0.0,1.0,0.0]}
{"display":"","ok":true,"result":null}
{"display":"[0, 1, 0]","ok":true,"result":[0.0,1.0,0.0]}
\end{lstlisting}

Changing the mass did \textbf{not} change the momentum --- momentum is
the canonical stored state; the \emph{velocity} is what changed. That
final subtlety is the union design showing through the wire protocol.

\subsection{Example 11 --- watching a Kepler orbit live, then running it backward}

The same two-body setup as Example 1 --- but this time we \emph{watch}
it. \kw{scene create} opens the scene window in the browser;
\kw{scene start} sets it in motion; \kw{scene reverse} runs the movie
backward.

\begin{lstlisting}
In[1]:= new point { mass = 1e9, position = [0, 0, 0] }
Out[1]= obj0
In[2]:= new point { mass = 1, position = [0.4, 0, 0], velocity = [0, 2, 0] }
Out[2]= obj1
In[3]:= set system.g_constant = 1e-9
In[4]:= set system.softening = 0
In[5]:= scene create 7878
Out[5]= scene window created: http://127.0.0.1:7878/
(opened in your browser; if no window appeared, open that address yourself)
showing 2 entities; SCENE START begins the evolution — HELP lists all scene commands
In[6]:= scene set_time_step 0.005
Out[6]= scene time step dt = 0.005
In[7]:= scene start
Out[7]= scene playback: running
In[8]:= scene pause
Out[8]= scene playback: paused
In[9]:= scene status
Out[9]= scene: http://127.0.0.1:7878/  (1 window(s) connected)
mode = paused, t = 0.4550000000000003, dt = 0.005, steps = 91, history = 91 frame(s)
entities = 2 (hidden: none)
camera: yaw = -60°, pitch = 55°, dist = 12, target = [0, 0, 0]
In[10]:= scene reverse
Out[10]= scene playback: reversing
In[11]:= scene status
Out[11]= scene: http://127.0.0.1:7878/  (1 window(s) connected)
mode = reversing, t = 0.15500000000000005, dt = 0.005, steps = 91, history = 31 frame(s)
entities = 2 (hidden: none)
camera: yaw = -60°, pitch = 55°, dist = 12, target = [0, 0, 0]
In[12]:= scene events
Out[12]= window connected (1 total)
In[13]:= scene close
Out[13]= scene closed (http://127.0.0.1:7878/)
\end{lstlisting}

\emph{What to notice.} Between \code{In[7]} and \code{In[8]} a few
real seconds passed while the orbit ran in the window --- playback
happens on a background thread at $\sim$30 frames per second, so
\textbf{the notebook prompt never blocks}: cell 9's \kw{status} shows
the \emph{playback copy} has advanced 91 steps to $t \approx 0.455$
while your notebook state still sits untouched at $t = 0$ (that is the
``playback copy'' design: the window animates its own synchronized
copy of your system, so \code{GET system.time} here would still print
\code{0}). After \kw{reverse} (cell 10), the status in cell 11 shows
time flowing \emph{backward} ($t \approx 0.155$) and the history
buffer draining ($91 \to 31$ frames): each reversed frame is an exact
recorded snapshot, so the rewind retraces the orbit bit-for-bit. Cell
12 drains the asynchronous event queue --- the window announced itself
when the browser connected. In JupyterLab you would not even need cell
12: events surface on their own as \code{[scene] ...} lines
(\S\ref{sec:jlab}).

\subsection{Example 12 --- driving the camera from the notebook}

Every gesture the mouse can make in the window has a command twin, so
a \emph{script} can compose the exact view you want --- useful for
repeatable demonstrations and screenshots.

\begin{lstlisting}
In[1]:= new sphere { mass = 3, radius = 0.6, position = [2, 0, 0] }
Out[1]= obj0
In[2]:= new cuboid { mass = 1, half_extents = [0.4, 0.3, 0.2], position = [-2, 0, 1] }
Out[2]= obj1
In[3]:= scene create 7878
Out[3]= scene window created: http://127.0.0.1:7878/
(opened in your browser; if no window appeared, open that address yourself)
showing 2 entities; SCENE START begins the evolution — HELP lists all scene commands
In[4]:= scene translate 2 0
Out[4]= camera target = [2, 0, 0]
In[5]:= scene rotate 30 -10
Out[5]= camera yaw = -30°, pitch = 45°
In[6]:= scene zoom in
Out[6]= camera distance = 9.6
In[7]:= scene zoom 2
Out[7]= camera distance = 4.8
In[8]:= scene zoom out
Out[8]= camera distance = 5.999999999999999
In[9]:= scene hide 0
Out[9]= 1 object(s) hidden
In[10]:= scene show all
Out[10]= 0 object(s) hidden
In[11]:= scene redraw
Out[11]= scene redraw queued for every window
In[12]:= scene status
Out[12]= scene: http://127.0.0.1:7878/  (0 window(s) connected)
mode = stopped, t = 0, dt = 0.01, steps = 0, history = 0 frame(s)
entities = 2 (hidden: none)
camera: yaw = -30°, pitch = 45°, dist = 5.999999999999999, target = [2, 0, 0]
In[13]:= scene close
Out[13]= scene closed (http://127.0.0.1:7878/)
\end{lstlisting}

\emph{What to notice.} The camera is an \textbf{orbit camera}: it
circles a \emph{look-at point} at a \emph{distance}, aimed by
\emph{yaw} (compass direction) and \emph{pitch} (elevation).
\kw{translate} \code{2 0} moves the look-at point onto the sphere
(this is what the arrow keys do in the window); \kw{rotate}
\code{30 -10} adds $30^\circ$ of yaw to the default $-60^\circ$ and
$-10^\circ$ of pitch to the default $55^\circ$ (this is what
left-dragging does); the three zooms multiply the distance by
$1/1.25$, then $1/2$, then $1.25$ --- watch the arithmetic:
$12 \to 9.6 \to 4.8 \to 6$ (this is the mouse wheel). \kw{hide}
\code{0} blanks the sphere out of every connected window without
deleting anything --- \kw{show all} brings it back. Because this
transcript ran headless, cell 12 reports \code{0 window(s) connected}
--- the commands work regardless, and any window that connects later
receives the composed view. Note also the two-argument form
\code{scene translate 2 0}: the optional \code{dz} defaulted to 0, and
\code{-10} in cell 5 was one negative argument, not a subtraction
(scene arguments are \emph{terms} --- the second subtlety of \S3).

\subsection{Example 13 --- the box of shapes: every body type in a rigid, infinitely massive box}

The finale scene: \kw{box} 4 builds the rigid room (\S\ref{sec:box})
and one of every shape goes inside ($R = 1$, $M = 1$) --- a
\textbf{torus} (mass $M$, inner radius 1, outer radius 2), a
\textbf{point} ($M$, the only mover: $v = (100, 200, 100)$), a
\textbf{sphere} ($2M$, $r = \tfrac12$), an ideal zero-thickness
\textbf{disk} ($2M/3$, $r = 1$), a \textbf{cube} ($5M/3$, side 1) and
a \textbf{cylinder} ($2M$, $r = \tfrac12$, height $3/2$). An
axis-aligned torus of outer radius 2 would \emph{exactly inscribe} the
4-box (its outer equator touching four walls), so the torus is tilted
with its axis along $(1,1,1)/\sqrt3$: its extent per axis is
$1.5\sqrt{2/3} + 0.5 \approx 1.7247 < 2$ --- clearance 0.2753. The
other positions are random, drawn by a documented LCG
($x \leftarrow 1664525\,x + 1013904223 \bmod 2^{32}$, seed 20260724,
sequential rejection at $\ge 0.05$ separation) --- see
\code{scripts/collisions/11\_box\_of\_shapes.posim}, the script this
transcript replays.

\begin{lstlisting}
In[1]:= set system.g_constant = 0
In[2]:= box 4
Out[2]= box: inner size 4 x 4 x 4 — six static walls obj0, obj1, obj2, obj3, obj4, obj5 with inverse_mass = 0 (infinitely massive); objects collide elastically off the inside faces
In[3]:= new torus { mass = 1, inner_radius = 1, outer_radius = 2, orientation = [0.888073833977115, -0.325057583671868, 0.325057583671868, 0] }
Out[3]= obj6
In[4]:= new point { mass = 1, position = [1.406590, -0.995859, 0.569601], velocity = [100, 200, 100] }
Out[4]= obj7
In[5]:= new sphere { mass = 2, radius = 1/2, position = [1.424704, 1.367496, -0.493612] }
Out[5]= obj8
In[6]:= new disk { mass = 2/3, radius = 1, position = [-0.677386, -1.041493, -1.091679], orientation = [0.900447102352677, 0.307567078752479, -0.307567078752479, 0] }
Out[6]= obj9
In[7]:= new cuboid { mass = 5/3, half_extents = [0.5, 0.5, 0.5], position = [1.074397, 0.816223, 1.102099] }
Out[7]= obj10
In[8]:= new cylinder { mass = 2, radius = 1/2, height = 3/2, position = [-1.027024, -1.403890, -0.485619], orientation = [0.968912421710645, 0, 0.247403959254523, 0] }
Out[8]= obj11
In[9]:= list
Out[9]= obj0: cuboid he=[1, 4, 4], mass=0, charge=0, pos=[3, 0, 0] [wall: static, inverse_mass=0]
obj1: cuboid he=[1, 4, 4], mass=0, charge=0, pos=[-3, 0, 0] [wall: static, inverse_mass=0]
obj2: cuboid he=[4, 1, 4], mass=0, charge=0, pos=[0, 3, 0] [wall: static, inverse_mass=0]
obj3: cuboid he=[4, 1, 4], mass=0, charge=0, pos=[0, -3, 0] [wall: static, inverse_mass=0]
obj4: cuboid he=[4, 4, 1], mass=0, charge=0, pos=[0, 0, 3] [wall: static, inverse_mass=0]
obj5: cuboid he=[4, 4, 1], mass=0, charge=0, pos=[0, 0, -3] [wall: static, inverse_mass=0]
obj6: torus ring=1.5 tube=0.5, mass=1, charge=0, pos=[0, 0, 0]
obj7: point, mass=1, charge=0, pos=[1.40659, -0.995859, 0.569601]
obj8: sphere r=0.5, mass=2, charge=0, pos=[1.424704, 1.367496, -0.493612]
obj9: disk r=1, mass=0.6666666666666666, charge=0, pos=[-0.677386, -1.041493, -1.091679]
obj10: cuboid he=[0.5, 0.5, 0.5], mass=1.6666666666666667, charge=0, pos=[1.074397, 0.816223, 1.102099]
obj11: cylinder r=0.5 h=1.5, mass=2, charge=0, pos=[-1.027024, -1.40389, -0.485619]
In[10]:= collide
Out[10]= collisions ON (51 collidable pair(s); 0 impulse(s) so far)
In[11]:= get obj0.inverse_mass
Out[11]= 0
In[12]:= get obj6.inertia_tensor
Out[12]= [[1.28125, 0, 0], [0, 1.28125, 0], [0, 0, 2.4375]]
In[13]:= energy
Out[13]= 30000
In[14]:= momentum
Out[14]= [100, 200, 100]
In[15]:= run 0.1 steps 100
Out[15]= t = 0.1 (2121 solver steps, 100 snapshots, |dE/E| = 2.040e-10, 119 collision(s) — CONTACTS lists them)
In[16]:= energy
Out[16]= 29999.999993880127
In[17]:= momentum
Out[17]= [146.48911126803657, 102.1478121131382, 46.01601291636828]
In[18]:= get system.collisions
Out[18]= 119
In[19]:= get system.box
Out[19]= 4
\end{lstlisting}

\emph{What to notice.} Cell 3 sizes the torus by the order-independent
\code{inner\_radius} + \code{outer\_radius} pair; cell 8 uses
\code{height = 3/2} --- the \emph{full} height, so \code{half\_height}
reads 0.75. \kw{collide} counts \textbf{51 collidable pairs}: 12
objects make 66 pairs, minus the 15 wall--wall pairs (two static
bodies can never collide). \code{Out[12]} is the torus's analytic
inertia $\mathrm{diag}(1.28125, 1.28125, 2.4375)$ --- exactly
$I_{xy} = m(\tfrac12 c^2 + \tfrac58 a^2)$,
$I_z = m(c^2 + \tfrac34 a^2)$ for $c = 1.5$, $a = 0.5$. The energy
starts at $E_0 = \tfrac12\cdot 1\cdot|v|^2 = 30000$ \textbf{exactly}
and, after \textbf{119 collisions} in 0.1\,s, is conserved to
$|dE/E| = 2.040\times10^{-10}$ --- but momentum went from
$(100, 200, 100)$ to $(146.49, 102.15, 46.02)$: \textbf{not
conserved}, because the infinitely massive walls absorb momentum
without moving --- the physical signature of infinite mass. Along the
way the point particle can \emph{thread the torus hole} and passes
through the ideal zero-thickness disk (a measure-zero contact) --- the
ball-vs-anything collision tier is exact SDF geometry, not an
approximation --- while every finite-size body bounces off both. The
walls end the run bit-identically at rest.

\subsection{Example 14 --- two dumbbells from a user-defined function}

The two-release finale: a user function (\S\ref{sec:userdefs}) that
builds \textbf{named dumbbells} (\S\ref{sec:new}), called twice to
launch a pair of tumbling dumbbells at each other, off-center and
spinning. No walls this time --- so \emph{all three} conserved
quantities must survive the collisions. The script is
\code{scripts/collisions/12\_two\_dumbbells.posim}; this transcript
replays it.

\begin{lstlisting}
In[1]:= set system.g_constant = 0
In[2]:= def create_dumbell(name, m1 = 1, m2 = 1, m_rod = 0.5, r1 = 0.25, r2 = 0.25, rod_radius = 0.1, length = 1, position = [0, 0, 0], velocity = [0, 0, 0], angular_velocity = [0, 0, 0]) {
  new dumbbell as name { m1 = m1, m2 = m2, m_rod = m_rod, r1 = r1, r2 = r2, rod_radius = rod_radius, length = length, position = position, velocity = velocity, angular_velocity = angular_velocity }
}
Out[2]= function create_dumbell(11 parameter(s)) defined — 1 body line(s)
In[3]:= create_dumbell("dumbell0", 1, 2, 0.5, 0.25, 0.25, 0.1, 1, [-2, 0.15, 0], [1.5, 0, 0], [0, 0, 0.6])
Out[3]= obj0 as dumbell0
In[4]:= create_dumbell("dumbell1", 2, 1, 0.4, 0.3, 0.2, 0.08, 1.2, [2, -0.15, 0], [-1.5, 0, 0], [0.4, 0, 0])
Out[4]= obj1 as dumbell1
In[5]:= list
Out[5]= obj0: dumbbell r1=0.25 r2=0.25 rod_r=0.1 len=1, mass=3.5, charge=0, pos=[-2, 0.15, 0]
obj1: dumbbell r1=0.3 r2=0.2 rod_r=0.08 len=1.2, mass=3.4, charge=0, pos=[2, -0.15, 0]
In[6]:= get dumbell0.m1
Out[6]= 1
In[7]:= get dumbell1.m_rod
Out[7]= 0.3999999999999999
In[8]:= get dumbell0.vx
Out[8]= 1.5
In[9]:= energy
Out[9]= 7.865310611764706
In[10]:= momentum
Out[10]= [0.15000000000000036, 0, 0]
In[11]:= angmom
Out[11]= [0.4443030588235295, 0, -1.5059999999999998]
In[12]:= run 3 steps 60
Out[12]= t = 3 (3861 solver steps, 60 snapshots, |dE/E| = 9.463e-11, 2 collision(s) — CONTACTS lists them)
In[13]:= energy
Out[13]= 7.865310611020375
In[14]:= momentum
Out[14]= [0.15000000000000036, 0, 0]
In[15]:= angmom
Out[15]= [0.44430305882373156, -0.000000000000009547918011776346, -1.505999999999855]
In[16]:= get system.collisions
Out[16]= 2
\end{lstlisting}

\emph{What to notice.} Cell 2 is one \kw{def} whose single body line
does everything: \code{new dumbbell as name \{ ... \}} forwards all
eleven parameters, and because \code{name} arrives as a string
(\code{"dumbell0"}, \code{"dumbell1"}), each call registers the name
it was given --- \code{Out[3]}/\code{Out[4]} show both handles,
\kw{list} shows the part geometry, and cells 6--8 read the parts and
the \code{.vx} shorthand through the names. (\code{Out[7]} reads 0.4
back as \code{0.3999999999999999}: the parts are stored as mass
\emph{fractions} of the total, so one reconstruction rounding step
shows through --- 15 correct digits.) Then the physics: two tumbling,
asymmetric rigid bodies meet off-center, collide \textbf{twice} (cell
16), and every conserved quantity survives the real CVODE event
handling. Energy: $7.865310611764706 \to 7.865310611020375$ --- the
run banner prints \code{|dE/E| = 9.463e-11}. Momentum:
\code{[0.15000000000000036, 0, 0]} \textbf{bit-identical}. Angular
momentum about the origin --- the delicate one, orbital $r \times p$
plus spin, exchanged between the two at every off-center impact:
conserved to $\sim\!10^{-13}$ per component (\code{Out[11]} vs.\
\code{Out[15]}), because the part-wise exact narrow phase applies the
action--reaction impulse pair at one shared contact point. Contrast
Example 13: there the infinitely massive walls absorbed momentum;
here, with no walls, $E$, $P$ \emph{and} $L$ all hold at solver
precision.

\subsection{Example 15 --- a particle in a box, solved inside the language}

The special functions are not decoration: with \code{eigenvalues} you
can set up and solve a small quantum problem in the notebook itself.

Take a particle in an infinite square well of width $L=1$ with
$\hbar=m=1$. Discretise $-\tfrac12\psi''=E\psi$ on 4 interior points,
spacing $h=0.2$. The second-derivative stencil puts $1/h^2$ on the
diagonal and $-1/2h^2$ beside it, so the Hamiltonian is a $4\times4$
tridiagonal matrix you can type out. The exact answer is
$E_n=n^2\pi^2/2$, so $E_1=\pi^2/2$.

\begin{lstlisting}
In[1]:= let h = 0.2
Out[1]= h set
In[2]:= let k = 1 / (h * h)
Out[2]= k set
In[3]:= eigenvalues([[k, -0.5*k, 0, 0], [-0.5*k, k, -0.5*k, 0], [0, -0.5*k, k, -0.5*k], [0, 0, -0.5*k, k]])
Out[3]= [4.774575140626312, 17.274575140626325, 32.72542485937371, 45.225424859373675]
In[4]:= pi * pi / 2
Out[4]= 4.934802200544679
\end{lstlisting}

The ground state comes out at 4.7746 against an exact 4.9348 ---
\textbf{3.2\,\% low}. That is not a bug, and it is worth understanding
rather than tightening a tolerance until it goes away. A 3-point stencil
on 4 points is a \emph{coarse} grid, and its eigenvalues have the closed
form $k\bigl(1-\cos(n\pi/5)\bigr)$ with $k=1/h^2=25$ --- which
reproduces all four numbers above exactly
($25(1-\cos36^\circ)=4.7746$). The discretisation error is real, known,
and second order in $h$: refine the grid and it falls off as $h^2$.

Note the row shapes. Each row here has four entries, so it lexes as a
\emph{quaternion} rather than a list --- the bracket literal is
overloaded (see \S\ref{sec:specialfns}). It works anyway, because every
bracket shape is accepted where a numeric list is wanted. You should not
have to think about quaternions to type a matrix.

The rest of the family is there to be checked against its own
identities, which is the honest way to use a special-function library
you did not write:

\begin{lstlisting}
In[5]:= let t = 0.7
Out[5]= t set
In[6]:= chebyshev_t(5, cos(t)) - cos(5 * t)
Out[6]= -0.00000000000000011102230246251565
In[7]:= sph_j(0, 1.3) - sin(1.3) / 1.3
Out[7]= 0
In[8]:= legendre_p(3, 0.4) - 0.5 * (5 * 0.4 * 0.4 * 0.4 - 3 * 0.4)
Out[8]= 0.00000000000000011102230246251565
In[9]:= bessel_j_array(4, 2)
Out[9]= [0.22389077914123567, 0.5767248077568734, 0.35283402861563773, 0.12894324947440208, 0.03399571980756844]
\end{lstlisting}

\code{In[6]} is $T_n(\cos\theta)=\cos n\theta$; \code{In[7]} is
$j_0(x)=\sin x/x$; \code{In[8]} is $P_3(x)=(5x^3-3x)/2$. All three land
at zero or one part in $10^{16}$ --- a single rounding step.
\code{In[9]} returns the whole table $J_0\ldots J_4$ in one call, which
is what a Bessel-expanded propagator actually needs.

Finally, the thing the library refuses to do:

\begin{lstlisting}
In[10]:= hermite_h(2.5, 1)
Err[10]: hermite_h(): argument 1 must be a whole number (an integer order), got 2.5
\end{lstlisting}

There is no Hermite polynomial of order 2.5. The alternative ---
silently computing $H_2$ --- would hand you a plausible number with no
indication anything went wrong, and you would carry it through the rest
of the calculation.

\subsection{Example 16 --- quantum mechanics: bound states, then tunnelling}

The \kw{qm} family solves an actual quantum problem in the notebook.
Start with the harmonic oscillator, whose answer everyone knows: with
$\hbar=m=\omega=1$ the spectrum is $E_n=n+\tfrac12$. The potential is an
ordinary user function.

\begin{lstlisting}
In[1]:= def v(x) { 0.5 * x * x }
Out[1]= function v(1 parameter(s)) defined -- 1 body line(s)
In[2]:= qm grid -8 8 250
Out[2]= grid [-8, 8] with 250 interior points, h = 0.063745 (potential and psi cleared)
In[3]:= qm potential v
Out[3]= potential `v` sampled at 250 points, V in [0.000507928445580228, 31.49207155441977] (psi cleared)
In[4]:= qm states 4
Out[4]= 4 lowest bound state(s):
  E[0] = 0.499872985615
  E[1] = 1.499364798834
  E[2] = 2.498348101796
  E[3] = 3.496822505539
\end{lstlisting}

0.4999, 1.4994, 2.4983, 3.4968 against an exact 0.5, 1.5, 2.5, 3.5. The
error grows with $n$ --- that is not noise, it is the three-point
stencil: higher states oscillate faster and a coarse grid resolves them
less well, with the error tracking $2n^2+2n+1$.

Now the sharpest test available. A stationary state is \emph{stationary},
so loading one and propagating it must change nothing:

\begin{lstlisting}
In[5]:= qm state 1
Out[5]= psi = bound state 1, E = 1.499364798834, t reset to 0
In[6]:= qm energy
Out[6]= 1.499364798833686
In[7]:= qm run 4 steps 400
Out[7]= t = 4 (400 step(s) of dt = 0.01), <E> = 1.499364798834, norm drift = 4.219e-15
In[8]:= qm energy
Out[8]= 1.499364798833692
\end{lstlisting}

Four hundred Crank--Nicolson steps moved the energy by $6\times10^{-15}$
--- one bit --- and the norm by $4\times10^{-15}$. The Cayley operator is
unitary for \emph{any} time step, so that drift measures the linear
solver rather than \code{dt}; and the energy holding still couples the
eigensolver to the propagator, so it would break if either were wrong or
if they disagreed about the Hamiltonian.

Now tunnelling, the problem with no classical analogue. A packet of
central energy $E_0=k_0^2/2=2$ is fired at a barrier of height 2.5 ---
\emph{above} its energy, so classically nothing gets through:

\begin{lstlisting}
In[1]:= qm grid -100 100 2000
Out[1]= grid [-100, 100] with 2000 interior points, h = 0.099950 (potential and psi cleared)
In[2]:= qm potential barrier 2.5 0 1
Out[2]= potential `barrier 2.5 on [0, 1]` sampled at 2000 points, V in [0, 2.5] (psi cleared)
In[3]:= qm packet -25 2 2
Out[3]= psi = Gaussian packet at x0 = -25, sigma = 2, k0 = 2, t = 0
In[4]:= qm run 30 steps 3000
Out[4]= t = 30 (3000 step(s) of dt = 0.01), <E> = 2.023971780446, norm drift = 3.018e-13
In[5]:= qm prob 1 100
Out[5]= 0.327563019391693
In[6]:= qm prob -100 0
Out[6]= 0.672436408645103
\end{lstlisting}

\textbf{33\,\% of the particle got through a barrier it did not have the
energy to cross.} The two channels add to 1.000000 --- nothing was lost,
and nothing reached a wall (there was no warning). The domain is 200
wide precisely so the reflected packet has nowhere to bounce from within
$t=30$.

Note the grid sizes. Bound states used 250 points because that
calculation diagonalises a dense matrix and costs $O(n^3)$; scattering
used 2000 because propagation is $O(n)$ per step and can afford the
resolution. Choosing both to be the same would waste one or starve the
other.

\subsection{Example 17 --- tunnelling, with the barrier as a user function}

Comparison operators exist so that a piecewise potential can be an
ordinary function. This is the canonical 1-D quantum problem written
that way, and it is the notebook
\code{dynamic\_notebooks/tunneling.posim}.

\begin{lstlisting}
In[1]:= def barrier(x) { 2.5 * (x > 0) * (x < 1) }
Out[1]= function barrier(1 parameter(s)) defined -- 1 body line(s)
In[2]:= barrier(-1)
Out[2]= 0
In[3]:= barrier(0.5)
Out[3]= 2.5
In[4]:= barrier(2)
Out[4]= 0
\end{lstlisting}

\code{(x > 0) * (x < 1)} is 1 inside the interval and 0 outside --- an
indicator function built from arithmetic on truth values.

A packet of central energy $E_0=k_0^2/2=2$ is fired at a barrier of
height 2.5, \emph{above} its energy:

\begin{lstlisting}
In[5]:= qm grid -100 100 2000
Out[5]= grid [-100, 100] with 2000 interior points, h = 0.099950 (potential and psi cleared)
In[6]:= qm potential barrier
Out[6]= potential `barrier` sampled at 2000 points, V in [0, 2.5] (psi cleared)
In[7]:= qm packet -25 2 2
Out[7]= psi = Gaussian packet at x0 = -25, sigma = 2, k0 = 2, t = 0
In[10]:= qm run 30 steps 3000
Out[10]= t = 30 (3000 step(s) of dt = 0.01), <E> = 2.023971780446, norm drift = 3.018e-13
In[11]:= qm prob 1 100
Out[11]= 0.327563019391693
In[12]:= qm prob -100 0
Out[12]= 0.672436408645103
\end{lstlisting}

\textbf{33\,\% crossed a barrier it had no classical right to cross},
and the two channels account for all but $10^{-6}$ of the probability.

Note that \code{qm potential barrier} picked the \emph{user's} function,
not the built-in shape of the same name. A bare name is always yours;
the built-in needs arguments.

\subsubsection*{Watching it happen}

\begin{lstlisting}
In[14]:= qm animate "scatter.html" 32 frames 140
Out[14]= wrote scatter.html -- 140 frames over t = 32 (dt = 0.011429, 667 points per frame), worst norm drift 5.822e-13. Open it in a browser.
\end{lstlisting}

A self-contained page --- nothing fetched from the network --- showing
$|\psi|^2$ against $x$ with the potential overlaid, play/pause, a scrub
bar, and a live norm and transmitted readout computed in the browser.
Its final transmitted value is 0.327562 against the notebook's 0.327563:
two independent implementations of the same integral.

\subsubsection*{The same answer on half the domain}

\begin{lstlisting}
In[15]:= qm grid -45 45 1350
In[17]:= qm absorb 15 3
Out[17]= absorbing edges: width 15, strength 3, power 2. Propagation is NO LONGER unitary.
In[19]:= qm run 20 steps 2000
Out[19]= t = 20 (2000 step(s) of dt = 0.01), <E> = 2.027672292154, norm drift = 1.056e-4
In[20]:= qm prob 1 30
Out[20]= 0.327689667523621
\end{lstlisting}

0.327690 against 0.327563 --- four parts in ten thousand, on a domain
\textbf{less than half the size}.

\subsubsection*{A double barrier, and an honest negative result}

\begin{lstlisting}
In[24]:= def double(x) { 2.5 * ((x > 0) * (x < 1) + (x > 3) * (x < 4)) }
In[31]:= qm prob 4 100
Out[31]= 0.271498902829932
\end{lstlisting}

Two barriers with a gap form a resonant cavity. At this energy it
transmits \textbf{0.2715, less than the single barrier's 0.3276}, with
about 0.8\,\% of the probability left in the gap. That trapped fraction
\emph{is} the cavity; but resonant enhancement occurs only at its
quasi-bound energies, and those peaks are narrower than this packet's
momentum spread ($\sigma=2$ gives $\Delta k=0.25$), which averages
straight over them. A scan over $k_0$ from 1.4 to 2.9 rises
monotonically with no peak. See \code{TUNNELING\_RESULTS.md}.

\subsection{Example 18 --- the double slit, in two dimensions}

The canonical 2-D quantum problem, and the notebook
\code{dynamic\_notebooks/double\_slit.posim}. The wall is an ordinary
user function of two arguments, built entirely from comparisons:

\begin{lstlisting}
In[1]:= def openings(y) { (y > 1.5) * (y < 2.5) + (y > -2.5) * (y < -1.5) }
In[2]:= def slit(x, y) { 60 * (x > 0) * (x < 0.4) * (1 - openings(y)) }
In[3]:= slit(0.2, 0)
Out[3]= 60
In[4]:= slit(0.2, 2)
Out[4]= 0
\end{lstlisting}

\begin{lstlisting}
In[8]:= qm2 grid -12 24 450, -16 16 400
In[10]:= qm2 absorb 4, 8
In[11]:= qm2 packet -4.5, 0, 1.2 4, 8 0
In[12]:= qm2 energy
Out[12]= 31.008336923103123
In[15]:= qm2 run 2.6 steps 800
Out[15]= t = 2.6 (800 ADI step(s) of dt = 0.00325), <E> = 30.990650438623, norm drift = 7.796e-1
In[18]:= qm2 norm
Out[18]= 0.220410865359174
\end{lstlisting}

Those three lines together say something worth pausing on:
\textbf{78\,\% of the probability was absorbed at the walls, and the
energy moved by 0.06\,\%}. The absorber removed the reflected wave ---
the wall is 60 high against a packet of energy 32, so most of it bounces
--- without touching the energy of what remains. $\langle E\rangle$ is
$\langle\psi|H|\psi\rangle/\langle\psi|\psi\rangle$, divided by the norm
on purpose; undivided it would have fallen with the norm and looked like
an energy leak.

The fringes, counted in bands on a screen at $11<x<19$:

\begin{center}
\begin{tabular}{lll}
\textbf{band} & \textbf{P} & \textbf{expected}\\
\hline
$0\pm0.5$ & \textbf{0.01577} & central maximum\\
0.5--1.5 & 0.00452 & first minimum, $y=1.46$\\
2.5--3.5 & \textbf{0.01742} & first maximum, $y=2.96$\\
3.5--4.5 & 0.00471 & second minimum, $y=4.56$\\
5.5--6.5 & \textbf{0.01399} & second maximum, $y=6.32$\\
\end{tabular}
\end{center}

With slit separation $d=4$ and $\lambda=2\pi/k=0.785$, maxima sit at
$\sin\theta = m\lambda/d$ with the screen $L\approx14.8$ from the slits,
giving 2.96 and 6.32. \textbf{Every predicted maximum and minimum lands
in the right band.}

Two corrections got there. A first attempt used $k=4$, $d=2$, putting the
first maximum at $52^\circ$ --- off the screen, so the scan showed a
featureless lobe. A second launched the packet \emph{inside} the
absorber, and 86\,\% of it was eaten before reaching the slits.

\begin{lstlisting}
In[28]:= qm2 animate "double_slit.html" 3 frames 100
\end{lstlisting}

A self-contained heat-map page. Its own in-page analysis of the final
frame finds fringe maxima at $y=0.04$ and $y=3.23$, matching the band
scan and the theory independently of the Rust that produced it.

\section{Error message tour}

\begin{center}
\begin{tabular}{p{5.4cm}p{9.2cm}}
\textbf{you type} & \textbf{you get}\\\hline
\code{get obj7.mass} (no obj7) & \code{no object obj7}\\
\code{get obj0.bogus} & \code{unknown object field `bogus' --- see HELP for the field list}\\
\code{[1,0,0] * [0,1,0]} & \code{cannot multiply vec3 by vec3 (for vec3*vec3 use dot()/cross())}\\
\code{set = 3} & \code{parse error at column 5: expected a path root (`objN` or `system`), found `=`}\\
\code{run 1} under SPRK with $B\ne0$ & \code{SPRK method requires a separable Hamiltonian: magnetic field B must be zero ... use METHOD ADAMS or BDF}\\
\code{bogusname} & \code{unknown name `bogusname` (define it with LET, pass it as a function parameter, or use `bogusname.field` for a registered object)}\\
\code{get d.m1} (only \code{ball} is registered) & \code{no object named `d` (registered names: ["ball"]; NEW ... AS <name> creates one; positional paths are objN.field, contactK.field, system.field)}\\
\code{new sphere as d} (name taken) & \code{the name `d` already refers to obj0 --- DEL it or pick another name}\\
\code{new sphere as system} & \code{`system` is reserved}\\
\code{show nosuch} & \code{no user function `nosuch` --- FUNCS lists the defined ones}\\
\end{tabular}
\end{center}

Every error names the column or the exact field/feature at fault, and
never aborts the session.

\section{The numerical engine underneath (SUNDIALS 7.8.0)}

You never name a solver library when you type \kw{step} or \kw{run} ---
but it is worth one page knowing what answers you.

\textbf{SUNDIALS} is Lawrence Livermore National Laboratory's suite of
differential-equation solvers, in production use in physics codes since
the 1990s. \code{sundials\_rs/} in this repository is a \emph{pure-Rust
translation of SUNDIALS 7.8.0}: no \code{unsafe}, no FFI, no C compiler,
no crates.io dependency. It is vendored byte-for-byte from
\code{once-ere/SUNDIALS\_7\_8\_Rust\_port\_for\_Linux@780b916}; the
upgrade from the 7.7.0 engine the predecessor repository used is
recorded, with its evidence, in \code{PORT\_7.8.0\_PROVENANCE.md}.

\subsection{Which solver runs which command}

\begin{center}
\begin{tabular}{p{4.6cm}p{7.4cm}p{2.2cm}}
\textbf{you type} & \textbf{what runs} & \textbf{from}\\\hline
\kw{method adams} (default) & CVODE, variable-order variable-step
Adams--Moulton, Newton iteration, dense difference-quotient Jacobian &
\code{cvode\_rs}\\
\kw{method bdf} & CVODE BDF --- the stiff-problem family &
\code{cvode\_rs}\\
\kw{method sprk} \emph{table} [\emph{dt}] & ARKODE SPRKStep, symplectic
partitioned Runge--Kutta at a fixed step & \code{arkode\_rs}\\
\kw{collide on} with a collidable pair & the same solver with
\emph{rootfinding} armed on the pairwise signed separations; a downward
zero crossing \emph{is} the moment of impact &
\code{CVodeRootInit}\\
\end{tabular}
\end{center}

There is deliberately no fourth option. Every trajectory this simulator
prints came out of an error-controlled or symplectic solver; the legacy
Euler and Verlet steppers were not carried over.

\subsection{What ``rootfinding'' buys you}

A naive collision check samples the geometry once per output frame and
asks ``are they overlapping now?''. If the ball moved farther than the
wall is thick, the answer is no on both sides and the ball passes
through. Rootfinding asks a different question: it hands the solver a
continuous function --- here, the signed gap between two surfaces ---
and the solver \emph{interpolates back to the instant that function
crossed zero}. The bounce happens at the exact time of impact, not at
the next frame boundary.

\begin{lstlisting}
In[1]:= new sphere { mass = 1, radius = 0.5, position = [0, 4.5, 0] }
In[2]:= new cuboid { mass = 1, half_extents = [10, 0.05, 10],
                     position = [0, -0.05, 0], inverse_mass = 0 }
In[3]:= set system.uniform_gravity = [0, -10, 0]
In[4]:= collide on
In[5]:= run 1 steps 100
Out[5]= t = 1 (209 solver steps, 100 snapshots, |dE/E| = 2.890e-11,
        1 collision(s) - CONTACTS lists them)
In[6]:= contacts
Out[6]= contact0: obj0 <-> obj1 at t = 0.8944271909999157
  point  = [0, 0, 0]
  normal = [0, -1, 0]  (from obj0 toward obj1)
  approach speed = 8.94427190999916, impulse = 17.88854381999832
\end{lstlisting}

Work out the answer by hand. The plate's top face is at $y=0$ (centre
$-0.05$, half-thickness $0.05$); the ball's centre starts at $4.5$ and
its radius is $0.5$, so the surfaces meet after the centre falls $4.0$.
Under $g=10$ that is $t=\sqrt{2\cdot 4/10}=\sqrt{0.8}$:

\begin{lstlisting}
analytic  sqrt(0.8) = 0.8944271909999159
measured            = 0.8944271909999157
relative error      = 1.24e-16
\end{lstlisting}

\textbf{One part in $10^{16}$} --- a single unit in the last place of a
64-bit float. Nothing about the output cadence produced that: the run
asked for 100 snapshots, and none of them falls anywhere near
$0.894\ldots$. The solver interpolated to the crossing.

\subsection{All six families, and what reaches each}

\begin{center}
\begin{tabular}{p{2.6cm}p{6.6cm}p{4.6cm}}
\textbf{crate} & \textbf{solves} & \textbf{reached by}\\\hline
\code{cvode\_rs}  & ODE initial-value problems (Adams / BDF) & \kw{method adams}, \kw{method bdf}\\
\code{arkode\_rs} & Runge--Kutta: explicit, implicit, IMEX, multirate, symplectic & \kw{method sprk}\\
\code{ida\_rs}    & differential-algebraic systems $F(t,y,\dot y)=0$ & \kw{constrain} + \kw{method ida}\\
\code{kinsol\_rs} & nonlinear algebraic systems, Anderson acceleration & \kw{equilibrium}\\
\code{cvodes\_rs} & CVODE plus forward and adjoint sensitivity analysis & \kw{sensitivity}\\
\code{idas\_rs}   & IDA plus sensitivity analysis & \kw{sensitivity} while constrained\\
\end{tabular}
\end{center}

Every one of them is diffed against the upstream C reference output ---
see \code{sundials\_rs/VERIFICATION.md}. The last four arrived with the
7.8.0 engine and were wired into the language afterwards; the shape of
that wiring is in ARCHITECTURE.md \S3.9.

\section{Constrained dynamics, equilibrium and sensitivity}

Everything so far answers \emph{``what happens next?''}. Three more
commands ask different questions, each answered by a different solver.

\begin{center}
\begin{tabular}{p{5.4cm}p{5.6cm}p{2.8cm}}
\textbf{you type} & \textbf{question} & \textbf{solver}\\\hline
\kw{step} / \kw{run} & what happens next? & CVODE / ARKODE\\
\kw{constrain} + \kw{method ida} + \kw{run} & \dots with this geometry held exactly & IDA\\
\kw{equilibrium} & where does it come to rest? & KINSOL\\
\kw{sensitivity} & how much does the answer depend on an input? & CVODES / IDAS\\
\end{tabular}
\end{center}

\subsection{The four joints}

\begin{center}
\begin{tabular}{p{4.6cm}p{1.2cm}p{4.4cm}p{3.4cm}}
\textbf{command} & \textbf{rows} & \textbf{holds} & \textbf{freedoms left}\\\hline
\kw{constrain} \emph{a b} [\emph{len}] & 1 & a fixed distance & 5\\
\kw{ball} \emph{a b} & 3 & a shared point & 3\\
\kw{universal} \emph{a b u w} & 4 & a shared point, two shafts square & 2\\
\kw{hinge} \emph{a b axis} & 5 & a shared point and an axis & 1 (the swing)\\
\end{tabular}
\end{center}

The last three grip \emph{orientation}, which is why \kw{method ida}
carries the full 13-numbers-per-object state rather than positions
alone. \textbf{The pivot is the midpoint of the two bodies as they stand
when you make the joint}, carried into each body's frame --- the same
``freeze what you have'' rule a bare \kw{constrain} follows, so the joint
is satisfied the instant it is made.

A door is an anchor, a slab and a hinge:

\begin{lstlisting}
In[4]:= new sphere as jamb { mass = 1, radius = 0.02,
                             position = [0, 0, 0], inverse_mass = 0 }
In[5]:= new cuboid as door { mass = 1, half_extents = [0.2, 0.4, 0.2],
          position = [0.0199986666933331, -0.9998000066665778, 0] }
In[6]:= hinge jamb door [0, 0, 1]
In[9]:= run 1 steps 10
Out[9]= t = 1 (70 solver steps, 10 snapshots, |dE/E| = 1.613e-9)
In[10]:= constraints
Out[10]= constraint0: hinge obj0 <-> obj1, 5 row(s)
worst |g| = 2.73750133672479e-10, worst |g_dot| = 2.3769240437118873e-9
In[11]:= get door.angular_momentum
Out[11]= [0, 0, 0.003742219622967182]
\end{lstlisting}

The door swings about \emph{z and nothing else} --- the two rows a hinge
has over a ball joint are exactly the ones that forbid the other two
axes --- and the joint is held to $2.7\times10^{-10}$ after 70 steps. A
hinged rigid body is a \emph{compound} pendulum, with period
$T = 2\pi\sqrt{I_{\text{pivot}}/(mgd)}$ and
$I_{\text{pivot}} = I_{\text{com}} + md^2$; the simulator reproduces it
to about $3\times10^{-8}$ of a full swing.

\textbf{A joint constrains velocity, not just position.} A ball joint
says the two bodies share a point, so at the velocity level it says
$v_i + \omega_i\times r_i = v_j + \omega_j\times r_j$ --- a body turning
about a pivot offset from its centre \emph{must have its centre moving}.
Give it a spin and leave its velocity at zero and the state is not on the
constraint manifold at all. The run \textbf{projects} the starting
velocities onto it --- the smallest mass-weighted change that satisfies
the joint, which is exactly the impulse a real coupling delivers when
clutched onto a spinning shaft --- and reports how big that change was.
An already-consistent state is left exactly alone. (A rod has no angular
Jacobian, so spin never enters its $\dot g$; that is why rods carried
spinning bodies from the start and hid this.)

\textbf{One real limit.} Orientation joints carry a tolerance floor of
$\texttt{rtol}=10^{-6}$, because the differential-algebraic system a
hinge produces is \emph{index 2} and has an accuracy ceiling no tolerance
can push past --- measured across twelve compound pendulums, $10^{-6}$
converges in every one and $10^{-8}$ in none. \kw{run} reports when the
floor was applied.

\kw{ball}, \kw{hinge} and \kw{universal} are \emph{contextual} keywords
--- commands only at the start of a line --- so \code{new sphere as ball}
and \code{get ball.mass} keep working.

\subsection{\kw{constrain} --- a rod, not a spring}

\begin{lstlisting}
CONSTRAIN <a> <b> [length]   rigid rod between two objects
CONSTRAIN OFF                drop every rod
CONSTRAINTS                  list them, and how well they are held
\end{lstlisting}

\code{<a>} and \code{<b>} are objects, written either positionally
(\code{obj0}) or by a name registered with \kw{new} \dots\ \kw{as}.
\textbf{With no length the rod freezes the separation the two bodies
already have}, which is always consistent and is what you almost always
want.

You could model a rod as a very stiff spring. That is not the same
thing: a stiff spring is an approximation that vibrates, needs a tiny
step, and still lets the length wander. A constraint is an
\emph{algebraic equation} the motion must satisfy exactly. Adding one
changes the problem from an ODE into a \textbf{differential-algebraic
equation}, and a constrained system refuses every method but IDA, by
name:

\begin{lstlisting}
In[7]:= run 1 steps 5
Err[7]: this system has 1 rigid constraint(s), which only the DAE
        integrator can hold: use METHOD IDA (or remove them with
        CONSTRAIN OFF). The current method is Adams
\end{lstlisting}

\textbf{A body with \code{inverse\_mass = 0} is an anchor.} It never
moves and it absorbs the rod's reaction. Anchor + bob + rod is a
pendulum.

\textbf{What is held, and how well.} \kw{constraints} reports the worst
$|g|$ (how far the rod is from its length) and $|\dot g|$ (how fast that
is changing). Both stay at roundoff for the whole run, because the
formulation carries them \emph{both} as equations. A cheaper scheme that
constrains only the acceleration lets $g$ drift quadratically, and by
the time you notice, the answer is quietly wrong.

\textbf{Scope.} Constraints act on positions. A spinning rigid body, or
an external torque, is refused with a message naming it.

\subsection{\kw{equilibrium} --- where does it come to rest?}

Finds a configuration where every free body has zero net force, every
anchor is where it started, and every rod is the right length; moves the
bodies there and stops them. \code{system.time} is untouched --- this is
not an integration.

\begin{lstlisting}
In[7]:= equilibrium
Out[7]= equilibrium found in 17 Newton iteration(s), 67 residual
        evaluation(s); largest net force on any free body =
        7.459152323898993e-13, worst |g| = 1.9317880628477724e-14
\end{lstlisting}

The starting guess is wherever the bodies already are, so it finds
\emph{the nearby} rest state. \textbf{It says nothing about stability}
--- a pencil balanced on its point is an equilibrium. Perturb the answer
and \kw{run} if you need to know. And a system where every body is free
has no isolated equilibrium at all (translate it and nothing changes);
the refusal says exactly that.

\subsection{\kw{sensitivity} --- how much does the answer depend on the input?}

\begin{lstlisting}
SENSITIVITY <t> "<param>" ["<param>" ...]
\end{lstlisting}

Runs for \code{<t>} \emph{and} reports $\partial(\text{state})/\partial
p$ for each parameter. Names are quoted strings because \code{mass 0} is
two tokens: \code{"g\_constant"}, \code{"mass 0"}, \code{"charge 1"},
\code{"gravity.y"}, \code{"e\_field.x"}, \code{"b\_field.z"}.

Running twice with slightly different inputs and subtracting is the
difference of two nearly equal numbers and loses most of its digits.
This integrates the derivative \emph{alongside} the state instead.
CVODES normally; IDAS when the system is constrained, because only then
are the equations a DAE.

Free fall gives $y(T)=\tfrac12 gT^2$, so $\partial y/\partial g = T^2/2$,
which at $T=3$ is exactly $4.5$:

\begin{lstlisting}
In[15]:= sensitivity 3 "gravity.y" "mass 0"
Out[15]= t = 3 (CVODES, 129 solver steps)
d/d(gravity.y):
  obj0 position [0, 4.500000056696235, 0]
d/d(mass 0):
  obj0 position [0, 0, 0]
\end{lstlisting}

The second answer is worth a moment. In uniform gravity every mass
accelerates equally, so the trajectory does not depend on the mass
\textbf{at all} --- and the derivative comes back as exactly zero, not as
a small number. That is the difference between a real sensitivity
calculation and a finite difference.

\subsection{The one rule that changed for anyone reading the Rust}

The 7.8.0 translation models C's opaque pointers directly. An
\code{N\_Vector} is a handle, and its payload is reached through
\code{N\_VGetArrayPointer}, which returns a \emph{borrow guard}:

\begin{lstlisting}
let y = N_VNew_Serial(neq, &sunctx).ok_or("N_VNew_Serial returned NULL")?;
{
    let mut d = N_VGetArrayPointer(&y).expect("serial vector");
    d[0] = 1.0;
}                       // <- guard dropped here, ON PURPOSE
CVode(&cvode_mem, tout, &y, &mut t, CV_NORMAL);
\end{lstlisting}

Holding that guard across the \code{CVode} call would be a live borrow
of a vector the solver is about to write.
\code{physical\_object/src/integrate.rs} therefore wraps every access in
\code{with\_data} / \code{with\_data\_mut}, which take the guard, do the
work and drop it. Nothing about the command language changes.

\section{Browser videos: watching a run after it has finished}

\kw{scene create} (\S\ref{sec:scene}) gives you a \emph{live} window: it
needs a running posim behind it. The other thing you often want is a
file --- something you can mail to a colleague, open next year on a
laptop with no Rust toolchain, and still scrub frame by frame.

\code{tools/record\_video.py} makes one:

\begin{lstlisting}
cargo build --release -p posim          # once
recorder/src/record_video.py videos/scenes/kepler_ellipse.posim \
     -o videos/kepler_ellipse.html \
     --frames 360 --dt 0.02 \
     --title "Kepler orbit, e = 0.6" \
     --caption "CVODE Adams, 360 frames of dt = 0.02."
\end{lstlisting}

\subsection{How it works, and what it does not do}

\begin{enumerate}
\item It starts \code{posim --machine} --- the same JSON protocol
      JupyterLab uses (\S\ref{sec:jlab}).
\item It runs your setup script one line at a time.
\item Then it loops: ask for \code{\{"op":"state"\}}, keep the answer,
      send \code{step <dt>}, repeat.
\item It writes one HTML file with those frames embedded and a plain
      canvas player around them.
\end{enumerate}

\textbf{Every advance in step 3 is a real SUNDIALS step.} The tool has
no integrator of its own; it is a camera, not a physics engine. And the
page it writes fetches nothing: no CDN, no font server, no analytics.
Open it with \code{file://} on a machine with no network and it works.

\subsection{The player}

\begin{center}
\begin{tabular}{ll}
\textbf{control} & \textbf{does}\\\hline
Play / Pause, or Space & run the recording\\
$\blacktriangleleft\blacktriangleleft$ / $\blacktriangleright\blacktriangleright$, or $\leftarrow$ $\rightarrow$ & one frame back / forward\\
scrub bar & jump anywhere\\
speed & $0.25\times$ to $4\times$\\
drag & orbit the camera\\
wheel, or $+$ / $-$ & zoom\\
$\uparrow$ $\downarrow$ & pan\\
Reset view & back to the framing it opened with\\
trails / labels / contacts & motion trails, body names, gold contact-normal arrows\\
\end{tabular}
\end{center}

The readout in the corner is live per frame: the frame number, $t$, the
total energy, $|P|$, $|L|$, the running collision count, and the method
and $dt$ the recording was made with. \textbf{This is the point of the
format}: you can stop on the frame where something looks wrong and read
the conserved quantities off it.

\subsection{The thirteen shipped recordings}

\begin{center}
\begin{tabular}{p{3.6cm}p{6.4cm}p{4.2cm}}
\textbf{file} & \textbf{what to watch for} & \textbf{measured}\\\hline
\code{videos/kepler\_ellipse.html} & the speed swinging between
perihelion and aphelion on an $e=0.6$ ellipse &
$|dE|/E = 9.8\text{e-}8$, $|dL|/|L| = 1.3\text{e-}7$\\
\code{videos/tumbling\_racket.html} & the Dzhanibekov flip: a
torque-free cuboid spun about its \emph{intermediate} axis turns over,
and over & $|d|L||/|L| = 0$ \textbf{exactly}; $|dE|/E = 6.4\text{e-}9$\\
\code{videos/box\_of\_shapes.html} & a cylinder, a disk and a cuboid
rattling in a rigid \kw{box} 4; gold arrows are the analytic contact
normals, sized by impulse & 36 collision events,
$|dE|/E = 3.4\text{e-}16$\\
\code{videos/double\_pendulum\_hinges.html} & two \kw{hinge} joints
assembled into the chaotic linkage; the gold rings are the joints &
the joints hold to $|g| = 5.6\text{e-}8$; energy wanders 3 parts in
10,000\\
\code{videos/universal\_joint.html} & a \kw{universal} joint carrying a
driven shaft's rotation to a second shaft; the bend flattens out
straight and folds back, and the speed across the joint swings with it
& the bend stops at $\cos\beta = 0.6000004$ against a geometric bound
of exactly $0.6$; the three joints hold to $|g| = 4.0\text{e-}7$\\
\code{videos/ball\_joint\_chain.html} & four links on \kw{ball}
joints, whirling as they collapse; the chain leaves the plane it
started on, which a hinged chain cannot & the four joints hold to
$|g| = 3.3\text{e-}9$; $|z|$ runs from exactly $0$ to $1.7147$\\
\code{videos/rod\_pendulum\_chain.html} & four bobs on four
\kw{constrain} rods, the cheapest linkage there is at one row each,
going chaotic & run continuously at the default tolerance the rods
hold to $|g| = 5.4\text{e-}15$; this recording, 250 cold restarts,
holds $|g| = 7.8\text{e-}8$\\
\code{videos/spinning\_top.html} & a top held at its tip by a
\kw{ball} joint, precessing under gravity & precesses at $1.020440$
rad/s against a closed form of $1.020408$, three parts in $100{,}000$,
without nutating\\
\code{videos/gyroscope\_gimbal.html} & a rotor slung in two gimbal
rings on three perpendicular \kw{hinge} axes; the push goes in about
one axis and comes out about another & total $L\cdot\hat y$ conserved
to $6.4\text{e-}15$; no centre moves further from the pivot than
$7\text{e-}35$\\
\code{videos/cardan\_gear.html} & a wheel inside a ring of twice
its radius, rolling on a \kw{gear} row: the rim point runs along a
\textbf{straight line}, the degenerate hypocycloid & the line is held to
$1.1\text{e-}8$, against $1.8\text{e-}3$ for the same mechanism with the
ratio merely imposed\\
\code{videos/rack\_and\_pinion.html} & a weight on a \kw{rack}
winding up a flywheel, guided by a \kw{prismatic} & the rack falls at
exactly $g/2$, and at the same rate for two different pitch radii\\
\code{videos/piston\_crankshaft.html} & the slider-crank:
\kw{hinge} + two \kw{ball}s + \kw{prismatic}, free-running & follows
$x = a\cos\theta + \sqrt{L^2 - a^2\sin^2\theta}$ to $8.4\text{e-}8$;
stroke exactly $L-a$ to $L+a$\\


\code{videos/cardan\_compass.html} & the same two rings, but with a
\textbf{pendulous} bowl, so gravity is the restoring torque and the
card seeks level & two physical-pendulum periods, $1.878587$ and
$2.307339$ s, measured $1.883426$ and $2.313653$\\
\end{tabular}
\end{center}

\subsection{A trap worth knowing about}

The box recording sets \code{system.g\_constant = 0} on its first line,
and that line is not decoration. posim's default $G$ is 1, so three
bodies rattling inside a small box also \emph{attract each other}, and
the softened pairwise force at \code{softening = 1e-6} is very nearly
singular when two surfaces touch. With gravity left on, the same
scenario drifts \textbf{3.2\,\%} in energy through 25 events; with
$G=0$ it holds to $5.1\text{e-}16$ through 36. Neither number is a
solver defect --- they are different physical systems.

The lesson generalizes: \textbf{a conservation claim carries its system
settings with it.} When a run's $|dE/E|$ surprises you, check
\kw{energy}, \code{system.g\_constant}, \code{system.softening} and
\code{system.uniform\_gravity} before suspecting the integrator.

\subsection{What a BALL joint buys, measured exactly}

\kw{ball} and \kw{hinge} are easy to describe and easy to confuse: both
hold a point, and a hinge additionally holds an axis. The ball-chain
recording turns that sentence into a number.

Take four links laid end to end and start them as a \textbf{rigid
rotation} of the whole assembly about the vertical,
$v = \omega \times r$ with $\omega = [0, 1.5, 0]$, so
$v = [0, 0, -1.5x]$. A rigid motion moves nothing relative to anything,
so it violates no joint of any kind --- as a \emph{position} constraint.
The velocity constraint is where the two joints part company:

\begin{center}
\begin{tabular}{p{6.4cm}p{2.6cm}p{2.6cm}}
\textbf{the same four links, joined by} & \textbf{worst $|g|$} &
\textbf{worst $|\dot g|$}\\\hline
four \kw{ball} joints & $0$ exactly & $\mathbf{0}$ \textbf{exactly}\\
four \kw{hinge} joints about $z$ & $0$ exactly & $\mathbf{1.5}$\\
\end{tabular}
\end{center}

$1.5$ is not approximately anything. It is $\Omega$, the whirl rate,
because the whirl is precisely the component a hinge about $z$ forbids,
and the velocity residual of a forbidden motion is its own magnitude.

\textbf{What follows from it.} A start with $|\dot g| \neq 0$ is not on
the constraint manifold, so before integrating anything the solver must
project the velocities onto it, and that projection \emph{changes the
motion you asked for}. The hinge chain cannot be given this whirl; it
can only be given whatever remains after the whirl is removed. The ball
chain takes it untouched, and \code{initial\_velocity\_projected} is
$0$.

Watch the recording's $z$ axis for the consequence. Every link starts
exactly on the plane $z = 0$, and the chain reaches $|z| = 1.7147$. A
hinged chain would still be at exactly zero, for ever, because that is
what fixing an axis means.

\subsection{One number decides what a suspension does}

Three of the recordings hang a rotor in a mount, and the only thing
that really differs between them is where the centre of mass sits
relative to the pivot. Read them as one experiment with one dial.

\begin{center}
\begin{tabular}{p{2.4cm}p{3.7cm}p{3.2cm}p{3.7cm}}
& \code{spinning\_top} & \code{gyro\_gimbal} & \code{cardan\_compass}\\\hline
mount & one \kw{ball} at the tip & three \kw{hinge}s & two \kw{hinge}s\\
rows & 3 on 6 freedoms & 15 on 18 & 10 on 12\\
centre of mass & $r = 0.6$ \textbf{from} the pivot & \textbf{on} the
pivot & $d = 0.12$ \textbf{below} it\\
gravity & lever arm $r$: a torque & lever arm 0: none & lever arm $d$: a
\emph{restoring} torque\\
result & steady precession $Mgr/(I_3\omega_3)$ & no gravitational
precession & swings back to level, $2\pi\sqrt{I/Mgd}$\\
\end{tabular}
\end{center}

All three carry a \code{uniform\_gravity}. In the middle column it does
nothing whatever, and that is the point rather than an oversight: a
lever arm of zero produces a torque of zero, so a gimballed gyroscope
is not balancing against gravity, it is free of it. Move the centre of
mass off the pivot and gravity comes back --- \emph{above} the pivot it
would tip the thing over, \emph{below} it, it rights it. That last case
is the whole of a ship's compass: it is not held level by its bearings,
it is held level by its own weight.

\textbf{The compass periods.} With the bowl pendulous each axis is an
ordinary physical pendulum. In pitch only the bowl swings, giving
$I = M(3R^2+4h^2)/12 + Md^2 = 0.21046667$ and $T = 1.878587$ s; in roll
ring and bowl swing together, $I = 0.28574980$ against a restoring
$(M_b - M_r)gd$, giving $T = 2.307339$ s. Measured: $1.883426$ and
$2.313653$. The $2.8\times10^{-3}$ excess is not solver error --- halve
the kick and it falls by $4.01$ and $4.02$, the signature of a term in
$\theta^2$. Drive one axis alone and the attribution is exact: the
pitch period comes out $4.99\times10^{-4}$ long against a predicted
$\theta^2/16 = 5.02\times10^{-4}$.

\textbf{Both rings lie flat, and that is not a styling choice.} A
gimbal ring pivots about a \textbf{diameter} --- a line in its own plane
--- so its plane has to contain both hinge axes. Orient it instead so
its symmetry axis \emph{is} the hinge axis and ``swinging'' becomes a
spin about its own axis of revolution: a torus turned that way neither
shows anything nor does anything, and the roll inertia silently becomes
the axial $M_r(c^2 + 3a^2/4)$ instead of the diametral
$M_r(c^2/2 + 5a^2/8)$. That is a factor of two in the ring's
contribution, and it moves the roll period from $2.307339$ to
$2.582727$ s. The recording caught it: the ring sat there visibly inert
while the bowl rocked.

\textbf{A cost of the midpoint pivot.} Two hinges sharing one point
force $p_{\text{frame}} = p_{\text{bowl}} = -p_{\text{ring}}$, so
hanging the bowl $d$ below the centre puts the \emph{ring's} centre $d$
above it. That is not cosmetic: the ring is then top-heavy and
\textbf{subtracts} from the restoring torque, which is why the roll
period carries $(M_b - M_r)$ and not $(M_b + M_r)$, and why the ring is
made light.

\textbf{And a hinge cannot torque about its own axis.} That is
precisely the freedom it grants --- but read the other way round it is
a conservation law, and the gimbal has two. The outermost hinge turns
about the vertical, so nothing can torque the system about the
vertical, and the total $L\cdot\hat y = 0.52688$ is conserved to
$1.4\times10^{-14}$. The innermost hinge's axis \emph{is} the rotor's
spin axis, so nothing can torque the rotor about it either, and its
axial spin $I_3\omega_3 = 3.75$ is conserved to 1 part in $100{,}000$.

The first is held to machine precision rather than to solver tolerance,
and the reason is structural: angular momentum is integrated
\emph{state} in the 13N packing, not a quantity derived afterwards from
positions. The second involves the orientation as well, so it inherits
the tolerance instead.

\textbf{A gimbal holds a point, exactly.} All four bodies are
concentric, so each hinge's midpoint pivot lands on the same origin and
the three axes meet there --- which is what makes this a gimbal rather
than a linkage. Over the whole recording no centre moves further from
that origin than $1.2\times10^{-34}$. The rings only \emph{look} nested.

\textbf{What to watch.} The assembly starts turning about the vertical
at 1 rad/s. Something without a gyroscope in it would simply keep
turning. This does not: the outer ring never gets more than
$7.43^\circ$ from where it began, and the rotation reappears as a
\textbf{tilt} --- the rotor's axis swinging up to $16.33^\circ$ away, at
right angles to the push. Rotation in about one axis, rotation out
about another, is the whole of gyroscopic behaviour, and here it is a
consequence of the two laws above rather than a separate rule.

\subsection{The right motion versus the right mechanism}

The Cardan-gear recording exists to make one distinction concrete, and
it is the distinction the \kw{gear} joint was added for.

A wheel of radius $r$ rolling inside a ring of radius $2r$ sends every
point of its rim along a straight line --- a diameter of the ring. The
hypocycloid of ratio 2 does not approximate a line, it \textbf{is} one.

\textbf{You can get that motion without a constraint.} Set the carrier
turning one way and the wheel the other at the same rate --- which for
this radius ratio is exactly the rolling condition --- and both rates
are conserved, so the ratio persists and the line comes out. Nothing in
the picture tells you it was never enforced.

\textbf{Then disturb it.} The same scene, measured four ways:

\begin{center}
\begin{tabular}{p{8.0cm}p{4.6cm}}
& \textbf{rim point off the line}\\\hline
ratio as an initial condition & $1.795\times10^{-3}$\\
ratio enforced by \kw{gear} & $1.059\times10^{-8}$\\
initial condition, plus a torque on the wheel & $\mathbf{7.468\times10^{-1}}$\\
\kw{gear}, plus the same torque & $1.152\times10^{-8}$\\
\end{tabular}
\end{center}

Undisturbed, the constraint is worth five orders of magnitude.
Disturbed, it is worth everything: the imposed version does not degrade,
it \textbf{collapses}, and the straight line is gone.

\textbf{The general point.} A relation you can only set up in the
initial conditions gives the right motion and the wrong mechanism. The
two are indistinguishable while nothing happens, and a single
disturbance separates them.

\textbf{What the exercise taught.} The joint set is now enough for the
mechanisms this chapter set out to build --- a wheel rolling
\emph{along the ground} is \kw{prismatic} for the line plus \kw{rack}
for the rolling, a piston is \kw{prismatic}, a cam follower is
\kw{prismatic} against whatever drives it. What outlasts the list is
that the two hard joints were the two whose natural coordinate is an
angle nobody counted, and each needed a different trick chosen by what
else was to hand: \kw{gear} has two angles and nothing else, so it hides
the wrapping in a sine and pays with a rational ratio; \kw{rack} has an
angle \emph{and an unbounded distance}, so it unwraps the angle from the
distance and pays nothing; \kw{prismatic} has no angle in its statement
and needed neither. The difficulty was never the mechanism, but whether
the constraint could be written as a function of the state --- and when
it could not directly, whether some other coordinate of the same
mechanism could say what the state had forgotten.

\subsection{A closed form that is exact, and how to actually hit it}

The spinning-top recording is the one place in this chapter where the
answer is a formula you can look up, so it is worth being precise about
what is being checked.

A symmetric top spinning at $\omega_3$ about its own axis, with its
centre of mass a distance $r$ from the pivot, precesses about the
vertical at $\Omega = M g r / (I_3 \omega_3)$. Textbooks label this the
\textbf{fast-top approximation}. The exact condition for steady
precession at tilt $\theta$ from the vertical is
\[
  M g r = \Omega I_3 \omega_3 - I_1 \Omega^2 \cos\theta,
\]
so the shortcut is exact whenever $\cos\theta = 0$ --- a top whose axis
is \textbf{horizontal}. That is why this scene mounts the gyroscope on
its side, with the axis along $z$ and gravity along $-y$: not for the
look of it, but so the thing being asserted is an identity rather than
a limit.

\begin{center}
\begin{tabular}{p{7.0cm}p{7.0cm}}
$M = 1$, $r = 0.6$, $g = 3$, radius $0.42$ & $I_3 = MR^2/2 = 0.0882$\\
$\omega_3 = 20$ & $\Omega = 1.0204081632653061$ rad/s\\
measured over the recording & $1.020440$ rad/s, \textbf{3 parts in
$100{,}000$}\\
tilt out of horizontal & $4.3\times10^{-5}$ --- it does not nutate\\
\end{tabular}
\end{center}

\textbf{Hitting it takes both velocities.} A steady precession is not
``spin it and let go'': give the top only its spin and it is not on the
steady orbit at all. It dips and bobs its way round instead, and at
this spin rate the \emph{average} precession misses the formula by
\textbf{7.8\,\%}. Steady means the whole body is already turning about
the vertical through the pivot as well as spinning about its own axis,
so \code{angular\_velocity} carries $[0, \Omega, \omega_3]$ and
\code{velocity} carries $[\Omega r, 0, 0]$. Both are rigid motions that
leave the pivot point still, so $\dot g = 0$ exactly and nothing is
projected.

The general lesson: \textbf{an initial condition has to satisfy the
constraints at the velocity level, not just the position level}, and a
closed form that assumes a particular steady state has to be started in
that state or it is not the thing you are measuring.

\subsection{A restart is not a continuation}

The rod-chain recording exists to document a trap that costs an
afternoon if you meet it without warning.

\kw{constrain} is the best-conditioned joint in the language --- one
well-scaled scalar equation, $g = |d| - L$, with none of the index-2
orientation coupling that forces a tolerance floor on \kw{ball},
\kw{hinge} and \kw{universal}. So a rod-only system is correctly
\emph{not} floored, and runs at whatever you asked for. Four rods, five
seconds, default $\mathrm{rtol} = 10^{-10}$, $\mathrm{atol} =
10^{-12}$, give a worst $|g|$ of $5.4\times10^{-15}$: roundoff. Record
the same chain and it fails on the second frame, with \code{IDASolve}
returning $-4$.

Nothing about the physics changed. What changed is that \textbf{a
recording is not one integration.} It is one \kw{step} per frame, and
every \kw{step} is a \emph{cold restart} --- a fresh solver, a fresh
multiplier seed, no BDF history to lean on. The tolerance a continuous
run sustains is not the tolerance a restart sustains, and for this
chain the gap is four orders of magnitude:

\begin{center}
\begin{tabular}{p{5.0cm}p{3.6cm}p{4.0cm}}
\textbf{four rods, five seconds} & \textbf{continuous \kw{run}} &
\textbf{250 \kw{step}s}\\\hline
default $10^{-10} / 10^{-12}$ & $|g| = 5.4\times10^{-15}$ & fails on
restart 2\\
floor $10^{-6} / 10^{-8}$ & fine & $|g| = 7.8\times10^{-8}$\\
\end{tabular}
\end{center}

\kw{ball}, \kw{hinge} and \kw{universal} never meet this, and now the
reason is visible: they are floored to exactly $10^{-6} / 10^{-8}$
\emph{automatically}, so every orientation-joint recording has been
asking for restart-safe tolerances without knowing it. The floor those
joints get for \emph{accuracy} is also what makes them recordable frame
by frame. A rod-only scene has to ask by hand.

\textbf{How to recognize it.} A \kw{run} that succeeds and a loop of
\kw{step}s that fails, at the same tolerance on the same system, is
this and not a physics bug. Loosen the tolerance, or do not chop the
integration up.

\textbf{A related sharp edge.} This chain is chaotic, and its margin at
the default tolerance is thin enough to be moved by \emph{rounding the
initial positions}. Written to six decimals, a five-rod version
survives the continuous run; written exactly, it fails at $t = 0.43$.
Four rods is robust either way, which is why the scene ships with four
and writes its coordinates out in full. If a constrained run's success
depends on how you rounded the input, you are on a knife edge, not on a
result.

\subsection{Count your rows before you brace a mechanism}

The universal-joint recording is worth reading as a design exercise,
because the obvious way to build it does not work.

A \kw{universal} holds two things: one shared point, and one right
angle between the two trunnion axes. It does \emph{not} hold the two
shafts straight. The bend between them is free --- that is the whole
reason the joint exists --- so if the output shaft is left dangling it
does not transmit anything, it just swings down past the joint like a
pendulum and folds back on the input shaft at $176^\circ$.

So the output shaft has to be braced. The textbook drawing puts it in a
second bearing, and a bearing is a \kw{hinge}. Try it and IDA quits at
$t=0$ with \code{retval = -4}. Count the rows and the reason is
arithmetic, not numerics. Two free shafts are 12 freedoms, and
\[
  \underbrace{5}_{\kw{hinge}} + \underbrace{4}_{\kw{universal}}
  + \underbrace{5}_{\kw{hinge}} = 14 \text{ rows on } 12 \text{ freedoms.}
\]
Three of those rows are \emph{redundant}: once both shafts are held in
bearings whose axes meet, the shared point the universal joint asks for
is already implied by the geometry. Redundant rows are not merely
surplus --- they make $J M^{-1} J^{\mathsf T}$ singular, and the
constrained Newton solve has no answer to give. posim does no
redundant-row elimination, so it reports the failure rather than
guessing.

The fix is to brace with the \emph{smallest thing that does the job}. A
\kw{constrain} is one row:
\[
  5 + 4 + 1 = 10 \text{ rows on } 12 \text{ freedoms.}
\]
Two freedoms left, full rank, and the mechanism runs. Better still, the
rod fixes the bend in closed form. The output shaft's centre stays
$0.3$ from the cross and $0.4243$ from the post, so it rides the circle
where those two spheres meet, sweeping a cone of half-angle
$\arctan(0.3/0.6) = 26.565^\circ$ about a line itself $26.565^\circ$
off the $x$ axis. The bend therefore runs from $0$ to $53.130^\circ$
and no further:
\[
  \cos 53.130^\circ = 0.6 \quad\text{exactly.}
\]
Scrub the recording to the sharpest frame and the shafts are at
$0.6000004$. That number is never given to the integrator; it comes out
because the hinge, the universal joint and the rod all held at once.

\textbf{The rule.} Before adding a joint, add up the rows it will bring
and compare with the freedoms actually left. If the total meets or
exceeds the freedoms, you have either locked the mechanism or made it
singular, and a fatal return flag is the honest outcome either way.

\end{document}
